\documentclass[10pt,a4paper,twoside,bg=print,twocolumn,openany,nodeprecatedcode]{dndbook}
\usepackage[utf8]{inputenc}
\usepackage{tabulary}
\usepackage[normalem]{ulem}
\usepackage{color,soul}
\usepackage{float}
\usepackage{caption}
\usepackage{wrapfig}
\usepackage{tocloft}
\usepackage[toc]{multitoc}
\usepackage[hidelinks]{hyperref}
\raggedbottom
\extrafloats{2000}
\cftsetindents{section}{0em}{0em}
\cftsetindents{subsection}{0em}{0em}
\cftsetindents{subsubsection}{0.2em}{0em}
\setcounter{tocdepth}{1}
\renewcommand*{\multicolumntoc}{3}
\setlist[description]{nolistsep,listparindent=3pt,topsep=6pt}
\date{}
\begin{document}
\frontmatter
\tableofcontents
\mainmatter
\chapter{Foreword}
\DndDropCapLine{I}{} have always viewed the worlds of fantasy RPGs in terms of domains. The idea that the world is divided into spheres of influence and power, with specific people in charge of those things. Even as a teenager, I played in lots of different fantasy worlds, and I judged them all based on one thing: the map. I used maps as the keys that unlocked the different domains in a world.


Maps are one of the GM's most powerful tools. The map of the local area, the map of the region, the map of the city, the map of the world. Because maps invite exploration. A map explicitly asks the players to pore over it, point to something evocative or mysterious, and ask: "What's this?"

In the '80s and early '90s, there was an explosion of new, official campaign settings. New fantasy worlds for RPGs. And my friends and I were excited to try them. But whenever I saw the mapIn the '80s and early '90s, there was an explosion of new, official campaign settings. New fantasy worlds for RPGs. And my friends and I were excited to try them. But whenever I saw the map of a new world, I always wondered the same thing: "Who's in charge here?" I don't know why that was important to me. Probably, I was just looking for some kind of organizing principle. Some kind of lens to view all this new information through.

I didn't literally think, "Whose domain is this?" I wouldn't have used that language. And I never thought exclusively about barons and dukes and queens. If I pointed to a great forest, I wouldn't expect it to belong to a feudal noble. I know how fantasy works: it's a big forest, so it probably belongs to the elves!

But which elves? Surely a forest this big is an important asset, and some elven court considers it their realm. Their domain. And if not, that's probably because it's contested! A hot or cold war over ancient elven boundaries! Drama!

Or maybe the forest belongs to an order of powerful druids. Maybe they keep it deliberately wild and pool their resources to ensure no single civilization takes it over! Well, those druids would be just another domain.

That's how my mind worked then, and how it still works. It's okay if some places aren't controlled by anyone. But then, as now, when I see a map of a fantasy world and there are no political boundaries, no answers, no information to tell me, "Who's in charge?" anywhere, I get disappointed. I feel like...there's a layer of information missing.

To date, there's never been an edition of the world's most popular roleplaying game whose core rules took the idea of different domains—whether those are counts and dukes and queens, or druid enclaves, wizard circles, or knightly orders—and gave them mechanical features. Not just, "Oh, you get this one feature if you join this organization," but rules for letting different organizations engage in planning, conspiracy, and conflict the way people do.

That's what this book tries to do. I wanted to give you the tools to add organizations and warfare to your games and make doing so fun. When players think about their own party of heroes, even without changing anything about their characters, I hope they'll recognize at least one of the organizations or specializations in here as fitting their group.2 When GMs see all the different NPC realms in the book, I hope they'll want to run new adventures with a hag coven as the villain, building a secret army in the swamp. Or an evil empire run by an immortal dragon. A draconic tyranny!

When players choose an organization and a specialization, their characters gain new mechanical features—but they'll also feel more like a formal group. They'll have to communicate more to use their new domain powers, and will argue about how best to use domain features during a new phase of the game, called intrigue. Everyone will communicate, argue, and compromise—just like the members of a real organization.

This book will give the heroes and villains in your games more influence, more impact. Their reach will extend beyond the end of their swords or the range of their spells. The party's agents can go out and spy, research, build, and negotiate, all while the heroes are dungeon delving. And of course, when those heroes discover that another domain is building an army to march against the people who look to them for leadership, they can build their own army and prepare for war!

This book might not be for everyone. But if any of the above sounds cool to you, we think you'll like it.

—Matt Colville

Orange County, CA

2021

\addcontentsline{toc}{chapter}{Introduction}
\markboth{INTRODUCTION}{INTRODUCTION}
\chapter*{Introduction}
\DndDropCapLine{T}{}here's a lot in this book, but what's here ultimately boils down to two systems: intrigue and warfare.


\textbf{Intrigue} is conflict between two domains. A holy church versus a mystic circle. The Knights Templar against a dragon tyrant. A bunch of disorganized misfit adventurers against a nation of undead!

A \textbf{domain} is a pretty simple thing as a game concept. It's got four skills. It uses actions and reactions and bonus actions, just like characters and monsters. It has three defenses, but no hit points. The only way to "kill" a domain is to neutralize its officers—although some might escape and start a new domain for the sequel campaign! A domain could be a kingdom, thieves' guild, or druid circle!

An intrigue only ever takes place between two domains. That makes the system simple and easy to keep track of. Other domains might get involved, but purely as allies, lending aid. Of course, depending on the adventure, it's possible that a particular intrigue is just a prelude to another larger intrigue as the characters discover the villainous realm they just defeated was only a pawn in an even larger game....

Alongside the intrigue system, the book presents eight different heroic organizations (each of which can be customized with three specializations) and sixteen different NPC realms to use as allies, neutral observers, or enemies.

\textbf{Warfare} is conflict between units of soldiers and monsters. Along with rules for warfare, you'll find seventy-eight special units in this book and another sixty-four in the special unit deck available on our store (at \textit{mcdmproductions.com/warfare-deck}).

So while it might look like an enormous amount of stuff in here at first glance, it's really just two systems, each relatively straightforward, and then a ton of options for those systems.

The book also contains a handful of new items that can affect intrigue or warfare. And because they were popular in \textit{Strongholds \& Followers}, you get some new legendary artifacts called the codices. (Warning! They will unravel the fabric of reality!)

Strongholds \& Followers presented six courts of extraplanar creatures that the gods might send to aid characters in battle if their faith is strong enough. Those were so popular that we expanded them with three additional courts—devils, demons, and undead. New monsters! Some of these creatures use the action-oriented design philosophy mentioned in "Running the Game \#83," designed to make it easier for a GM to use a single monster to challenge an entire party.

Finally, there's an adventure designed to help you use and understand these rules. It's basically a big, detailed example of how the intrigue and warfare systems can work.

\section{How to Use This Book}
First, you should check out the \textbf{Glossary of Terms} section, which covers all the new terms the book uses. Don't worry about reading the whole thing, but just skimming it to familiarize yourself with the concepts will help a lot. And you can always refer back to the glossary if you're looking for a quick definition.

One of the design goals of this book was to make as few assumptions as possible about your game. You can start using these rules right now, with your existing characters in the middle of your current campaign—or even right in the middle of an adventure. In fact, we hope these rules take your group and your adventures and elevate them, granting everything a broader scope and a feeling of raising the stakes. In a world where the forces of powerful domains vie for control of ancient artifacts and forgotten lore, every adventuring party and their enemies can now be listed among those domains.

The easiest way to adopt these rules is to imagine that the heroes in your game are already an organization, with the players picking the specific type of organization from the eight in this book. Then set up the villains of the adventure as a villainous realm, which the GM picks from among the sixteen in the book.

You can stop there if you like. Domains have options that give their officers—both characters and NPCs—more things to do during an adventure or between adventures. But an organization is more than just its officers. There are followers, agents, and employees all ready to carry out the characters' and NPCs' orders. As such, leading a domain means players can have their organization conduct research or espionage while the characters continue exploring dungeons or fighting dragons.

Domains grant their officers access to domain features that become usable during intrigue. The GM will use the rules for intrigue to make the campaign more memorable as it moves toward one or more warfare battles that mark the intrigue's conclusion, with the heroes plotting and scheming against a villain as their organization uses its domain features and domain skills to prepare for battle—even as the villain does the same!

Domains also grant their officers access to domain powers. These are new special features that allow characters to work together during any combat that takes place during intrigue against the officers of an opposed domain—just as the villain and their lieutenants will be working together on that opposing side. This includes the combat that runs alongside the potential climactic final battle at the end of intrigue. (Though not all adventures need to have a final battle, it's great fun!)

The warfare rules in this book are robust enough to be used as their own game. This can be a fun way to play out small skirmishes or major battles, and it works fine without using the rules for intrigue.

But finally—and most ambitiously—all these rules work together as a whole. The heroes become an organization as they plot and scheme against the villain's realm, making allies and raising an army. Time and again, they confront enemy leaders in combat using their new domain powers while their loyal forces clash around them! And though they might not be something that ends up a part of every session, battles between the armies of domains can lend an epic feel to any adventure. It's a little more work, certainly, but that work pays off in the kind of conflict that pushes players and GMs toward the biggest possible adventure finale.

There might be other things you wish this system did. If we've done a good job, it'll be easy for you to design your own optional rules that expand on this framework. We look forward to seeing what you do with the content that follows!

\begin{DndSidebar}[float=!b]{Combat and Battle}
To keep clear the differences between warfare and characters engaging in combat, this book uses "combat" to refer to characters and monsters skirmishing using the game's normal rules. "Battle" is used to refer to armies clashing using the warfare rules.



\end{DndSidebar}

\section{Scale}
Another thing this book makes no assumptions about is the scale of your game—specifically, the units of distance you use on your campaign map. How much territory a domain controls is something left to the GM—if it ever comes up at all. Domains can come into conflict all on their own as villains try to blow up the world or some small part of it, and the heroes trying to stop them likely don't need to know, "Exactly how many square miles do they control?"

In addition to helping keep these rules flexible and easy to adopt, worrying about the literal area of a territory likely isn't something most people will find exciting in a game about pretending to be an elf who fights dragons. So in the end, the rules simply assume it takes about a week (however this is defined in your world) for a domain to carry out a single domain action. And even that can be easily tweaked to any length of time, as it's deliberately abstract. For a particular game, it might work to have a domain turn be a week at the beginning of an intrigue, but then only a single day at the end as the pace quickens and the stakes rise.

\section{The Operation}
"We sell it back to them," Narco said. "Hang on... then what are you gonna put on your bolts?"

Aurex nodded her head in agreement. "Need I remind you," the half-devil assassin said, "before your mind of chemicals and formulae overheats from its own cleverness, it took our people three weeks to steal those vials of \textit{dreamtime}. And they managed to get in, get it, and get out without anyone noticing. Those resources could have been put to better use if that drug doesn't end up on a bolt in Castagan."

Loroyan Thel, the drow brains of the Operation, peered at Narco from across the corner table of the quiet tavern. "We can handle the paladin with or without magic poison, demidevil," he said, almost under his breath.

"And get him to talk?" Aurex said. "Under \textit{dreamtime}, he'll talk while he sleeps—and we'll know it's the truth."

Thel sighed. "I'm aware of that."

"I thought the sleep juice was for your bolts," Zarek said. Not unusually, the half-orc enforcer felt two steps behind the conversation.

"Forget the bolts, Zarek," Aurex said. "I want to hear what our alchemist director of narcotics has in mind."

Narco turned to the half-devil. "Aurex, don't you see?" He smiled riotously, tossed a small vial into the air, then snatched it back again. "It's exactly as you said. Our agents got into the Tower of the Lens, got the drug, got out with no alarms. No searching, no deaths, no guards alerted. No trace."

No one spoke.

"No trace... of us. So we take this back to the Lens, and we tell them we lifted it off someone who tried to sell it to us."

"Ah..." Zarek said.

"No good." Thel shook his head. "They'll rumble that. Wizards are reliably clueless, but they're not stupid."

"Hey, what if we kill someone?" Zarek said. Everyone frowned at him, as this was his solution to everything. It would be on his family crest if fate ever conspired to make a noble of him.

Still, the half-orc had a point, so they let him continue. "Okay, check this out, yeah? If it's just us, sure, they'll rumble that like Thel says. But if we give them a body? And we say, yeah this piece of piss tried to sell us this juice, but we didn't even know what it was! We asked him where he got it, he got defensive. Says he stole it from you wizards, right? And we said, 'Oh, you trying to offload your scrab on us so you can go back to the Lens and frame us? Pull the other one.' Then he got pissy and tried something, and we sorta had to kill him or whatever!"

Everyone was looking at Zarek.

"Right? So they see this body and they start wondering, 'Who's this, right?' And if they're wondering that, they're not wondering... about us."

Thel slowly smiled and looked at Narco, who smiled back and waggled his eyebrows.

Then they turned to look at Aurex, whose visage was inscrutable. Slowly, she spoke.

"It might work."

"Good enough for me," Thel said. "Let's give it a shot. Well done, Zarek."

"Oh, thanks, boss," the enforcer said, beaming with pride.

"You people done scheming?" A small, snarling voice piped up as Garrote, the Operation's chief negotiator and resident goblin officer, hopped up onto a chair.

Thel explained the plan.

The goblin was impressed. "Hey, you came up with that all on your own?" he asked Narco.

"Well, I..."

"Not bad for a drug addict."

"Garrote," Thel warned. Narco looked hurt. Unlike the rest of the officers of the Operation, he wasn't a killer by trade.

The goblin looked around the table. "That it?" he said. "That all you got?"

Aurex took out her sin-metal blade and began polishing it again, looking meaningfully at the goblin.

"Oh yes, you're all very clever," Garrote said. "But unless you've forgotten..." He stabbed a dagger into the map of the temple the group intended to infiltrate, there to confront and combat the order of paladins who were, of late, bent on the destruction of the Operation. "We still got this shit to deal with."

"With all the intel Aurex ferreted out?" Thel said. "Their wizard's true name? The guards you blackmailed?"

"Thel, they have an army. They know we're coming for 'em. They're gonna surround this place so tight a mouse couldn't get in." Garrote pointed to Narco. "And don't get any ideas. I'm not gonna be a mouse again, so put those filthy potions away."

"Yeah, that is a pisser," Zarek said. "We gonna fight our way in, boss?"

"Well, it's been a month. I was hoping a solution to that would have presented itself." Thel frowned. "It's not too late to seek special assistance. We could call up the Crew."

"Don't you get it?" The goblin smiled, chuckling. "Narco here got halfway there, but it's up to the charming goblin to seal the deal."

"Get it over with," Aurex said. She held out her blade to see the dim light of the tavern glint off the shining gray metal before being absorbed.

Garrote's face fell. Aurex and her blade. "Yeah. Anyway. We take the \textit{dreamtime} back to the Tower of the Lens, we give 'em some corpse they can pin all their suspicion on, just like Narco and Zarek say. Then we say 'Goodness, would you like this back? Seems expensive.' But then we say, 'We don't want paying.' Maybe Aurex's people dig up some dirt on the Lens? Should they even have this stuff? It's illegal as shit. So, you know, we embarrass 'em, lie to 'em, give 'em someone else to be suspicious of, a little blackmail thrown in for good measure. Then we just offer to give it back to 'em. Free."

"To what end, Garrote?"

The goblin grinned. "So they volunteer to take care of our little military problem for us."

No one spoke. The officers of the Operation were all uniformly stunned. Thel looked at Aurex and Narco. "They could do it," he said.

Narco nodded. "Lend us a unit of elementals, show how grateful they are."

Aurex frowned, hating to admit the goblin was right. "You got it in you to brace a quaesitor of the Lens?"

Garrote made a "psh" noise. "Wizards are easy. You just gotta flatter 'em before you threaten 'em." He smiled. "By the time I'm done with 'em, they'll think it was their idea."

"I can try and dig up something on their quaesitor," Aurex said. "Give Garrote some ammo."

Thel put his hands flat on the table. "Intel says we got a week before the Order of the Black Pegasus comes after us, and I intend to go after them first. So...work your magics, and let's see if we can't convince the Lens they owe us a favor."

"Blackmail, intel, drugs." Garrote nodded at the three vials of \textit{dreamtime}. "Just a day's work for us, boss."

Zarek finally caught up. "Heh. This is good, you guys. Oh yeah, this is gonna be fun. This is real thieves' guild shit."

\section{Glossary of Terms}
This section collects the terminology underlying the domains, intrigue, and warfare systems at the heart of the book. As you read the rules for the first time, you might want to refer back here to check unfamiliar terms—or to flip forward to the \textit{Heroic Organizations} and \textit{NPC Realms} sections, and to the \textit{Warfare} chapter, to see how these terms play out in the rules.

\subsection{Domains \& Intrigue}
The following terms are used in domains and intrigue.

\textbf{Attitude:} The disposition of another domain toward yours. Used to establish alliances through diplomacy.

\textbf{Combat:} The normal round-by-round conflict between characters in the game, including the officers of domains.

\textbf{Communications:} One of a domain's three defenses. Communications represents how well agents and followers understand the goals and orders of the domain's leaders, and can relay the current state of the domain back to them.

\textbf{Decrement:} To decrease the number on a die by 1.

\textbf{Defense Levels:} A number that represents the strength and effectiveness of a domain's Communications, Resolve, and Resources. During intrigue, a domain wants to increase their own defense levels and decrease those of the opposed domain. Defense levels don't modify defense scores.

\textbf{Defense Score:} A number that represents the security of a domain's Communications, Resolve, and Resources.

\textbf{Development Point:} Spent by the officers in a domain to customize its stats.

\textbf{Diplomacy:} One of a domain's four skills, used to negotiate with other domains.

\textbf{Domain:} A group of characters or creatures working together to further some long-term goal. The term "domain" encompasses the leaders of the domain, the agents and followers who work for them, and the territory they influence.

\textbf{Domain Feature:} A special benefit or action used during an intrigue on behalf of a domain by one of its officers.

\textbf{Domain Power:} A special benefit or feature available to officers of a domain during combat against officers of an opposed domain.

\textbf{Domain Turn:} The time in which each domain in an intrigue conducts one domain action. Often a week, but could be any span of time determined by the GM.

\textbf{Espionage:} One of a domain's four skills, used to spy on or sabotage other domains.

\textbf{Found:} To create or establish, in regard to domains founding strongholds.

\textbf{Intrigue:} Conflict between two domains in which each takes domain actions, using domain features and domain skills to affect the defenses of the opposing side and prepare for battles. Also the period of time covering this conflict.

\textbf{Leader:} The NPC officer who runs a villainous realm.

\textbf{Lieutenants:} The NPC officers of a villainous realm who obey the realm's leader.

\textbf{Lore:} One of a domain's four skills, used to research obscure knowledge.

\textbf{Muster:} To raise new warfare units for a domain.

\textbf{Officer:} One of the player characters or NPCs in charge of a domain.

\textbf{Operations:} One of a domain's four skills. Operations is used to muster new units and to take actions not covered by any other domain skill.

\textbf{Organization:} A player-controlled domain.

\textbf{Party Sheet:} The record of an organization's stats, defenses, and features.

\textbf{Power Die:} A die that an officer can roll at the start of a battle. The die is placed with its result facing up in a domain's power pool.

\textbf{Power Pool:} The shared group of power dice all officers in a domain access when using their powers.

\textbf{Realm:} An NPC-controlled domain.

\textbf{Resolve:} One of a domain's three defenses, representing how committed a domain's followers are to the cause of the domain's leaders.

\textbf{Resources:} One of a domain's three defenses. Resources covers both wealth and the availability of anything a domain needs to operate.

\textbf{Stronghold:} A domain's headquarters, whatever form that might take.

\textbf{Test:} To resolve a skill roll for a domain. Used in place of "check" to make it clearer when the rules are talking about interactions between domains rather than between characters and other creatures.

\textbf{Title:} A special feature an officer gains as a result of commanding a domain.

\textbf{Villianous Realm:} An NPC realm in direct conflict with the players' organization.

\subsection{Warfare}
The following terms are used in warfare.

\textbf{Activation:} The act of selecting a unit, determining its movement and actions, and resolving them. A unit's activation in battle is effectively equivalent to a character's turn in combat.

\textbf{Adjacent:} Spaces and units are adjacent if one is above, below, left, or right of the other. Spaces and units that are diagonal to each other are not adjacent.

\textbf{Aerial:} Units that can fly.

\textbf{Allied Units:} All units that are part of an army or are friendly to that army.

\textbf{Artillery:} Siege engines and archers.

\textbf{Attack:} The bonus to a d20 roll representing a unit's ability to successfully execute an attack order and engage an opposed unit. Opposed by Defense.

\textbf{Battle:} Conflict between two armies.

\textbf{Battle Magic:} Magic items crafted by officers and given to units in an army.

\textbf{Bringing the Siege:} The test that determines which domain marches on the other at the start of a battle. This determines which side defends a stronghold and which side attacks.

\textbf{Broken:} A broken unit is removed from the battlefield, usually because it suffered its last casualty. Broken units can be reformed by rallying.

\textbf{Casualty:} The higher a unit's casualties, the more damage it can take. Think of casualties as hit points for units.

\textbf{Casualty Die:} The die placed on each unit card that tracks how many casualties the unit has remaining. A unit's casualty die is initially equal to its size.

\textbf{Cavalry:} Highly mobile ground troops, usually mounted.

\textbf{Center:} The most protected rank on the battlefield. Archers begin battle in their army's center rank.

\textbf{Combat:} The normal round-by-round conflict between characters in the game, including the officers of domains.

\textbf{Command:} The bonus to a d20 roll representing a unit's ability to correctly interpret complex orders and execute them successfully.

\textbf{Commander:} The character or NPC who controls a unit.

\textbf{Decrement:} To decrease the number on a die by 1.

\textbf{Defense:} A numerical value representing a unit's ability to maneuver in such a way as to avoid an opposed unit's attack. Used as the DC for an opposed unit's Attack test.

\textbf{Defenseless:} Lacking an army in a battle.

\textbf{Deployment:} The placing of units in legal spaces on the battlefield. Usually the first thing that happens in a battle.

\textbf{Diminished:} A unit is diminished when its current casualties are half or less than its starting casualties.

\textbf{Disband:} A unit that disbands is gone forever and cannot be rallied.

\textbf{Exposed:} A unit that is vulnerable to enemy cavalry because it is not properly protected by other units.

\textbf{Fortification:} A construction in one or more spaces on the battlefield, which grants units various bonuses during the battle.

\textbf{Front:} The vanguard rank of an army plus all of an opposed side's ranks.

\textbf{Increment:} To increase the number on a die by 1.

\textbf{Infantry:} Basic ground troops, trained to fight in melee and hard to kill.

\textbf{Inflict Casualties:} To damage a unit. If a unit inflicts 1 casualty on another unit, the target loses 1 casualty.

\textbf{Levies:} Laborers who briefly put down their pitchforks and ploughshares and pick up sword and pike to serve their domain.

\textbf{Maneuver:} A special action a unit can take, many of which require a Command test.

\textbf{Martial Advantages:} Special features that units gain based on the character class of the officer commanding those units.

\textbf{Morale:} A bonus to a d20 roll representing a unit's ability to maintain discipline in the face of overwhelming odds, magic, exotic enemies, and impending destruction.

\textbf{Opposed Units:} All units that are part of the army of the other side, or which are friendly to that army.

\textbf{Order of Battle:} The rules that describe which units can attack which other units.

\textbf{Power:} A bonus to a d20 roll representing a unit's physical strength and its ability to inflict casualties on an opposed unit. Opposed by Toughness.

\textbf{Rank:} A row on the battlefield.

\textbf{Rear:} The last rank on the battlefield, whose units defend the center rank from enemy cavalry. Usually populated by levies. (Wish them luck!)

\textbf{Reserve:} A well-protected rank on the battlefield where units can be stacked until needed.

\textbf{Special Unit:} A warfare unit that cannot be mustered through a normal Operations check, or that is unique to a specific domain.

\textbf{Suffer Casualties:} When a unit takes damage, it suffers casualties. If a unit suffers 1 casualty, it loses 1 casualty.

\textbf{Tier:} A measure of how powerful a unit is. Higher-tier units are more effective in battle, and harder to muster.

\textbf{Toughness:} A numerical value representing a unit's physical hardiness and its soldiers' ability to withstand attacks and keep fighting. Used as the DC for an opposed unit's Power test.

\textbf{Unit:} A group of soldiers, civilians, or monsters who make up an army.

\textbf{Unit Dependency:} The rule governing how many units of each tier can be fielded.

\textbf{Vanguard:} The first rank of an army, which faces the opposing army. Usually populated by infantry.

\chapter{Domains \& Intrigue}
\DndDropCapLine{E}{}very fantasy RPG campaign contains domains. The local church is one. The baron, their court, and the people of the barony are another. The thieves' guild. The secret society of rangers and druids. Whatever your choice of fantasy subgenre, your campaign contains many different domains feeding a network of alliances, allegiances, suspicions, and grudges.


The rules in this book give you a straightforward way to represent those domains with game mechanics, and to allow the player characters to step up onto a larger stage and translate their influence from adventuring into political power. Player-run domains are called \textbf{organizations}. NPC-run domains are called \textbf{realms}. This book presents eight different player organizations, each with three specialties, and another sixteen NPC realms.

When two domains come into direct conflict, this is called an \textbf{intrigue}. The purpose of intrigue is to move toward a climactic confrontation involving combat between the officers of opposed domains and battles between their armies. Some adventures might end with just a final warfare battle at the end of intrigue, while others will feature numerous smaller battles leading up to a final battle alongside a final combat between the characters and the officers of an opposed domain. A domain can be destroyed only by neutralizing its officers. As such, once an intrigue starts, it continues until initiative is called for the final showdown between the domains' officers.

\section{The Core Assumptions}
Running nations and guilds and churches and spy networks can be incredibly complex. It's easy to imagine a highly detailed game system that can model a world at this level. Tracking a barony's food supply, determining how much iron is available to a duchy, mapping out a spy network's elaborate web of agents, or even maintaining a list of a local religion's temples and shrines—it's all a complicated process. Games like that already exist, and are typically complex grand-strategy simulations that might even require a computer to run them.

This book takes a different approach.

First off, these rules understand that you're already playing a complex fantasy RPG. We all know that just managing a single character or a world of NPCs can be a lot of work. So ideally, any domain management system has to be something that sits lightly on top of the game that players and GMs are already enjoying.

Furthermore, the game almost always has a heavy focus on fighting monsters. Most monsters are defined by their combat prowess, and many of a character's best class features revolve around making them a better monster fighter. By assuming that training followers makes characters better at the stuff they're already good at—whether that means research, study, or fighting—running a domain provides another way to improve the effectiveness of player characters and NPCs in combat.

Finally, almost every adventure involves some villain plotting and scheming while the heroes try to stop them. This system assumes that the villain runs their own villainous realm, that the player characters' organization is actively trying to stop them, and that this conflict will eventually come to a head in an all-out battle—or perhaps more than one. This is the same way things would pan out if you weren't using the new rules in this book. But with these rules, every officer in a domain has new special features called domain powers. These powers require communication and cooperation to use, even as armies clash in warfare to all sides.

In simplest terms, the rules in this book assume that characters continue adventuring as they normally do. While the characters are adventuring, slowly working their way toward a final confrontation with the villain, their domain is looking for allies to lend aid, investigating the villain and their lieutenants to learn their secrets, and disrupting the workings of the villain's realm to weaken that realm's forces in a final showdown. Even for a campaign that's in the middle of an adventure right now, players and GMs can adopt the new rules in this book to provide a framework for the campaign's existing plots and intrigue—then raise the stakes for both heroes and enemies in a warfare battle!

The adventure included with this book—"The Regent of Bedegar"—serves as an example of how a straightforward adventure can include rewards such as military units or resistance to an opposing realm's features based on a domain's actions during the adventure.

\section{Characters and Officers}
\textit{Kingdoms \& Warfare} assumes that all the player characters in a campaign will be officers in a domain organization. These rules introduce a new dice pool (described below) shared by all the officers of an organization, which grants characters access to domain powers. Those new features are based on what type of organization the players choose, and require the players to communicate and cooperate as a team. (A villainous realm has its own dice pool that the leader and their lieutenants gain access to for their own domain powers.)

A character can be an officer in more than one organization, but they can be involved only in one intrigue at a time and can benefit only from the effects of one party sheet at a time. (That's the record of a domain's stats, defenses, and features.) A wizard might be a member of the party's noble court organization and run their own arcane order on the side. But during an intrigue involving the noble court, the character is focused on that organization while their arcane order tends to its own affairs.

\section{What is an Organization?}
An organization is a domain built around the officers who founded it—the player characters. The players decide what their characters' organization does, and who its initial allies and enemies are. Among those allies, an organization also includes the NPC followers, retainers, and lieutenants the characters attract as a result of rising fame or infamy earned through adventuring. For example, if the party's organization is an underworld syndicate, that organization includes all the members of the syndicate down to the lowest-level agents keeping their ears to the ground and feeding the characters information. Likewise, an organization set up as a noble court might include large numbers of farmers and laborers who rely on the heroes for protection, and who are ready to serve to show their thanks.

Using the rules in this book, an adventuring party becomes an organization when it founds a stronghold, typically by buying, building, discovering, or inheriting it. Founding a stronghold announces to the world that the adventurers are more than just mercenaries, and are ready to get involved with local affairs in one way or another.

\section{Describing a Domain}
Each fully detailed domain (whether a player character organization or a villainous realm) is described by its collection of skills, defenses, domain powers, and domain features, as well as its size.

A domain has four \textbf{skills}: Diplomacy, Espionage, Lore, and Operations, each of which has a modifier. A domain's officers use these skills during intrigue.

A domain has three \textbf{defenses}: Communications, Resolve, and Resources. Each defense has a numerical value that provides a target number for tests made with domain skills during an intrigue. Each defense also has a level that can be raised or lowered. Any domains can gain bonuses or take penalties during battle or combat based on their domain defense levels.

Each domain also grants its officers access to special \textbf{domain powers} they can use in combat. Each officer gets a \textbf{power die} that they roll at the beginning of any combat against one or more officers of an opposed domain. The roll goes into a \textbf{power pool} shared by all the officers of a domain, with these dice fueling each character's domain powers.

Every domain also provides a number of unique \textbf{domain features} that officers can make use of during intrigue, allowing the domain to affect the opposed domain's defenses and units, to muster special units, and make other preparations for battle.

Finally, each domain has a \textbf{size} that determines how many turns it can take during intrigue, how large its officers' power dice are, and the scope of certain benefits available to the domain's army (as detailed in the \textit{Warfare} chapter).

\section{Intrigue}
Conflict between domains is called intrigue, and serves as the backdrop to the warfare battles that can play out between the heroes and the villains. The characters can use their organization's skills outside of intrigue, making use of Diplomacy to forge alliances, Espionage to gather intelligence, and so forth. But once the heroes decide it's time to act and stop the villain (or once the villain decides to stop the heroes), the GM announces that intrigue has begun as its own special phase of the game.

Intrigue occurs between two domains—by default, the heroes' organization and the villain's enemy realm. NPC realms might also be involved, but they don't act on their own during intrigue. Rather, they lend aid to one side or the other.

Intrigue is divided up into domain turns, during which the players and the GM make use of the domain features and special actions available to them. Intrigue ends once both sides have completed all their domain turns, at which point the GM will set up a final showdown between domains— involving combat between the characters and their enemies, a climactic battle between the armies of powerful domains, or both!

\section{Founding an Organization}
The players create an organization when their characters acquire a stronghold. They might do so by spending time and money to buy or build a stronghold, or through other means such as discovering an old ruin, clearing it out, and fixing it up. Some campaigns might even begin with the heroes inheriting a stronghold!

The game's core rules note prices for different strongholds of different sizes, but there's also an entire book dedicated to this—\textit{Strongholds \& Followers}. (You don't need that book to use the rules in this book, but it might be fun.) 

Once the characters gain a stronghold, they naturally start to attract followers. Folks hear about the deeds they've done, notice the new headquarters, and volunteer to help or serve the characters. That stronghold and those followers are the foundation of the player characters' organization.

Of course, the GM can waive any of these requirements if doing so is a better fit for the campaign. The only thing that's really necessary is to have people working for the characters. But without some kind of headquarters, even if it's just the local tavern (an establishment according to the \textit{Strongholds \& Followers} rules), there's no physical structure to defend and nothing for an enemy to attack—and these things are important to the new rules in this book, as you will see.

\subsection{Choosing an Organization Type}
Eight different types of organizations for player character domains are presented in this book, and each of those has three specializations. Two of those organizations—the noble court and the adventuring party—are good for existing campaigns in which the players don't want to make new characters to use these rules, or for campaigns with a wide range of character types. The rest are themed more narrowly, and are best for new characters built around a specific organization type—and for parties where characters are focused on similar or even the same classes. (See \textbf{Granting Titles} below for more information.)

That said, even the most strongly themed organizations don't make any assumptions about what classes officers can or should be. With only a few exceptions, none of the domain powers that officers can take make reference to class features. Any group of characters could decide to be a thieves' guild or a knightly order. After all, every criminal enterprise needs wizards and clerics, and you don't have to be a heavily armored and chivalrous paladin to follow a knight's creed.

\subsection{Development Points}
Each organization is set up with a \textbf{party sheet} that defines the organization's starting skills, its defenses, its domain powers, and its domain features. The sheet even has a place for the characters' pool of power dice.

When characters found an organization, they begin with \textbf{8 development points}, which are available to spend on skills and defenses. The players take turns passing the party sheet around, with each player spending 1 point and checking off the appropriate blank box on the party sheet (see below) until all points have been spent.

Players can make choices for spending development points as a group, or each player can do what they like individually. The GM can also weigh in if necessary, depending on the circumstances of the campaign. For example, in a campaign where one character has inherited a barony from their parent and the other characters are the new baron's lieutenants, it might make sense for the new baron's player to spend most of the organization's initial development points. (That's less fun, though!)

Every time a domain levels up, the officers gain an additional 8 development points to spend on improving their organization. The players once more pass the sheet around, with each player spending 1 point until all points are spent.

Each development point spent lets a player mark off a box on the party sheet, moving to the right. When a box is marked off, the box on its right is marked off next. When a box with a value in it is marked off, the organization's skills or defenses become that number. For example, if an organization has +2 marked off in Diplomacy and the players mark off the next two boxes to reach +3, the organization's Diplomacy modifier increases to +3. If Resolve starts at 10 and the players mark off the next six boxes, the organization's Resolve increases to 15.

Each organization starts with its skills and defenses at specific values, representing what the organization is naturally good at based on its type. (See the Heroic Organizations section on page 31 for details.) Mark off the boxes with the starting values and all the boxes to the left of those boxes. Then players spending development points continue marking off boxes to the right.

Any time an effect changes a domain's skills or defenses, the effect always refers to the modifier, not the number of boxes. If a domain specialization grants an organization a +1 to Diplomacy, that means the domain's Diplomacy modifier increases by 1. Only development points are applied one box at a time.

\subsection{Raising an Army}
As a setup for using the warfare rules, when the characters found an organization, the players immediately muster four Tier I units of their choice to start up their army, from any ancestry the GM agrees the organization has access to. Each unit must be controlled by an officer of the organization, and each officer can command a number of units equal to their proficiency bonus. (The \textit{Warfare} chapter has more information on all these things.)

During intrigue, players can muster more units for their organization's army. Rules for this are covered under the \textbf{Operations} domain skill in this chapter (page 23) and \textbf{Building an Army} in the \textit{Warfare} chapter (page 100).

An army obeys the commands of an organization's officers, and can be deployed outside of intrigue or warfare for narrative purposes. As a rule, an army can perform any activity that an organization's officers could perform but which the characters are too high a level to bother with, from escorting a diplomat to quelling an uprising of cultists.

The Operations domain skill can be used to resolve tasks an army performs outside of battle.

\subsection{Granting Titles}
Each organization grants its officers access to five titles that are distinct to each type of organization. Each character claims a different title within an organization, with that title granting new features and proficiencies. A character can have more than one title, but they can benefit from only one at a time, switching titles during a long rest. If a party includes more than five player characters, the GM can allow specific titles in an organization to be duplicated or allow a character to take a title from a different type of organization.

Most organizations' titles are designed to support playing an organization whose officers are all of the same class, by giving those officers a wider range of features. For example, in a normal game, a party consisting of all rogues will deal formidable damage and be very stealthy indeed! At the same time, though, the lack of a tank and a healer makes such a group more fragile and less versatile than a typical group of adventurers. But in a thieves' guild organization whose officers are all rogues, titles help shore up these deficiencies.

As an optional rule, the GM might consider allowing titles only when the players create a single-class party. This isn't a strict requirement, but players and GMs should all be aware that some titles might become overpowered if stacked with similar benefits from class features or feats.

\section{Organizations and Warfare}
This book is built around the iconic idea of the heroes fighting an adventure's villain in one or more epic bouts of combat, while outside the characters' or the villain's sanctum, a battle rages between the armies controlled by both sides. If the characters run an arcane order or a druid circle, their army might be composed of elementals or treants. A thieves' guild might hire mercenaries or field units of elite scouts, harriers, and sappers. But regardless of which organization the players choose, \textbf{the rules assume they have an army}. The battle that rages alongside the combat has a mechanical impact on the characters. When the army of one domain or the other wins a battle, the officers of the victorious domain gain a morale surge that offers a one-time benefit in combat.

For some players, this idea fits the fantasy concept just fine. But for others, the idea of their characters' thieves' guild or arcane order fielding an army might seem strange, and out of line with their ambitions or style of play, or the way they imagine their organization. If this is ever the case, the players and the GM should feel free to ignore the warfare system in this book altogether. You can still use the intrigue system, and these rules will still work and be fun. Or you might try ignoring both warfare and intrigue to just use an organization's skills and powers! The power dice mechanic works well on its own, and will be fun even if the characters never raise an army.

Even so, the GM is encouraged to describe the villain's realm as still having an army! But those forces will just be off doing something terrible in grand cinematic fashion (and possibly creating the reason the heroes need to stop the villain in the first place). Some sort of final combat still happens, with the heroes fighting the Scion of Orcus or Lord Saxton or what have you. But the villain's armies are busy wreaking havoc elsewhere.

\section{Before Intrigue}
Domain-level play works the same way the skirmish-level game works. Just as each character has skills they can use outside of and during combat, an organization has skills it can use before and during intrigue. In the same way that characters decide to draw steel and initiate combat, an organization can decide to deploy its agents and begin intrigue. Likewise, the GM can decide, "You've pushed this NPC far enough! It's time for intrigue!" Or, "This villainous realm has decided the party is an active threat, and they initiate intrigue!"

In any event, the GM determines when intrigue begins, regardless of which side initiates it.

Before intrigue begins, the characters' organization can probe and test their enemy, but \textbf{they can't use their skills directly against a villainous realm}. That kind of action is what begins intrigue. This is a fine line, and it's up to the GM to adjudicate it. But in general, if the characters can make use of their organization's skills to deal with everyday people (for example, learning something about a villainous realm by making general inquiries using Diplomacy) or known information (for example, researching a villain's history using Lore), then those activities can remain part of the regular game. (See below for more information on these domain skills and their use.)

A Diplomacy test made against one of Lord Saxton's allies might provide some idea what kind of army he's preparing without starting intrigue. But an Espionage test to spy on Lord Saxton would initiate intrigue. Likewise, a Lore test to learn what Lord Saxton can do in combat might be allowable outside of intrigue if the GM decides Lord Saxton has shown off his combat skills often enough that there are people outside his domain who know that information. But any domain skill that requires targeting a villainous realm (including its agents, libraries, and lands) should start intrigue.

\section{Domain Skill Tests}
A domain skill test works exactly like a character or creature making an ability check. A d20 is rolled, modifiers are added, and the total is compared to the test's Difficulty Class, determined by the GM. If the total is equal to or higher than the DC, the test succeeds. Tests can likewise be made with advantage and disadvantage, just like ability checks. The rules use the word test instead of check to help differentiate domain skills from the abilities and skills used by the characters.

\subsection{Using Domain Skills Outside of Intrigue}
Players don't have to wait until intrigue begins to use the skills of their characters' domain. A domain and its agents are always standing by, waiting for orders. For example, a Diplomacy test—along with some good roleplaying and effort on the characters' part—could change the attitude of a clan of elves toward an organization, so that once intrigue begins, the DC to convince the elves to aid the organization is lower.

Outside of their place in intrigue, there are no hard and fast rules describing how domain skills can be used. Players can suggest ideas, and if the GM agrees that an idea makes sense, they can try it! The goal with skills outside of intrigue is to keep the rules purposefully broad, so each table can develop their own standards.

How long it takes to resolve a domain skill test outside of intrigue is also up to the GM. A good rule of thumb is one domain skill test per week, but this is entirely dependent on the timeframe of the game and the campaign. If you don't track the passage of time closely in your game, it might be easier to tie domain skill tests to major narrative events, such as the awarding of XP to characters, the end of a session, or even the end of a major encounter. (See \textit{Time and Domain Turns} below for more guidance in this area.)

Regardless of how the GM rules, characters use their domain's skills by giving orders to the organization's followers and waiting, sometimes for days or even weeks of campaign time. During that time, those followers work to carry out those orders. Making an Espionage test to determine what military units a villainous realm has mustered can take an enormous amount of effort to pull off, including long days spent in research and networking.

\section{Running Intrigue}
A villainous realm is defeated only when its leader is defeated in combat, whether that means death, or surrender, or stranding them in another part of the multiverse. The defeated realm's agents and followers might organize under a different leader later, but that will be a different domain with its own stats and features.

Before the final showdown between the characters (in their roles of officers of their domain organization) and the officers of the villain's realm, the characters' organization conducts an intrigue against that realm. During intrigue, the players can use their organization's skills to put their army together, or sabotage certain elements of the opposed domain's power structure to impose penalties on its skills or defenses. At the same time, agents of the villainous realm are making their own attempts to weaken the characters' organization.

Much like deciding when to call for initiative rolls in combat, the GM decides when intrigue starts. This usually happens once the heroes encounter one or more agents of the villainous realm, or the effects of those agents' villainy. This might occur at a different point in each adventure, but once the characters have direct evidence of the existence and operation of a villainous realm, intrigue can begin.

As a rule, the heroes' organization can't research or confront a villainous realm they don't know about, or which is beyond the reach of their influence. The characters can't use their friendly neighborhood thieves' guild to sabotage another domain on the other side of the planet.

An intrigue is focused on one boss villain. A short adventure with only one main villain will likely have only one intrigue, which begins as soon as the heroes encounter the agents of that villain (and recognize them for what they are) or the results of their villainy. For complex campaigns in which the characters will face off against increasingly powerful sub-bosses before getting to the main boss, it's a good rule of thumb to plan for additional intrigues—perhaps even one per sub-boss. Each sub-boss fight can mark the end of a chapter or act, and will have its own intrigue leading up to it. As such, each sub-boss should be treated as the leader of their own villainous realm, in service to the end-boss's realm.

\section{Domain Turns and Actions}
During an intrigue, each \textbf{domain can act a number of times equal to 4 + the domain's size}. Each time a domain acts during an intrigue is called a domain turn. This is true for the heroes' organization as well as the villain's realm.

\subsection{Choosing Actions}
During a domain turn, the characters' organization and the villainous realm act at the same time. The GM chooses a domain action for the villainous realm and keeps it secret. Then the characters choose a domain action for their organization. After the players declare the organization's domain action, the GM reveals the action of the villainous realm. Once actions are declared, the players roll any necessary skill tests for their characters' organization's domain action while the GM does the same for the villainous realm. Then the players share the results of their organization's actions followed by the GM sharing the villainous realm's results.

\subsection{Who Knows What?}
Spies from both domains monitor each other, so almost all actions a domain takes and the results of those actions are public knowledge. The exception is mustering units and using domain skills to improve a domain's stats (its domain skills and defenses). Enemies know when units are mustered or when a domain's stats are improved, but knowing which units a domain has or what the domain's current stats are requires a successful Espionage test. See \textbf{Domain Skills} on page 18 for more information.

\subsection{Taking Domain Turns}
Sometimes one domain in an intrigue has a greater domain size and can take more turns than the other. If the larger domain initiated the intrigue, those extra domain turns are taken at the start of the intrigue before the other domain can act. If the larger domain didn't initiate the intrigue, the extra domain turns are taken at the end of intrigue, after the other domain can no longer act.

\subsubsection{Domain Features}
When a domain takes a turn, one of its player character or NPC officers runs the domain turn, directing the domain to make use of its domain features, or to use domain skills as a domain action. Many features also involve the use of domain skills. When using a domain skill, if the officer is proficient with any of the domain skill's associated character skills, they can add their proficiency bonus to the skill test roll. However, this can be done only once per officer per intrigue.

There are three domain action types: domain action, domain bonus action, and domain reaction. An officer can make one domain bonus action in addition to a domain action during a domain turn. Each officer can take only one domain reaction during an intrigue. Any domain reactions that an officer can make have specific triggers noted in their descriptions.

Because domain features are used during intrigue, many of them relate directly to the warfare rules as they set up battles featuring the army of the characters' domain. The \textit{Warfare} chapter (page 92) has more information on the rules and terminology referenced in those domain features.

\subsubsection{Whose Turn Is It?}
There is no limit to the number of rolls an officer can make during intrigue—only on the number of domain turns each domain can take. Each officer issues orders to the domain's agents and followers, which can easily require multiple domain turns. However, because only one of that officer's skill tests can gain the officer's proficiency bonus during an intrigue, it's best to have different characters and NPCs take domain turns. The more one character must spread their attention between different activities, the less effective their leadership.

\subsubsection{Time and Domain Turns}
As noted above, the GM decides how often domains can take domain turns during regular adventuring. Intrigue begins when one domain threatens another; continues through periods of campaigning, battles, or both; and lasts until a final confrontation when the characters, the villains, and their armies face off for the last time. During the intrigue, an entire adventure might play out, featuring several encounters and multiple days or weeks of game time. One intrigue might last just a short period of time while another lasts a month or more, simply because of the different scope of the adventure.

The GM should feel free to use any of the following options when deciding when and how often to call for new domain turns:


\begin{itemize}
\item One domain turn per week.
\item One domain turn per day.
\item One domain turn after each combat encounter.
\item One domain turn per short rest, as the officers dash off orders to their agents.
\item Two domain turns every time experience is awarded.
\end{itemize}

The GM should also feel free to change these milestones on a per-adventure basis depending on how much time is passing. Or even to change them within an adventure to represent the changing pace of action! This might mean calling for two domain turns per day at the start, and then calling for a domain turn after every encounter as the end of intrigue gets closer. The final domain turn usually occurs just before the climactic confrontation begins!

Ultimately, there's no wrong way to do this, and the GM should feel free to adapt how and when they call for domain turns to suit the adventure they're running and the ways in which they track time in their game.

\section{Domain Skills}
Each domain has four skills: Diplomacy, Espionage, Lore, and Operations. These skills define how an organization or realm interacts with the world and with other domains. Skills can be used before and during intrigue. During intrigue, making a domain skill test is a domain action.

Each domain skill represents an entire department inside a domain, including followers and retainers who are experts in their fields. Those NPCs can mount operations on their own with just the direction of a domain's officers, with those operations not requiring any direct involvement of the heroes or the villain. As such, not only do domain skills allow a domain to do things while the characters are adventuring, they allow the domain to do things outside the scope of adventuring, from research to sabotage to negotiation.

Players are encouraged to be creative with their domain's skills, the same way players can improvise how they use a character's abilities and skills. Don't worry about "doing it right." Each group will develop their own standards for which skills are best and what the appropriate defense or DC is, in much the same way each table develops its own sense of when Dexterity (Acrobatics) is preferable to Strength (Athletics).

During intrigue, each skill can be used to improve a domain's defense levels or to target the opposed domain and reduce its defense levels. But it's up to the players and the GM to imagine how this plays out, by translating skill tests and their effects into game terms.

\subsection{Adjusting Defense Levels}
Once intrigue begins, each domain involved in the intrigue prepares for some sort of final confrontation between their officers—the player character heroes versus the GM's villain and their lieutenants. Each side can use domain actions to make use of their domain skills, with the goal of either improving a domain's defense levels, or of targeting the opposed side's defenses to make them less prepared for the warfare battles that take place during intrigue.

\subsubsection{Improving a Domain's Defenses}
Each domain has three defenses—\textbf{Communications}, \textbf{Resolve}, and \textbf{Resources}—each of which is rated at a level from 3 to −3 (as described in detail under \textit{Domain Defenses} on page 24). \textbf{The DC to improve one of a domain's defense levels is 13 + the defense's current level.} So a domain skill test to take its Resolve level from normal (level 0) to loyal (level 1) has a DC of 13. Improving Resolve level from loyal (level 1) to fanatic (level 2) is a DC 14 check, and going from revolt (level −3) to rebellious (level −2) is DC 10. The level of any defense cannot be increased above 3.

When the players want to affect one of their organization's defense levels using a domain skill, an acting officer chooses the intended defense level, a skill, and sets out how that skill is to be used. For example, Anna (playing a character named Judge) wants to raise the Communications level of her group's mercenary company, the Chain of Acheron, using that organization's Espionage domain skill. She tells the GM: "I want to require all our agents to use only verbal communications, nothing written down, and to use our secret battle language."

This is a great idea—but it's not necessarily easy to implement. The Chain of Acheron has a lot of soldiers, agents, and allies in the city, so there's no guarantee that these orders will work well enough to make a difference. To make the domain skill test, Anna rolls a d20 and adds the Espionage bonus for Acheron's Chain. She also adds Judge's proficiency bonus, since Judge is proficient in Investigation and has not already used their bonus during this intrigue. (See the domain skill descriptions below for more information on associated character skills.)

The Chain of Acheron's Communications defense is currently secure (level 1), which sets the check as DC 14. If Anna succeeds, the domain's Communications goes up one level. If she rolls badly, Communications remains secure but does not improve, with the GM likely interpreting this as meaning Judge's plans were too ambitious.

A defense level can normally be improved only by one level by making a domain skill test as a domain action. However, if a player rolls a 20 on the test, this represents an unexpectedly successful result that improves the defense by two levels.

Villainous realms can increase their defense levels the same way, with the GM making the decision as to which officer, skill, and specific approach are involved.

\subsubsection{Lowering a Domain's Defenses}
In the same manner, the players can target a villainous realm in the hopes of lowering that domain's defense levels. An acting officer picks a skill and one of the opposed domain's defenses, then describes how they intend to use that skill to compromise the defense. If the scenario makes sense to the GM, the officer rolls a d20, adds the domain skill's modifier, adds their proficiency bonus if applicable, and compares the result with the enemy realm's defense score. If the result is equal to or greater than the chosen defense's score, that defense is lowered one level. If the player rolls a 20 on the check, the chosen defense is lowered two levels. The level of any defense cannot be decreased below −3.

Villainous realms can target the defenses of the heroes' organization in the same way, with the GM making the decision as to which officer, skill, and specific approach are involved.

\subsection{Other Uses for Domain Skills}
The only limit to what benefits characters can earn through the use of their domain skills is the players' imaginations and the GM's sense of what's fair and reasonable. Each skill gives some ideas of what can be done with it, and a GM should refer to the DCs noted in the core rules when trying to set difficulties for an optimal balance between what's realistic and what's dramatic.

\subsection{Diplomacy}
\textit{Associated Character Skills: Insight, Persuasion}

The Diplomacy domain skill represents a domain's ability to get what it wants without resorting to war or espionage. It's primarily used to make alliances with other domains. An alliance is not a mechanical notion, but is just an agreement between two domains. The terms of the agreement are up to the officers of each domain, but usually an alliance means that two domains each agree to help the other in times of conflict.

Some alliances are purely defensive. "We agree to go to war against an enemy who attacks either of us." More rarely, alliances can be offensive. "Lord Saxton must be stopped! Will you help us?" Offensive alliances are more difficult to establish, though, and another domain might require great assurances before entering into one. "We might attack Lord Saxton if things are as you say. But what proof can you offer that Saxton is a villain?"

\subsubsection{Gaining Special Units}
As a domain action, a domain can make a Diplomacy test to petition an NPC realm for military aid for an upcoming battle. The DC depends on the attitude of the NPC realm, as noted on the following table.


\begin{DndTable}{XcX}

Attitude & Diplomacy DC & \\
Hostile & 18 & \\
Suspicious & 15 & \\
Neutral & 13 & \\
Friendly & 10 & \\
Allied & 8 & \\
\end{DndTable}

This Diplomacy test can be used more than once per intrigue, but each NPC realm can only grant one unit per intrigue.

Each NPC realm has a special unit that another domain can gain as a result of an alliance. This is the same special unit the domain can muster using a domain action, as noted in its description. For example, a successful Diplomacy test made by the characters' organization to seek aid from a giant jarldom NPC realm gets the organization a Pet Roc special unit to command in an upcoming battle. Each NPC realm has only one special unit to provide, so characters had better recruit them before their enemies do! A special unit gained through an alliance with an NPC realm returns to the NPC realm after the battle in which it is used.

\subsubsection{Allies and Enemies}
A hostile NPC realm is not the same thing as a villainous realm. A villainous realm is actively trying to destroy the heroes' organization, whereas a hostile realm is merely one with a tradition of enmity toward the heroes' domain or other domains like it. A local druid circle might have a troubled history with a noble court that rules the land near their forest, creating potential problems for the different noble court run by the characters even if the characters have never interacted with the druids before.

At the GM's discretion, a realm that is allied with the characters' organization can send units to aid the organization in battle even without a Diplomacy test. But sometimes allies are busy or occupied with other conflicts, in which case the GM will ask for a test to be made.

Diplomacy can also be used to convince an NPC realm currently allied with an enemy to sit a battle out, or to convince a neutral realm to lend aid beyond units. Such aid might include a powerful lieutenant, or information that would benefit the domain requesting assistance.

If a domain fails a Diplomacy test to influence or gain allies from an NPC realm, the domain can make a new Diplomacy test to achieve the same goal on a new domain turn. However, if a Diplomacy test fails by 5 or more, the NPC realm's attitude toward the Domain moves one step toward hostile.

\subsection{Espionage}
\textit{Associated Character Skills: Investigation, Stealth}

Espionage allows a domain to learn the secrets of other domains, including their plans, their current activities, and who they are allied with. Successful Espionage tests often decrease the level of an opposing domain's Communications defense, reflecting how successful attempts at destabilizing the opposed domain have been.

Espionage can be used to learn the makeup of an opposed domain's units and the disposition of its allies in advance of a battle. The DC is usually set by targeting one of the opposed domain's defenses, but the GM can also set a DC arbitrarily (as per the guidelines for setting difficulty in the core rules). Espionage does not reveal the nature of specific magic in use by an opposed domain, but it might reveal that an opposed domain has some magical resources or secrets, which a successful Lore test can then ferret out.

When an opposed domain makes a skill test to muster units or raise its defenses, the GM doesn't automatically reveal the specifics of a successful result. The characters learn that the opposed domain mustered units, but not how many units or what kind. Likewise, the specifics of which defense was raised and by how much remain a mystery. If the characters want to know these things, they need to make a successful Espionage test as a domain action, typically against the opposed domain's Communications. On a success, the players can choose to know one of the following pieces of information:


\begin{itemize}
\item The target domain's current defense scores
\item The target domain's current defense levels
\item The target domain's skill bonuses
\item The target domain's current units
\item Another statistic or piece of information, at the GM's discretion
\end{itemize}

Knowing the kinds of units an opposed domain has mustered helps a domain plan its own army. For example, if a villainous realm is building aerial units, the characters' organization will want some aerial units or artillery of their own to take them on. Likewise, if an opposed domain keeps boosting its Resources defense levels, the characters might decide to try to lower that defense and undo some of that progress!

If a domain feature requires a domain to know a secret piece of information to work properly, then the domain learns the information when they successfully use the feature. A separate Espionage test is not necessary. For example, a fey court domain can duplicate an opposed domain's highest-tier unit with the Your Own Worst Enemy feature. When the court succeeds on the Lore test to use the feature, they automatically learn what the target domain's highest-tier unit is.

\subsection{Lore}
\textit{Associated Character Skills: Arcana, History, Religion}

Lore represents a domain's ability to research both magical and historical knowledge. This is a broad and wide-ranging skill, which can be used outside of intrigue to unlock arcane discoveries, dig up the details of forgotten history, discover the existence of lost spells, or reveal the answers to ancient mysteries. The only limit on what a Lore check can accomplish is what the GM rules is reasonable.

As a rule, the Lore domain skill should not be used as a substitute for adventuring. But at the GM's discretion, it can be used to supplement the normal adventures of a domain's player character officers. For example, a team of NPC agents might be sent out to recover the legendary headpiece for the Staff of Ra, even as the characters quest after the location of the map room in the hidden Well of Souls. In general, if it would be easy but time-consuming for a domain's officers to acquire important knowledge or rare antiquities, then the Lore domain skill can be used to delegate that task. But if doing so is merely difficult, then the officers should undertake that task themselves.

Lore can be used to learn an opposed domain's secrets when those secrets are explicitly magical, such as the nature of supernatural attacks or defenses, or the presence of any famous magic items or artifacts in a domain's arsenal. Likewise, the presence and features of extraordinary units in a domain's army is something a Lore test can ferret out.

\subsection{Operations}
\textit{Associated Character Skills: Athletics, Insight}

Operations represents a domain's ability to muster new units, as well as to perform many of the basic functions of maintaining a domain. This is intended to be a catchall domain skill for any activity the players or GM might think of that doesn't easily fit into another skill.

\subsubsection{Building Fortifications}
Heroic organizations and NPC realms can acquire strongholds and other fortifications during the course of a campaign in any number of ways. The characters might seize enemy fortifications after a successful battle (or have their own fortifications seized by a villainous realm), claim fortifications as a reward for adventuring, or find some other means of acquiring an existing fortified site.

A domain can also build specific battle-ready fortifications (as discussed in \textbf{Fortifications} in the \textit{Warfare} chapter) in any space on a battlefield the domain controls. Doing so during intrigue requires a successful Operations test made as a domain action. A successful DC 11 test is required to build a stone fence, while a successful DC 13 test is required to build a tower. Multiple fences and towers can then be put together in different configurations to build castles, keeps, and other fortifications.

The GM decides the time scale it takes to build a fortification, but here are some typical guidelines:


\begin{itemize}
\item If the Operations test to build the fortification succeeds by 5 or more, the fortification is built at the end of the current domain turn, or after 1 week if the domain is not in intrigue.
\item If the Operations test succeeds by 4 or fewer, the fortification is built at the end of the next domain turn, or after 2 weeks if the domain is not engaged in intrigue.
\item If the Operations test fails by 4 or fewer, the domain's fortification is still constructed, but it takes 3 domain turns (or 3 weeks) to do so.
\item If the test fails by 5 or more, the fortification can't be built, and the players or the GM must wait for a new domain turn (or for 1 week to pass if not engaged in intrigue) before they can try again.
\end{itemize}

\subsubsection{Mustering Units}
A domain's officer can use a domain action during intrigue to muster a new unit. The Operations skill covers this activity, though an Operations test is normally not required unless the GM feels there's a chance that the mustering might fail. This new unit belongs to the officer who mustered it. An officer can control a number of units equal to their proficiency bonus.

The units an officer musters can be from any ancestry the GM agrees that character has access to, and which conform to the rules for unit dependency. (See \textbf{Building an Army} on page 100 in the \textit{Warfare} chapter for information on ancestries, unit dependency, and much more.)

Units mustered through alliances and Diplomacy tests are controlled by a character but not owned by that character. They return to their home NPC realm after the battle in which they are used.

\subsubsection{Modifying Units}
A domain can spend money to upgrade a unit's equipment from light to medium, medium to heavy, and so forth. This requires a domain action and is considered a use of the Operations skill, but normally no Operations check is required. A domain can upgrade as many units as are controlled by the officer using a domain action. Upgrades to a unit can be made to any level a domain can afford, from light to super-heavy.


\begin{DndTable}{XX}

Upgrade from... & Cost\\
Light to medium & 500 gp × the tier of the unit\\
Medium to heavy & 1,000 gp × the tier of the unit\\
Heavy to super-heavy & 2,000 gp × the tier of the unit\\
\end{DndTable}

A domain cannot upgrade a unit that was mustered through alliances and other Diplomacy checks. Likewise, levies cannot have their gear upgraded. Upgraded units retain their original tier.

\subsubsection{Disbanding Units}
Eventually, a domain will max out the number of units it controls, making it impossible to muster any new units. When this happens, a domain can disband any number of units to free up space in its army. This isn't an action, but is simply an order an officer can issue under the umbrella of the Operations skill. The domain does not regain any money spent on equipment upgrades for units that are disbanded.

\subsubsection{Special Units}
Every type of domain has an action it can use to summon a special unit consisting of rare and powerful creatures. These include such units as the Crew, a unique group of roguish heroes that an underworld syndicate organization can muster; the drake-mounted knights known as the Praetores Draconis, mustered by a draconic empire realm; and many more. See \textit{Heroic Organizations} (page 31) and \textit{NPC Realms} (page 65) later in this chapter for full details.

A domain can muster a special unit once per intrigue. If an attempt to muster a special unit fails, additional attempts can be made in the same intrigue until successful. When the battle in which a special unit is used ends, the unit leaves the control of the domain, heading off on important business, returning to its home plane, and so forth. A special unit can be mustered again during a new intrigue, responding when duty calls.

A domain can make use of other special units in addition to the unit it can muster, by making use of the Diplomacy domain skill (as discussed earlier in this chapter). By talking to the leaders of NPC realms, the players (and their enemies) can see if those realms are willing to lend their own special units for a battle.

\subsubsection{Beyond Intrigue}
The Operations domain skill covers many mundane activities outside of intrigue, including building roads or watchtowers, sending out a unit of scouts to explore the local wilderness, and so forth. Anything a domain might reasonably be able to do but which isn't covered by another domain skill can be accomplished using Operations.

\section{Domain Defenses}
Every domain has three defenses, all of which can be targeted by an opposing domain: \textbf{Communications}, \textbf{Resolve}, and \textbf{Resources}. Defenses operate at one of seven levels, each of which has a name and a numerical rating from 3 to −3. A defense's level represents how robust and functional that defense is at any given moment. Each defense also has a numerical score used as the DC for the domain skill test of an opposed domain wanting to sabotage that defense.

At each level above normal (level 0), the level benefits for a defense are cumulative with the levels below. For example, if a domain's Communications levels are unbreakable (level 3), that domain gains all the bonuses for secure, coded, and unbreakable Communications. Similarly, at each level below normal (level 0), the level penalties for a defense are cumulative with the levels above.

At the beginning of intrigue, all of a domain's defenses start at level 0. The players and the GM can use domain actions during intrigue to raise the level of a domain's defenses or to lower the levels of an opposed domain's defenses (see \textit{Adjusting Defense Levels} on page 19). The level of any defense cannot be increased above 3 or decreased below −3.

Once intrigue ends, and after any final battle between the characters' organization and the villainous realm is over, the defense levels of both domains slowly return to normal, moving one step closer to 0 each week.

\subsection{Communications}
The Communications defense is a measure of how rapidly and accurately information is transmitted between a domain's officers and followers. Every domain, from a hard-as-nails mercenary company to a grove of peaceful druids, relies on its network of followers, retainers, and servants to accomplish its goals. And without effective communications, those goals can easily break down.

Communications directly affects an army's ability to coordinate its activities. When Communications is high, an army can efficiently maneuver into a better starting deployment—the arrangement of units under the warfare rules. (See the \textit{Warfare} chapter for more information.) If Communications is poor or compromised, a domain's enemies know what its officers and agents are doing, and can interfere with deployment by sending false signals to the domain's units.

\begin{DndTable}[header=Communications Levels]{cX}

Level & Effect\\
3 & Unbreakable\\
2 & Coded\\
1 & Secure\\
0 & Normal\\
-1 & Compromised\\
-2 & Garbled\\
-3 & Broken\\
\end{DndTable}

\textbf{Unbreakable (3).} At the end of the next deployment, this domain chooses any allied infantry or artillery unit on the battlefield, then moves that unit to any unoccupied space on the battlefield.

\textbf{Coded (2).} At the end of the next deployment, this domain chooses two allied units in any rank on the battlefield and swaps those units' locations.

\textbf{Secure (1).} At the end of the next deployment, this domain chooses an allied unit in any rank and moves that unit into any adjacent space.

\textbf{Normal (0).} No effect.

\textbf{Compromised (−1).} At the end of the next deployment, the opposed domain chooses one of this domain's units in any rank on the battlefield and moves it into any adjacent space.

\textbf{Garbled (−2).} At the end of the next deployment, the opposed domain chooses one of this domain's cavalry units, which is removed from battle and does not deploy until the end of the first round of battle.

\textbf{Broken (−3).} At the end of the next deployment, the opposed domain chooses two of this domain's infantry units, which are removed from battle and do not deploy until the end of the first round of battle.

\subsection{Resolve}
Resolve measures the commitment that the followers of a domain have to its cause, and depends on many factors—including how well those followers understand the domain's cause. A domain's followers and army, by default, share the philosophy of the domain's officers and leaders, whatever that philosophy is. For example, the followers of an arcane order domain value neutrality and knowledge, while the soldiers of a despotic regime domain revel in the glory and victories produced by their iron-fisted leaders.

A domain's Resolve has a direct effect on its officers. If everyone working for the domain is committed and engaged, it makes training easier and gives followers hope and confidence in victory. But the reverse is equally true. If everyone working for a domain is convinced of the domain's imminent failure, it makes training more exhausting and affects officer confidence. High Resolve means that a domain's officers are energized, coordinated, and convinced of victory, granting benefits to those officers in combat against the officers of an opposed domain. Poor resolve means that officers are distracted and tired, and their performance against foes from an opposed domain will suffer.

\begin{DndTable}[header=Resolve Levels]{cX}

Level & Effect\\
3 & Zealous\\
2 & Fanatic\\
1 & Loyal\\
0 & Normal\\
-1 & Discontented\\
-2 & Rebellious\\
-3 & Revolt\\
\end{DndTable}

\textbf{Zealous (3).} During any combat against officers of an opposed domain, each of this domain's officers has advantage on attack rolls until the end of their first turn.

\textbf{Fanatic (2).} As a reaction when an enemy starts their turn during the first round of any combat against officers of an opposed domain, one of this domain's officers who has not yet acted can cast a spell or make a weapon attack.

\textbf{Loyal (1).} The speed of each of this domain's officers increases by 10 feet during the first round of any combat against officers of an opposed domain.

\textbf{Normal (0).} No effect.

\textbf{Discontented (−1).} The speed of each of this domain's officers decreases by 5 feet during the first round of any combat against officers of an opposed domain.

\textbf{Rebellious (−2).} During any combat against officers of an opposed domain, the first saving throw made by one of this domain's officers against a spell or effect directed by an enemy has disadvantage.

\textbf{Revolt (−3).} The first attack roll made by each of this domain's officers during any combat against officers of an opposed domain has disadvantage.

\subsection{Resources}
The wealth of a domain is measured as Resources, though this defense represents more than just money. Resources includes whatever a domain values and collects, whether gold, knowledge, secrets, or things more esoteric. When a domain's Resources is high, it directly affects the gear of that domain's units, granting them improved power or damage. Poor Resources affects morale. Troops that haven't been fed or whose armor and weapons are in dire need of repair become agitated.

\begin{DndTable}[header=Resolve Levels]{cX}

Level & Effect\\
3 & Booming\\
2 & Abundant\\
1 & Surplus\\
0 & Normal\\
-1 & Low\\
-2 & Poor\\
-3 & Bankrupt\\
\end{DndTable}

\textbf{Booming (3).} During the first round of battle, each of this domain's artillery units that inflicts casualties inflicts 1 extra casualty.

\textbf{Abundant (2).} Each of this domain's cavalry units has advantage on Power tests until the end of the first round of battle.

\textbf{Surplus (1).} Each of this domain's infantry units has advantage on Power tests until the end of the first round of battle.

\textbf{Normal (0).} No effect.

\textbf{Low (−1).} Each of this domain's artillery units has disadvantage on Morale and Command tests until the end of the first round of battle.

\textbf{Poor (−2).} Each of this domain's cavalry and aerial units has disadvantage on Morale and Command tests until the end of the first round of battle.

\textbf{Bankrupt (−3).} Each of the domain's infantry units has disadvantage on Morale and Command tests until the end of the first round of battle.

\section{Domain Size}
Size determines how far a domain's power extends, and the size of the power die used by the domain's officers. Size is a relative measure of the reach and influence of a domain, though not an absolute measure of an area of land the domain controls. In one GM's campaign, where the map covers a whole region hundreds of miles across, a size 3 domain might extend its influence for 60 miles out from its stronghold. In another campaign that takes place entirely in one district in a large city, a size 3 domain might cover only a few blocks.

An organization starts at size 1 if the characters are its founders (but might start at a larger size if the GM has the characters take over an established organization). \textbf{Every time the characters' organization defeats another domain, the organization's domain size increases by 1}. The size of the organization's power die likewise increases, and the organization gains more development points to spend on the party sheet, as shown on the Domain Size table. (See \textit{Development Points} earlier in this chapter for more information.)

A domain cannot have a size greater than 5. However, at the GM's discretion, an organization of domain size 5 can still continue to gain new development points by defeating villainous realms.

Villainous realms are not built the same way the characters' organization is. The size of a villainous realm is thus determined by the GM. (See \textit{Building a Villainous Realm} on page 28 for full information.)

\begin{DndTable}[header=Domain Size]{cccX}

Domain Size & Power Die & Development Points & \\
1 & d4 & 8 (starting points) & \\
2 & d6 & +8 & \\
3 & d8 & +8 & \\
4 & d10 & +8 & \\
5 & d12 & +8 & \\
\end{DndTable}

\section{Domain Powers}
Domain powers represent the new features that a domain's officers earn as a result of the research and training they do with their agents. The officers of a thieves' guild work and train with their bravos and apprentice thieves, and as a result, become better at what they themselves do. Likewise, the stewards of a druid circle spend time between adventures studying, tending to nature, and training their acolytes, and so gain greater insight into the summoning of nature spirits.

Each type of domain—both heroic organizations and NPC realms—grants a number of unique domain powers, whose mechanics are broken out in the domain type's description. But all domain powers rely on the use of a shared resource—a pool of power dice.

\subsection{Power Dice}
Each officer in a domain—player characters and the villain and lieutenants of an opposed domain alike— gets one power die, with the die type determined by the domain's size (see above). Any officer can choose to roll their power die immediately after they roll initiative at the start of any combat (no action required). Once a power die is rolled, it cannot be rolled again until the officer who rolled it finishes an extended rest. (An extended rest is defined in \textit{Strongholds \& Followers} as 1 week of study and training spent at the stronghold of the officer's domain. The GM might use this as a guideline, or set some other parameters for what an extended rest means in the campaign.)

Each domain power allows the officers in the player characters' organization to use some or all of the dice in their shared pool to produce crazy new effects in combat. At the same time, the officers of a villainous realm will use their own domain powers and power dice to fuel their domain's push for victory over the characters.

Any power dice that aren't used are removed from a domain's pool at the end of the combat in which they were rolled.

For example, Anna, Lars, Grace, and Tom are playing the officers in a thieves' guild (one of the options for an \textbf{underworld syndicate organization}; see page 61). They're just starting out, so each of their characters has a d4 power die. In a fight against a local enemy thieves' guild (an NPC realm) known as the Clock, Lars waits until initiative is determined, and decides to roll his power die, getting a 3. He adds this to the empty pool on the party sheet.

Following suit, the other officers, including officers of the Clock, all decide to roll their power dice. Anna rolls a 4, Tom a 1, and Grace another 4, all of which are added to their pool. The heroes' pool now has four dice in it: a 1, a 3, and two 4s. Whichever hero acts first can take any or all of those dice out of the pool, depending on which domain power they intend to use.

Rolling power dice together in the same combat provides a potential benefit by increasing the number of dice in the pool, and increased chances of high rolls. But some players might want to not roll their power dice, so as to save them for another upcoming combat. Likewise, the GM might decide to not roll power dice even when the players are, if they think their villainous officers will face the characters again before everyone has time to take an extended rest.

\subsubsection{Decrementing}
When a player or GM activates a domain power, the rules for that power sometimes instruct them to decrement the power die, usually at the end of each turn of the character who used the power. Decrementing the power die means to decrease the number on the die by 1. When a power die showing a 1 is decremented, the die is spent and the power activated with that die is no longer in effect.

If a domain power that calls for decrementing a power die has been used and is currently active on an officer, that power cannot be activated again for the same officer until the power die being decremented is spent.

\section{Villainous and NPC Realms}
The player characters' enemies also run their own domains, known as villainous realms. A villainous realm is run by a \textbf{leader}, whose chief agents are known as \textbf{lieutenants}. Leaders and lieutenants are the domain's officers, and are referred to specifically by certain domain features.

\subsection{Building a Villainous Realm}
The GM builds villainous realms in much the same way the players build the characters' organization. However, the process for the GM has more options, and allows for the creation of a domain that is larger than the size 1 domain the characters must start with. GMs can establish realms at the start of the campaign or at any point within it, and NPC and villainous realms can rise, fall, and evolve as the GM sees fit.

A GM builds a villainous realm using the following process:


\begin{itemize}
\item The GM picks a domain type from the \textit{NPC Realms} section, or a domain and a specialization from the domain types presented in the \textit{Heroic Organizations} section (both of which follow later in this chapter). This allows for the creation of nefarious holy churches, druid circles, and mercenary companies, in addition to the traditionally antagonistic realms of despotic regimes, fey courts, and others.
\item The GM chooses the domain size of the villainous realm. In most cases, it's best to set the realm's size within 1 of the size of the characters' organization by the time both sides get into an intrigue.
\item The GM spends development points improving the villainous realm's skills and defenses. The GM has a number of points to spend equal to 8 × the villainous realm's domain size.
\item The GM picks stat blocks for the villainous realm's leader and lieutenants, and gives one of the domain's titles to the leader. (If the domain has been chosen from the \textit{NPC Realms} section, it has only one title.) The lieutenants of the realm do not get titles (because the GM already has enough to worry about), but all officers might have access to the realm's domain powers and domain features.
\item The GM picks a stronghold for the villainous realm.
\end{itemize}

For example, consider a group of heroes running a size 2 holy church. With that organization's latest foes vanquished, the GM wants to reveal the vampire Countess Sanguin's villainous realm, which has observed the characters' last battle from the shadows and wants to remove the threat their organization poses. The GM knows that Sanguin's servants include vampire spawn, wights, and other undead, so they make the villainous realm an undead dominion (page 86). Since Sanguin's domain has been established for some time, the GM makes it size 3. It poses a threat to the heroes' organization, but not an overwhelming one.

The GM then spends 24 development points to improve the undead dominion's domain skills and defenses, using the same party sheet the players use for the characters' domain. They choose the \textbf{vampire} stat block for Countess Sanguin and give her the deathlord title. The GM gives Sanguin three \textbf{vampire spawn} lieutenants and names each of them, then gives the undead dominion a stronghold: an ancient castle.

\subsection{Building an NPC Realm}
Building an NPC realm is easier than building a heroic organization or a villainous realm, because NPC realms don't directly influence intrigue or battles. They might be persuaded by the characters' organization or the villainous realm to lend aid during a battle, but the GM can easily set that up without finalizing all of an NPC realm's stats. (That said, a GM is free to flesh out any NPC realm with the same level of detail as a villainous realm—especially if the NPC realm has enemy potential!)

The baseline process for creating an NPC realm is as follows:


\begin{itemize}
\item As with creating a villainous realm, the GM picks either a domain and a specialization from the \textit{Heroic Organizations} section or a domain from the \textit{NPC Realms} section. This choice doesn't imply a specific role for the NPC realm, though, which might be friendly or antagonistic toward the characters' organization as the GM decides.
\item The GM picks the NPC realm's starting attitude toward the characters' organization, using the table under \textit{Diplomacy} in \textbf{Domain Skills} (page 20). This attitude should be based on the NPC realm's history not just with the characters, but also with domains similar to the characters' organization.
\item The GM names the NPC realm's leader. This NPC doesn't necessarily need a stat block, but it helps to note a few details about their appearance and personality.
\item The GM chooses the NPC realm's domain size (which need not bear any relationship to the size of the characters' organization), and picks a stronghold the NPC realm controls.
\item The GM can freely swap out the special unit the NPC realm can normally muster with a special unit from another domain. Perhaps a fey court of elves in the campaign has a roguish flair and is able to bring together the Crew (normally associated with an underworld syndicate domain) instead of the Court Jesters.
\end{itemize}

For example, consider a GM who wants to create a domain of secluded stone giants living in the mountains near the characters' stronghold. The GM begins with the giant jarldom NPC realm and gives the isolationist giants a suspicious attitude toward the characters. Smallfolk warred with the giants centuries ago, and while the humanoids might have forgotten that past, the stone remembers. The stone giants are led by Jarl Klanga. She's young and defensive minded, after her parents died in the same landslide that resulted in her losing an arm. The GM decides the realm doesn't hold much influence given their isolationist nature, and gives it a size of 1. The domain has a stronghold: a mountain fortress called Citadel Adamantine.

\subsection{Creating NPC Realms Together}
In a campaign in which the GM enjoys letting the players join in on the fun of setting creation, it can be cool to let each of the players design an NPC realm for the GM to use. The GM can assign each of the players a domain size and an attitude that their NPC realm has toward the characters' organization, then let them determine the rest of the details. The players get in on the fun of building the easiest type of domain, and the GM has to do less work. Win-win!

\begin{DndSidebar}[float=!b]{Reskinning Domains}
GMs might find that their worlds include a villainous or NPC domain that's not an automatic fit for any of the concepts this book provides. For example, consider a campaign that features a powerful faction of gnolls. Even though there's no NPC realm set up as a gnoll war band, there's almost certainly another domain that fits. If the gnolls are proud warriors, the GM can use the orc clan domain and call it something else. If the gnolls are religious zealots, their domain could be built as a hidden cult religious order. Or for gnolls with a culture of invention and innovation, the gnomish kingdom might be a perfect fit!



Likewise, GMs shouldn't feel locked into building NPC realms based on ancestry. The goblinoid realm in a campaign doesn't have to be a goblinoid coalition domain. The GM can make it an arcane order, a despotic regime, an undersea colony, or whatever makes sense for their world.



\end{DndSidebar}

\section{Heroic Organizations}
When a group of characters founds a domain, the players choose one of the eight organization types described in this section. Each organization features three specializations that the players also choose from, creating a wide range of possibilities for the heroes' domain.

Though this section is primarily for player characters, and is written to speak to the players, a GM can also use the domains presented in this section to create enemy and NPC realms.

\begin{DndSidebar}[float=!b]{Gaining New Proficiencies}
If a character gains a proficiency from a domain title that they already gain from another source, that character can take another proficiency of the same kind (skill or tool) instead.



\end{DndSidebar}

\section{Adventuring Party}
Mercenaries. Troubleshooters. Rag-tag heroes trying to get by—except "heroes" might not be quite the right word. Your characters have banded together as an adventuring party, a haphazard group of individuals traveling the realm with a common goal. Though you might not be the most diplomatic or noble of groups, you get the job done when it counts. It just might be a bit...messy.


\begin{itemize}
\begin{multicols}{2}
\item \textbf{Skills}
\item \textbf{Diplomacy:} -1
\item \textbf{Espionage:} +0
\item \textbf{Lore:} +1
\item \textbf{Operations:} +2
\item \textbf{Defenses}
\item \textbf{Communications:} 11
\item \textbf{Resolve:} 12
\item \textbf{Resources:} 10
\end{multicols}

\end{itemize}

\subsection{Domain Titles}
Adventuring party officers gain access to the following titles, each of which grants a character the noted additional features.


\begin{itemize}
\item Cap'n
\item Doc
\item Know-It-All
\item Smart Mouth
\item Weirdo
\end{itemize}

\subsection{Domain Power}
Adventuring party officers gain access to the following domain power.

\paragraph{Optional Feature}
Never Tell Me the Odds|KaW

\subsection{Domain Features}
Adventuring party officers gain access to the following domain features.

\paragraph{Optional Feature}
Call In a Favor|KaW

\paragraph{Optional Feature}
To Arms, Fellow Adventurers (Special Unit)|KaW

\subsection{Specialization}
When you found your adventuring party, you'll also choose a specialization. Are you an explorers' society, dedicated to seeking out the unseen frontiers of the world? Are you a mercenary company, completing dangerous jobs for coin? Or maybe you're just a band of disorganized misfits, doing all you can to help those in need?

\subsection{Disorganized Misfits}
You're...well, you're trying your best. You might not be the most disciplined or most knowledgeable. You all fight a lot and rarely agree on anything. But you're hard to kill. More importantly, though, you all have good hearts, and you're willing to do whatever it takes to help those in need.

{\par \vspace { 9pt plus 3pt minus 3pt } \noindent  \DndFontTableTitle{Skill and Defense Bonuses} \nopagebreak }

\begin{itemize}
\begin{multicols}{6}
\item \textbf{Operations:} +1
\item \textbf{Resolve:} +1
\end{multicols}

\end{itemize}

\subsection{Domain Power}
Disorganized misfits officers gain access to the following domain power.

\paragraph{Optional Feature}
Avenge Me|KaW

\subsection{Domain Features}
Disorganized misfits officers gain access to the following domain features.

\paragraph{Optional Feature}
Like in the Great Stories|KaW

\paragraph{Optional Feature}
Whatever it Takes|KaW

\subsection{Explorers' Society}
As the name suggests, your adventuring party is a band of explorers, dedicated to traversing uncharted forests and plumbing the depths of mysterious caverns in search of new discoveries.

{\par \vspace { 9pt plus 3pt minus 3pt } \noindent  \DndFontTableTitle{Skill and Defense Bonuses} \nopagebreak }

\begin{itemize}
\item \textbf{Lore:} +2
\end{itemize}

\subsection{Domain Power}
Explorers' society officers gain access to the following domain power.

\paragraph{Optional Feature}
What Does This Button Do?|KaW

\subsection{Domain Features}
Explorers' society officers gain access to the following domain features.

\paragraph{Optional Feature}
I've Seen Weirder|KaW

\paragraph{Optional Feature}
Research, Research, Research|KaW

\subsection{Mercenary Company}
You're a team of soldiers for hire, fighting the battles no one else wants to. For the right price, of course.

{\par \vspace { 9pt plus 3pt minus 3pt } \noindent  \DndFontTableTitle{Skill and Defense Bonuses} \nopagebreak }

\begin{itemize}
\begin{multicols}{6}
\item \textbf{Operations:} +1
\item \textbf{Resources:} +1
\end{multicols}

\end{itemize}

\subsection{Domain Power}
Mercenary company officers gain access to the following domain power.

\paragraph{Optional Feature}
Fighting Dirty|KaW

\subsection{Domain Features}
Mercenary company officers gain access to the following domain features.

\paragraph{Optional Feature}
Can We Get a Raise?|KaW

\paragraph{Optional Feature}
Survive Till Payday|KaW

\section{Martial Regiment}
Armed to the teeth and with battle strength to match, your characters have come together to form a martial regiment. All of you are skilled warriors, willing to fight the battles others shirk from. And while your group might not always act as the most subtle of organizations, you are well-known for your tactical strategy and efficiency in combat.


\begin{itemize}
\begin{multicols}{2}
\item \textbf{Skills}
\item \textbf{Diplomacy:} +1
\item \textbf{Espionage:} +0
\item \textbf{Lore:} -1
\item \textbf{Operations:} +2
\item \textbf{Defenses}
\item \textbf{Communications:} 12
\item \textbf{Resolve:} 11
\item \textbf{Resources:} 10
\end{multicols}

\end{itemize}

\subsection{Domain Titles}
Martial regiment officers gain access to the following titles, each of which grants a character the noted additional features.


\begin{itemize}
\item Ambush Captain
\item Bravura Commander
\item Field Medic
\item Tactical Marshal
\item War Mage
\end{itemize}

\subsection{Domain Power}
Martial regiment officers gain access to the following domain power.

\paragraph{Optional Feature}
Brute Force|KaW

\subsection{Domains Features}
Martial regiment officers gain access to the following domain features.

\paragraph{Optional Feature}
Armament|KaW

\paragraph{Optional Feature}
The Professionals (Special Unit)|KaW

\subsection{Specialization}
When you found your martial regiment, you'll also choose a specialization. Are you a localized regiment, keeping your homes safe as members of a city watch? Are you members of a military squadron, striking as a crew of highly skilled experts on the most dangerous battlefields? Or are you perhaps a knightly order, united by your dedication to a creed or religion?

\subsection{City Watch}
You might not be the best trained or most efficient of soldiers. But as members of a city watch, you are dedicated to protecting the place you call home—no matter the cost.

{\par \vspace { 9pt plus 3pt minus 3pt } \noindent  \DndFontTableTitle{Skill and Defense Bonuses} \nopagebreak }

\begin{itemize}
\begin{multicols}{6}
\item \textbf{Communications:} +1
\item \textbf{Resources:} +1
\end{multicols}

\end{itemize}

\subsection{Domain Power}
City watch officers gain access to the following domain power.

\paragraph{Optional Feature}
Steel Resolve|KaW

\subsection{Domain Features}
City watch officers gain access to the following domain features.

\paragraph{Optional Feature}
Community Effort|KaW

\paragraph{Optional Feature}
They Will Not Breach This Wall|KaW

\subsection{Knightly Order}
Knightly orders are composed of warriors united under a single ideal. Some have religious ties, while others ascribe to more abstract creeds not bound to any church. Regardless of its members' affiliation, though, every knightly order delineates a certain sense of honor and pride.

{\par \vspace { 9pt plus 3pt minus 3pt } \noindent  \DndFontTableTitle{Skill and Defense Bonuses} \nopagebreak }

\begin{itemize}
\begin{multicols}{6}
\item \textbf{Diplomacy:} +1
\item \textbf{Resolve:} +1
\end{multicols}

\end{itemize}

\subsection{Domain Power}
Knightly order officers gain access to the following domain power.

\paragraph{Optional Feature}
Sworn to Protect|KaW

\subsection{Domain Features}
Knightly order officers gain access to the following domain features.

\paragraph{Optional Feature}
Gallant Company|KaW

\paragraph{Optional Feature}
Stay Strong|KaW

\subsection{Military Squadron}
As members of a military squadron, you are an archetypal team of soldiers. Though you can be discreet when missions call for it, your main focus is always on honing your deadly and precise skills in combat.

{\par \vspace { 9pt plus 3pt minus 3pt } \noindent  \DndFontTableTitle{Skill and Defense Bonuses} \nopagebreak }

\begin{itemize}
\begin{multicols}{6}
\item \textbf{Espionage:} +1
\item \textbf{Operations:} +1
\end{multicols}

\end{itemize}

\subsection{Domain Power}
Military squadron officers gain access to the following domain power.

\paragraph{Optional Feature}
Skirmisher|KaW

\subsection{Domain Features}
Military squadron officers gain access to the following domain features.

\paragraph{Optional Feature}
No Mercy|KaW

\paragraph{Optional Feature}
The War Room Where It Happens|KaW

\section{Mercantile Guild}
Your characters' organization is built on the back of economic prosperity, whether you acquire income from producing, selling, or plundering goods. You might be merchants, artisans, or pirates, looking to expand and defend your empire of industry, and with a network of agents ready to battle for profit.


\begin{itemize}
\begin{multicols}{2}
\item \textbf{Skills}
\item \textbf{Diplomacy:} +1
\item \textbf{Espionage:} +0
\item \textbf{Lore:} -1
\item \textbf{Operations:} +2
\item \textbf{Defenses}
\item \textbf{Communications:} 11
\item \textbf{Resolve:} 10
\item \textbf{Resources:} 12
\end{multicols}

\end{itemize}

\subsection{Domain Titles}
Mercantile guild officers gain access to the following titles, each of which grants a character the noted additional features.


\begin{itemize}
\item Acquisitions Expert
\item Chief of Security
\item Executive Manager
\item Fixer
\item Safety Officer
\end{itemize}

\subsection{Domain Power}
Mercantile guild officers gain access to the following domain power.

\paragraph{Optional Feature}
Outgunned|KaW

\subsection{Domain Features}
Mercantile guild officers gain access to the following domain features.

\paragraph{Optional Feature}
Soldiers of Fortune (Special Unit)|KaW

\paragraph{Optional Feature}
Spared No Expense|KaW

\subsection{Specialization}
When you found your mercantile guild, you'll also choose a specialization. Are you a monopoly, looking to not just dominate but to own an entire industry? Are you a pirate band, sailing the seas looking for loot and plunder? Or perhaps you're a trade guild—a group of artisans banded together to establish the best business deals and working conditions?

\subsection{Monopoly}
Your organization seeks to dominate the market, squashing any resistance you encounter. As such, you invest in making your enterprise the only option, letting your expansionist views force prospective competition to fold into your ranks or fight you to survive. Depending on your market, you might be viewed as an evil corporation or a force of positive change. But at the end of the day, your organization is going to come out on top.

{\par \vspace { 9pt plus 3pt minus 3pt } \noindent  \DndFontTableTitle{Skill and Defense Bonuses} \nopagebreak }

\begin{itemize}
\begin{multicols}{6}
\item \textbf{Espionage:} +1
\item \textbf{Resources:} +1
\end{multicols}

\end{itemize}

\subsection{Domain Power}
Monopoly officers gain access to the following domain power.

\paragraph{Optional Feature}
Action Plan|KaW

\subsection{Domain Features}
Monopoly officers gain access to the following domain features.

\paragraph{Optional Feature}
Corporate Espionage|KaW

\paragraph{Optional Feature}
Embargo|KaW

\subsection{Pirate Band}
You are a pirate band, robbing others of their goods to be resold as you see fit. Though you might ally yourselves with other pirate crews, follow a code of conduct, or work for governments as a privateer, the members of your band are truly beholden only to yourselves. Whatever the reasons you adopted a life of piracy, your organization is viewed with caution wherever you travel.

{\par \vspace { 9pt plus 3pt minus 3pt } \noindent  \DndFontTableTitle{Skill and Defense Bonuses} \nopagebreak }

\begin{itemize}
\begin{multicols}{6}
\item \textbf{Espionage:} +1
\item \textbf{Resolve:} +1
\end{multicols}

\end{itemize}

\subsection{Domain Power}
Pirate band officers gain access to the following domain power.

\paragraph{Optional Feature}
Pillage|KaW

\subsection{Domain Features}
Pirate band officers gain access to the following domain features.

\paragraph{Optional Feature}
Commandeer|KaW

\paragraph{Optional Feature}
Fight Dirty|KaW

\subsection{Trade Guild}
You oversee a network of like-minded artisans who pool their interests together for maximum potential. You might own businesses, or you could facilitate the businesses of others. Whatever the case, you have many skilled individuals investing in your organization to protect their own interests. While a trade guild might seem harmless to a casual observer, the number of resources at your disposal can summon a threat to rival any ruling body.

{\par \vspace { 9pt plus 3pt minus 3pt } \noindent  \DndFontTableTitle{Skill and Defense Bonuses} \nopagebreak }

\begin{itemize}
\begin{multicols}{6}
\item \textbf{Diplomacy:} +1
\item \textbf{Resources:} +1
\end{multicols}

\end{itemize}

\subsection{Domain Power}
Trade guild officers gain access to the following domain power.

\paragraph{Optional Feature}
Healthcare|KaW

\subsection{Domain Features}
Trade guild officers gain access to the following domain features.

\paragraph{Optional Feature}
Business Connections|KaW

\paragraph{Optional Feature}
Well Fed|KaW

\section{Mystic Circle}
Your organization is an association of arcane practitioners whose agents include wizards, sorcerers, scribes, and sages. Not all your associates cast magic—your order needs guards and spies as much as anyone else—but every member of your mystic circle respects the power of the magical arts. The collections of scrolls and tomes in your libraries contain knowledge that others seek, and unearthing the secrets of your enemies isn't so different from uncovering forbidden spells.


\begin{itemize}
\begin{multicols}{2}
\item \textbf{Skills}
\item \textbf{Diplomacy:} +0
\item \textbf{Espionage:} +1
\item \textbf{Lore:} +2
\item \textbf{Operations:} -1
\item \textbf{Defenses}
\item \textbf{Communications:} 12
\item \textbf{Resolve:} 10
\item \textbf{Resources:} 11
\end{multicols}

\end{itemize}

\subsection{Domain Titles}
Mystic circle officers gain access to the following titles, each of which grants a character the noted additional features.


\begin{itemize}
\item Augur
\item Bewitcher
\item Medium
\item Shade
\item Spellsword
\end{itemize}

\subsection{Domain Power}
Mystic circle officers gain access to the following domain power.

\paragraph{Optional Feature}
Universal Energy Field|KaW

\subsection{Domain Features}
Mystic circle officers gain access to the following domain features.

\paragraph{Optional Feature}
Forgotten Rite of Animation (Special Unit)|KaW

\paragraph{Optional Feature}
Unravel Sorcery|KaW

\subsection{Specialization}
When you found your mystic circle, you'll also choose a specialization. Are you an arcane order, tirelessly seeking answers to the mysteries of the universe? Are you a secret cabal, hiding your dark schemes from the world? Or are you a theatrical troupe, using the art of magic to enhance your performances?

\subsection{Arcane Order}
You are a public institution of magic, seeking to expand the understanding of the cosmic forces which ebb and flow through reality. You might be an arcane academy, a great library, or an assembly of notorious mages. Local lords approach you for counsel, unless you prefer to stay neutral in mundane affairs.

{\par \vspace { 9pt plus 3pt minus 3pt } \noindent  \DndFontTableTitle{Skill and Defense Bonuses} \nopagebreak }

\begin{itemize}
\begin{multicols}{6}
\item \textbf{Lore:} +1
\item \textbf{Resources:} +1
\end{multicols}

\end{itemize}

\subsection{Domain Power}
Arcane order officers gain access to the following domain power.

\paragraph{Optional Feature}
Your Staff is Broken|KaW

\subsection{Domain Features}
Arcane order officers gain access to the following domain features.

\paragraph{Optional Feature}
Find True Name|KaW

\paragraph{Optional Feature}
Midnight Oil|KaW

\paragraph{Optional Feature}
Wards of Protection Against Missiles|KaW

\subsection{Secret Cabal}
Your organization is clandestine. People might know you exist, but your work stays hidden to protect your members, keep your studies confidential, or better allow you to pull strings from behind the scenes. No matter what the reasons for your secrecy, one thing is certain: You look out for your own first and foremost.

{\par \vspace { 9pt plus 3pt minus 3pt } \noindent  \DndFontTableTitle{Skill and Defense Bonuses} \nopagebreak }

\begin{itemize}
\begin{multicols}{6}
\item \textbf{Espionage:} +1
\item \textbf{Communications:} +1
\end{multicols}

\end{itemize}

\subsection{Domain Power}
Secret cabal officers gain access to the following domain power.

\paragraph{Optional Feature}
Elemental Bargain|KaW

\subsection{Domain Features}
Secret cabal officers gain access to the following domain features.

\paragraph{Optional Feature}
Curse of Vitiation|KaW

\paragraph{Optional Feature}
We Were Never Here|KaW

\subsection{Theatrical Troupe}
Magic and artifice have always gone hand in hand, and the art of spellcasting makes for some impressive stage effects. Whether you are a performance company with an established base of operations or a band of entertainers travelling far and wide, you endeavor to make magic accessible to all! Your organization doesn't believe in keeping sorcery hidden away in secluded towers, underground lairs, or exclusive clubs. Showing magic to the world encourages others to take up the art, dispels fears of spellcasting, and can bring joy to the common folk.

{\par \vspace { 9pt plus 3pt minus 3pt } \noindent  \DndFontTableTitle{Skill and Defense Bonuses} \nopagebreak }

\begin{itemize}
\begin{multicols}{6}
\item \textbf{Diplomacy:} +1
\item \textbf{Operations:} +1
\end{multicols}

\end{itemize}

\subsection{Domain Power}
Theatrical troupe officers gain access to the following domain power.

\paragraph{Optional Feature}
Magic Misdirection|KaW

\subsection{Domain Features}
Theatrical troupe officers gain access to the following domain features.

\paragraph{Optional Feature}
Combat Anthem|KaW

\paragraph{Optional Feature}
You Didn't Hear This From Me|KaW

\section{Nature Pact}
Your organization lives and thrives in the untamed wilds, whether forests, deserts, tundra, or even oceans. Nature cannot be controlled, but you know how to channel primal magic. Barbarians, druids, rangers, beasts, and nature spirits are your domain's agents, and with these forces on your side, you are unstoppable.


\begin{itemize}
\begin{multicols}{2}
\item \textbf{Skills}
\item \textbf{Diplomacy:} +1
\item \textbf{Espionage:} +0
\item \textbf{Lore:} +2
\item \textbf{Operations:} -1
\item \textbf{Defenses}
\item \textbf{Communications:} 10
\item \textbf{Resolve:} 12
\item \textbf{Resources:} 11
\end{multicols}

\end{itemize}

\subsection{Domain Titles}
Nature pact officers gain access to the following titles, each of which grants a character the noted additional features.


\begin{itemize}
\item The Connected
\item The Pack Leader
\item The Primal
\item The Speaker
\item The Stalwart
\end{itemize}

\subsection{Domain Power}
Nature pact officers gain access to the following domain power.

\paragraph{Optional Feature}
Vine Entrapment|KaW

\subsection{Domain Features}
Nature pact officers gain access to the following domain features.

\paragraph{Optional Feature}
Frog of War (Special Unit)|KaW

\paragraph{Optional Feature}
Natural Disaster|KaW

\subsection{Spcialization}
When you found your nature pact, you'll also choose a specialization. Are you a barbarian tribe, escaping the conventions and customs of settled life by making your own rules in the wild? Are you a druid circle, living as an extension of the natural world and unable to fathom a home anywhere else? Or are you a hunter conclave defending nature from evil's corruption and industry's greed?

\subsection{Barbarian Tribe}
Power, ferocity, and strength define the barbarian warriors who make up your organization, allowing you to push through enemies and knock all obstacles out of your path.

{\par \vspace { 9pt plus 3pt minus 3pt } \noindent  \DndFontTableTitle{Skill and Defense Bonuses} \nopagebreak }

\begin{itemize}
\begin{multicols}{6}
\item \textbf{Operations:} +1
\item \textbf{Resolve:} +1
\end{multicols}

\end{itemize}

\subsection{Domain Power}
Barbarian tribe officers gain access to the following domain power.

\paragraph{Optional Feature}
Impenetrable Defense|KaW

\subsection{Domain Features}
Barbarian tribe officers gain access to the following domain features.

\paragraph{Optional Feature}
Ready to Slay|KaW

\paragraph{Optional Feature}
Trial by Pyre|KaW

\subsection{Druid Circle}
As the leaders of a druid circle, you gain the allegiance of beasts, plants, and even the elements to your cause. Your agents learn how to use terrain to their advantage, hiding in plain sight and creating an information network unlike any other.

{\par \vspace { 9pt plus 3pt minus 3pt } \noindent  \DndFontTableTitle{Skill and Defense Bonuses} \nopagebreak }

\begin{itemize}
\begin{multicols}{6}
\item \textbf{Espionage:} +1
\item \textbf{Lore:} +1
\end{multicols}

\end{itemize}

\subsection{Domain Power}
Druid circle officers gain access to the following domain power.

\paragraph{Optional Feature}
Primal Conjuration|KaW

\subsection{Domain Features}
Druid circle officers gain access to the following domain features.

\paragraph{Optional Feature}
Eagle Eye|KaW

\paragraph{Optional Feature}
Torrential Terrain|KaW

\subsection{Hunter Conclave}
You lead a hunter conclave, whose members are dedicated to taking on aberrations, monstrosities, undead, and other unnatural creatures that threaten the wilds.

{\par \vspace { 9pt plus 3pt minus 3pt } \noindent  \DndFontTableTitle{Skill and Defense Bonuses} \nopagebreak }

\begin{itemize}
\begin{multicols}{6}
\item \textbf{Diplomacy:} +1
\item \textbf{Espionage:} +1
\end{multicols}

\end{itemize}

\subsection{Domain Power}
Hunter conclave officers gain access to the following domain power.

\paragraph{Optional Feature}
Rapid Assault|KaW

\subsection{Domain Features}
Hunter conclave officers gain access to the following domain features.

\paragraph{Optional Feature}
Marked Targets|KaW

\paragraph{Optional Feature}
Ruinous Fire|KaW

\section{Noble Court}
A noble court domain is unique among the organizations in this section, as it allows the characters to take responsibility for governing a land and the citizens who live and work there. (And if your noble court reaches size 5, it could become a kingdom! Yay, it's the title of the book!) Farmers, smiths, carpenters, and shepherds might be the agents of a rural noble court, while clockmakers, luthiers, butchers, or musicians might serve you in a city. Whoever your followers are, they all look to you for guidance and protection.


\begin{itemize}
\begin{multicols}{2}
\item \textbf{Skills}
\item \textbf{Diplomacy:} +2
\item \textbf{Espionage:} +0
\item \textbf{Lore:} -1
\item \textbf{Operations:} +1
\item \textbf{Defenses}
\item \textbf{Communications:} 10
\item \textbf{Resolve:} 12
\item \textbf{Resources:} 11
\end{multicols}

\end{itemize}

\subsection{Domain Titles}
Noble court officers gain access to the following titles, each of which grants a character the noted additional features.


\begin{itemize}
\item Court Mage
\item Court Minstrel
\item High Priest
\item Master Assassin
\item Master-at-Arms
\end{itemize}

\subsection{Domain Power}
Noble court officers gain access to the following domain power.

\paragraph{Optional Feature}
Mantle of Authority|KaW

\subsection{Domain Features}
Noble court officers gain access to the following domain features.

\paragraph{Optional Feature}
Prepare the Steeds (Special Unit)|KaW

\paragraph{Optional Feature}
Skilled Negotiators|KaW

\subsection{Specialization}
When you found your noble court, you'll also choose a specialization. Are you a court of war, focused on growing your holdings and resources through military might? Are you a political administration, engaging in trade and detente to provide for your people? Or are you a regent state, honoring the traditions of your land and the family that leads it?

\subsection{Court of War}
Your noble court focuses on conquest, seeking growth through martial power. Like the subjects of so many great empires, your people are kept happy through the prosperity gained by annexing new territory. But those you conquer might view your actions as those of a tyrant rather than a hero.

{\par \vspace { 9pt plus 3pt minus 3pt } \noindent  \DndFontTableTitle{Skill and Defense Bonuses} \nopagebreak }

\begin{itemize}
\begin{multicols}{6}
\item \textbf{Operations:} +1
\item \textbf{Resolve:} +1
\end{multicols}

\end{itemize}

\subsection{Domain Power}
Court of war officers gain access to the following domain power.

\paragraph{Optional Feature}
Conqueror|KaW

\subsection{Domain Features}
Court of war officers gain access to the following domain features.

\paragraph{Optional Feature}
Forewarned is Forearmed|KaW

\paragraph{Optional Feature}
Outmaneuvered|KaW

\subsection{Political Administration}
Your organization is focused on the skills and traditions of statecraft. Its officers are politicians first, pulling strings and manipulating neighbors, gathering favors and making deals. You prefer to use your diplomatic skills to avoid war—but if that proves impossible, you ensure that the battle happens on your terms, surrounded by allies while your enemy stands alone.

{\par \vspace { 9pt plus 3pt minus 3pt } \noindent  \DndFontTableTitle{Skill and Defense Bonuses} \nopagebreak }

\begin{itemize}
\begin{multicols}{6}
\item \textbf{Diplomacy:} +2
\end{multicols}

\end{itemize}

\subsection{Domain Power}
Political administration officers gain access to the following domain power.

\paragraph{Optional Feature}
Timely Aid|KaW

\subsection{Domain Features}
Political administration officers gain access to the following domain features.

\paragraph{Optional Feature}
Diplomacy for Intel|KaW

\paragraph{Optional Feature}
Enchantment Economy|KaW

\subsection{Regent State}
A regent holds power in trust to someone else, usually an absent monarch or an heir too young to rule. But a regent can also hold power in trust to the people, always working to act in their best interests. Regents concern themselves first with their people's well-being, and work hard to establish justice, prosperity, and security.

{\par \vspace { 9pt plus 3pt minus 3pt } \noindent  \DndFontTableTitle{Skill and Defense Bonuses} \nopagebreak }

\begin{itemize}
\begin{multicols}{6}
\item \textbf{Resolve:} +1
\item \textbf{Resources:} +1
\end{multicols}

\end{itemize}

\subsection{Domain Power}
Regent state officers gain access to the following domain power.

\paragraph{Optional Feature}
Voices of the Past|KaW

\subsection{Domain Features}
Regent state officers gain access to the following domain features.

\paragraph{Optional Feature}
Backdoor Negotiations|KaW

\paragraph{Optional Feature}
Volunteer Blacksmiths|KaW

\section{Religious Order}
Your characters have formed a religious order, assembled in service to a higher power such as a god, archfiend, or unfathomable elder entity. You are no stranger to poring over ancient scriptures, and you can mobilize your tight-knit community of acolytes, priests, and champions to care for the poor and vulnerable, defend the faith, or conduct occult rituals. Whether congregating beneath the scintillating glow of stained glass windows or engaging in lost rites before ancient altars, your faith is strong, and you know that through prayer, all things are possible.


\begin{itemize}
\begin{multicols}{2}
\item \textbf{Skills}
\item \textbf{Diplomacy:} +0
\item \textbf{Espionage:} -1
\item \textbf{Lore:} +2
\item \textbf{Operations:} +1
\item \textbf{Defenses}
\item \textbf{Communications:} 11
\item \textbf{Resolve:} 12
\item \textbf{Resources:} 10
\end{multicols}

\end{itemize}

\subsection{Domain Titles}
Religious order officers gain access to the following titles, each of which grants a character the noted additional features.


\begin{itemize}
\item Conduit
\item Crusader
\item Herald
\item Omen
\item Prophet
\end{itemize}

\subsection{Domain Power}
Religious order officers gain access to the following domain power.

\paragraph{Optional Feature}
Beseech the Most High|KaW

\begin{DndSidebar}[float=!b]{Freedom of Religion}
The nature or specialization of your religious order has no bearing on which type of planar creature— aberration, celestial, or fiend—you can summon with the Beseech the Most High power. Think outside the box! Your holy church might petition a righteous deity for aid and be granted command of a monstrous devil bound by hallowed chains. Likewise, your hidden cult might worship an archdevil that torments souls in an eldritch demiplane, conjuring an aberration onto the battlefield from a portal of writhing tentacles.



\end{DndSidebar}

\subsection{Domain Features}
Religious order officers gain access to the following domain features.

\paragraph{Optional Feature}
Consult the Scriptures|KaW

\paragraph{Optional Feature}
Spread the Holy Word (Special Unit)|KaW

\subsection{Specialization}
When you found your religious order, you'll also choose a specialization. Are your rituals taboo or even illegal, forcing you to operate from the shadows as a hidden cult? Does your god command you to grow their faith by operating as a holy church, sending missionaries throughout the land? Or perhaps you're part of a monastic order, guarding ancient secrets in some secluded domain of faith?

\subsection{Hidden Cult}
Society is rife with close-minded individuals who disapprove of your causes or your methods, so you must conduct your rituals in secret. The agents of your hidden cult are everyday people—acolytes and blacksmiths, guards and government officials—who mask their true natures behind ceremonial garb and closed-door meetings. Your stronghold is likely a front for occult activity, and your allies are won through temptation, manipulation, and intimidation. And though you try to keep a low profile, conflict is sometimes inevitable.

{\par \vspace { 9pt plus 3pt minus 3pt } \noindent  \DndFontTableTitle{Skill and Defense Bonuses} \nopagebreak }

\begin{itemize}
\begin{multicols}{6}
\item \textbf{Espionage:} +1
\item \textbf{Lore:} +1
\end{multicols}

\end{itemize}

\subsection{Domain Power}
Hidden cult officers gain access to the following domain power.

\paragraph{Optional Feature}
Penance|KaW

\subsection{Domain Features}
Hidden cult officers gain access to the following domain features.

\paragraph{Optional Feature}
Blood Sacrifice|KaW

\paragraph{Optional Feature}
Poison the Well|KaW

\subsection{Holy Church}
As a devout group of worshipers, you act as the instruments of a god. Your agents are dedicated and organized, proudly proclaiming their faith to would-be converts and constructing gleaming temples to receive them. Your clothing or armor is emblazoned with the holy symbol of your deity as a reminder that your cause has been ordained by the ultimate authority.

{\par \vspace { 9pt plus 3pt minus 3pt } \noindent  \DndFontTableTitle{Skill and Defense Bonuses} \nopagebreak }

\begin{itemize}
\begin{multicols}{6}
\item \textbf{Diplomacy:} +1
\item \textbf{Resources:} +1
\end{multicols}

\end{itemize}

\subsection{Domain Power}
Holy church officers gain access to the following domain power.

\paragraph{Optional Feature}
Turn the Tide|KaW

\subsection{Domain Features}
Holy church officers gain access to the following domain features.

\paragraph{Optional Feature}
Having the Gods on Our Side|KaW

\paragraph{Optional Feature}
Resplendent Armor|KaW

\subsection{Monastic Order}
As a monastic order, you create a secluded community devoted to a cause, doctrine, or philosophy. Your agents are servants, scholars, and sages who renounced the world in favor of a nobler path. Monasticism is a way of life, full of strict regimens designed to hone the body, mind, and spirit. But though you spend your days in peaceful meditation, your order's arduous training makes you a dangerous weapon.

{\par \vspace { 9pt plus 3pt minus 3pt } \noindent  \DndFontTableTitle{Skill and Defense Bonuses} \nopagebreak }

\begin{itemize}
\begin{multicols}{6}
\item \textbf{Lore:} +1
\item \textbf{Operations:} +1
\end{multicols}

\end{itemize}

\subsection{Domain Power}
Monastic order officers gain access to the following domain power.

\paragraph{Optional Feature}
Quaking Fist|KaW

\subsection{Domain Features}
Monastic order officers gain access to the following domain features.

\paragraph{Optional Feature}
Ancient Archives|KaW

\paragraph{Optional Feature}
Iron Will|KaW

\section{Underworld Syndicate}
Your organization is an underworld syndicate acting from the shadows, running agents who might be thugs, spies, killers, or cutpurses. The resources you control put you in a perfect position to undertake clandestine activities, morally ambiguous deals, and jobs not for the faint of heart. You go in, you get the job done, and you get out again, all without anyone knowing. Or, if you're especially good at what you do, you'll leave everyone thinking someone else was responsible.


\begin{itemize}
\begin{multicols}{2}
\item \textbf{Skills}
\item \textbf{Diplomacy:} +0
\item \textbf{Espionage:} +2
\item \textbf{Lore:} -1
\item \textbf{Operations:} +1
\item \textbf{Defenses}
\item \textbf{Communications:} 12
\item \textbf{Resolve:} 10
\item \textbf{Resources:} 11
\end{multicols}

\end{itemize}

\subsection{Domain Titles}
Underworld syndicate officers gain access to the following titles, each of which grants a character the noted additional features.


\begin{itemize}
\item Enforcement
\item Narcotics
\item Negotiations
\item Operations
\item Records
\end{itemize}

\subsection{Domain Power}
Underworld syndicate officers gain access to the following domain power.

\paragraph{Optional Feature}
Find Weakness|KaW

\subsection{Domain Features}
Underworld syndicate officers gain access to the following domain features.

\paragraph{Optional Feature}
Plans within Plans|KaW

\paragraph{Optional Feature}
The Crew (Special Unit)|KaW

\subsection{Specialization}
When you found your underworld syndicate, you'll also choose a specialization. Are you a venerable and prestigious assassins' college, training a proud line of professional killers? Are you an official spy network, your existence known to the general public even as your work is shrouded in mystery? Or are you a down-and-dirty thieves' guild, taking on any jobs for the right price?

\subsection{Assassins' College}
Assassination is a most misunderstood art. You kill people for money, sure—but only certain people and only under certain circumstances. There are rules, after all, and new assassins must learn those rules at the hands of masters.

{\par \vspace { 9pt plus 3pt minus 3pt } \noindent  \DndFontTableTitle{Skill and Defense Bonuses} \nopagebreak }

\begin{itemize}
\begin{multicols}{6}
\item \textbf{Espionage:} +1
\item \textbf{Resources:} +1
\end{multicols}

\end{itemize}

\subsection{Domain Power}
Assassins' college officers gain access to the following domain power.

\paragraph{Optional Feature}
Assassin's Strike|KaW

\subsection{Domain Features}
Assassins' college officers gain access to the following domain features.

\paragraph{Optional Feature}
Looks Like We Need a New Boss|KaW

\paragraph{Optional Feature}
The Best There Is|KaW

\subsection{Spy Network}
As part of a covert team of spies, you gather intelligence and engage in clandestine activities for a government, a specific organization—or the highest bidder. Your agents are well trained and savvy, able to operate in even the most challenging conditions.

{\par \vspace { 9pt plus 3pt minus 3pt } \noindent  \DndFontTableTitle{Skill and Defense Bonuses} \nopagebreak }

\begin{itemize}
\begin{multicols}{6}
\item \textbf{Diplomacy:} +1
\item \textbf{Espionage:} +1
\end{multicols}

\end{itemize}

\subsection{Domain Power}
Spy network officers gain access to the following domain power.

\paragraph{Optional Feature}
Traitor|KaW

\subsection{Domain Features}
Spy network officers gain access to the following domain features.

\paragraph{Optional Feature}
Create Vulnerability|KaW

\paragraph{Optional Feature}
False Orders|KaW

\subsection{Thieves' Guild}
Your organization is a second-story crew, focusing on getting into where you don't belong, taking what isn't yours, and earning your well-deserved rewards. You're not much for politics, spying, and information brokering—except when it helps you cut deals with the law.

{\par \vspace { 9pt plus 3pt minus 3pt } \noindent  \DndFontTableTitle{Skill and Defense Bonuses} \nopagebreak }

\begin{itemize}
\begin{multicols}{6}
\item \textbf{Espionage:} +1
\item \textbf{Operations:} +1
\end{multicols}

\end{itemize}

\subsection{Domain Power}
Thieves' guild officers gain access to the following domain power.

\paragraph{Optional Feature}
Poison Weapons|KaW

\subsection{Domain Features}
Thieves' guild officers gain access to the following domain features.

\paragraph{Optional Feature}
Stolen Supplies|KaW

\paragraph{Optional Feature}
They're Going to Have a Hard Time Making Payroll This Week|KaW

\section{NPC Realms}
When a domain is controlled by an NPC—whether a villain, an ally, or a mystery figure whose role the characters need to figure out—the GM sets that domain up as an NPC realm. This section presents sixteen NPC realms for the GM's use in a campaign, covering a wide array of domain types, from the explicitly villainous (it's hard to imagine how a despotic regime could be a good thing) to neutral and possibly friendly realms such as a fey court or a gnomish kingdom.

Unlike heroic organizations, \textbf{each NPC realm has only one title, given to the leader}. This reflects the fact that all the officers of a villainous realm (and any other domain outside the characters' domain) are run by the GM, and no GM needs to worry about five NPCs with special abilities! Also unlike organizations, NPC realms have no specialization. They have the same number of domain powers and domain features overall, but those powers and features are consistent for each type of realm. This allows a GM to quickly copy the details for an NPC realm onto the generic domain sheet (page 315) and be ready to go.

Some of an NPC realm's domain powers can be taken only by the domain's leader. Other powers can be taken by the leader or other officers.

There's no rule that says all the NPC domains in a game must be realms from this section. GMs can use the domains in the \textit{Heroic Organizations} section as well when populating a world with NPC domains. But by keeping the setup for these realms simpler and less customizable, the intent is to make them easier for the GM to use.

\section{Despotic Regime}
A despotic warlord rules through fear and propaganda. They are not interested in the health or well-being of their people, but only in making sure their army is always prepared for battle. Despots are cruel. Conquest is the tool they use to appease their soldiers, but power is all they care about. If it ever seems as though a despot is being reasonable during negotiations, this is a ruse. They are always scheming, even against their allies. And if they defer to a greater power, it is only because a despot fears anyone with more power than them.


\begin{itemize}
\begin{multicols}{2}
\item \textbf{Skills}
\item \textbf{Diplomacy:} -1
\item \textbf{Espionage:} +2
\item \textbf{Lore:} +0
\item \textbf{Operations:} +2
\item \textbf{Defenses}
\item \textbf{Communications:} 11
\item \textbf{Resolve:} 12
\item \textbf{Resources:} 10
\end{multicols}

\end{itemize}

\subsection{Domain Title}
The despotic regime leader gains the following title and the noted additional features.

\paragraph{Optional Feature}
Despot|KaW

\subsection{Domain Powers}
The despotic regime leader gains access to the following domain powers.

\paragraph{Optional Feature}
Make an Example|KaW

\paragraph{Optional Feature}
My Life for Yours|KaW

\subsection{Domain Features}
Despotic regime officers gain access to the following domain features.

\paragraph{Optional Feature}
Death to Spies|KaW

\paragraph{Optional Feature}
Readiness is All|KaW

\paragraph{Optional Feature}
The Crimson Guard (Special Unit)|KaW

\paragraph{Optional Feature}
Work the Prisoners to the Bone|KaW

\section{Draconic Empire}
Some dragons amass gold and jewels; others covet the wealth of magic. But the most formidable and cunning of dragons hoard whole kingdoms. A draconic empire is the end result of such unparalleled greed: an amalgamation of cities and lands brought together under a single dragon's rapacious purview. With a lifespan of over a thousand years, a draconic ruler maintains an iron hold over their realm, crushing any attempt to conquer or disband their empire with ease.

Though dragons and other creatures such as wyverns are found among a draconic empire's subjects, the empire absorbs many by way of conquest, and not every citizen of such a realm has draconic blood. At the same time, not every draconic sovereign is seen as a tyrant, and some of the empire's folk might even prefer the efficient and judicious rulings of a dragon overlord compared to the alternatives.


\begin{itemize}
\begin{multicols}{2}
\item \textbf{Skills}
\item \textbf{Diplomacy:} -1
\item \textbf{Espionage:} +0
\item \textbf{Lore:} +2
\item \textbf{Operations:} +2
\item \textbf{Defenses}
\item \textbf{Communications:} 10
\item \textbf{Resolve:} 11
\item \textbf{Resources:} 13
\end{multicols}

\end{itemize}

\subsection{Domain Title}
The draconic empire leader gains the following title and the noted additional features.

\paragraph{Optional Feature}
Imperator Draconis|KaW

\subsection{Domain Powers}
Draconic empire officers gain access to the following domain powers.

\paragraph{Optional Feature}
Glittering Armor|KaW

\paragraph{Optional Feature}
The Great and Terrible|KaW

\subsection{Domain Features}
Draconic empire officers gain access to the following domain features.

\paragraph{Optional Feature}
Pillage|KaW

\paragraph{Optional Feature}
Destruction Rains from Above|KaW

\paragraph{Optional Feature}
Furnishings of the Hoard|KaW

\paragraph{Optional Feature}
Praetores Draconis (Special Unit)|KaW

\section{Dwarven Thanedom}
Not all dwarven cultures are bellicose, but the dwarves' history of stalwart bravery lends them several advantages when a dwarven thanedom is drawn into conflict—or when its folk decide to start trouble themselves. Dwarven physiology and hardiness is well suited to land war. But even more importantly, many dwarves possess a strong analytical streak that underlies their success as crafters, engineers, and miners. With an instinctive sense of how various parts work together to create a machine, dwarves called to battle understand how to become a part in a war machine, making dwarf soldiers a force to be reckoned with on the field.


\begin{itemize}
\begin{multicols}{2}
\item \textbf{Skills}
\item \textbf{Diplomacy:} +0
\item \textbf{Espionage:} -1
\item \textbf{Lore:} +1
\item \textbf{Operations:} +2
\item \textbf{Defenses}
\item \textbf{Communications:} 11
\item \textbf{Resolve:} 12
\item \textbf{Resources:} 12
\end{multicols}

\end{itemize}

\subsection{Domain Title}
The dwarven thanedom leader gains the following title and the noted additional features.

\paragraph{Optional Feature}
Thane|KaW

\subsection{Domain Powers}
Dwarven thanedom officers gain access to the following domain powers.

\paragraph{Optional Feature}
Iron Tide|KaW

\paragraph{Optional Feature}
Natural Sprinter|KaW

\subsection{Domain Features}
Dwarven thanedom officers gain access to the following domain features.

\paragraph{Optional Feature}
Accelerated Industry|KaW

\paragraph{Optional Feature}
Blazing Forges|KaW

\paragraph{Optional Feature}
Siege Manufacturing|KaW

\paragraph{Optional Feature}
Warpath|KaW

\section{Fey Court}
The courts of the elves, the fairies who serve them, and the fey creatures of the wode are alien to most people of the world. Cruel and capricious, the fey court's rulers are driven by motivations unknowable. They are older than the world, and remember magics that mortal creatures were not meant to know. Draw the ire of a fey court, and the shadows you see, the sounds you hear, the world you touch could well become an illusion designed to lead you down a long, torturous road of doom.


\begin{itemize}
\begin{multicols}{2}
\item \textbf{Skills}
\item \textbf{Diplomacy:} +0
\item \textbf{Espionage:} +1
\item \textbf{Lore:} +3
\item \textbf{Operations:} -1
\item \textbf{Defenses}
\item \textbf{Communications:} 11
\item \textbf{Resolve:} 12
\item \textbf{Resources:} 11
\end{multicols}

\end{itemize}

\subsection{Domain Title}
The fey court leader gains the following title and the noted additional features.

\paragraph{Optional Feature}
Monarch of Mischief|KaW

\subsection{Domain Powers}
Fey court officers gain access to the following domain powers.

\paragraph{Optional Feature}
Now You See Me...and Me...and Me!|KaW

\paragraph{Optional Feature}
Watch Where You Swing|KaW

\subsection{Domain Features}
Fey court officers gain access to the following domain features.

\paragraph{Optional Feature}
Disillusioned Illusions|KaW

\paragraph{Optional Feature}
Fairy Circles|KaW

\paragraph{Optional Feature}
Surely You Jest (Special Unit)|KaW

\paragraph{Optional Feature}
Your Own Worst Enemy|KaW

\section{Giant Jarldom}
In warfare, size matters. Giant jarls often have little patience for ingenuity or diplomatic negotiations, preferring to throw their weight around and squash any puny fools who dare to challenge their will. In the eyes of a living siege weapon, most fortifications constructed by the smallfolk are laughable at best. Giant jarldoms are living engines of war, ready for battle at a moment's notice. Jarls can live hundreds of years, none of which are spent suffering fickle alliances, usurpers, or smallfolk kingdoms who've forgotten their place. The mere touch of a giant's barrel-sized fist can topple a soldier in full plate armor, and the tremors of their footsteps foreshadow the destruction that comes in their wake.


\begin{itemize}
\begin{multicols}{2}
\item \textbf{Skills}
\item \textbf{Diplomacy:} +0
\item \textbf{Espionage:} -1
\item \textbf{Lore:} +2
\item \textbf{Operations:} +2
\item \textbf{Defenses}
\item \textbf{Communications:} 10
\item \textbf{Resolve:} 13
\item \textbf{Resources:} 11
\end{multicols}

\end{itemize}

\subsection{Domain Title}
The giant jarldom leader gains the following title and the noted additional features.

\paragraph{Optional Feature}
Jarl|KaW

\subsection{Domain Powers}
Giant jarldom officers gain access to the following domain powers.

\paragraph{Optional Feature}
The Bigger I Am, the Harder You Fall|KaW

\paragraph{Optional Feature}
Thick Hide|KaW

\subsection{Domain Features}
Giant jarldom officers gain access to the following domain features.

\paragraph{Optional Feature}
Bring the Big Rocks|KaW

\paragraph{Optional Feature}
Indomitable Might|KaW

\paragraph{Optional Feature}
Runic Omens|KaW

\paragraph{Optional Feature}
There Are Giants in the Sky (Special Unit)|KaW

\section{Gnomish Kingdom}
Bustling with life, gnomish kingdoms are centers for arcane study and engineering breakthroughs. Here, elderly gnomes pass down stories learned across long life spans to bright-eyed young creators, while clever inventors in underground laboratories chip away at new experiments. Gnomes are often known for innovations that work in tandem with the earth, not against it, and their kingdoms are all the richer for it.

Gnomish kingdoms are home to many other types of creatures as well, including students from far-flung lands seeking lore and training. Elderly gnome engineers take on such apprentices happily, ready to share the tricks of their craft. But don't mistake this friendliness for weakness. When pushed to war, gnomish kingdoms conduct themselves with the same precision and ingenuity that their citizens display in their workshops. Their infantry can cleverly dismantle a fortress in a day, and their artillery squadrons wield powerful firearms unlike any other.


\begin{itemize}
\begin{multicols}{2}
\item \textbf{Skills}
\item \textbf{Diplomacy:} +0
\item \textbf{Espionage:} -1
\item \textbf{Lore:} +2
\item \textbf{Operations:} +2
\item \textbf{Defenses}
\item \textbf{Communications:} 12
\item \textbf{Resolve:} 10
\item \textbf{Resources:} 12
\end{multicols}

\end{itemize}

\subsection{Domain Title}
The gnomish kingdom leader gains the following title and the noted additional features.

\paragraph{Optional Feature}
Prime Sophist|KaW

\subsection{Domain Powers}
Gnomish kingdom officers gain access to the following domain powers.

\paragraph{Optional Feature}
Never Underestimate the Small|KaW

\paragraph{Optional Feature}
Thinking Outside the Box|KaW

\subsection{Domain Features}
Gnomish kingdom officers gain access to the following domain features.

\paragraph{Optional Feature}
Clever Together|KaW

\paragraph{Optional Feature}
Mirror Legion (Special Unit)|KaW

\paragraph{Optional Feature}
Raw Firepower|KaW

\paragraph{Optional Feature}
Warfare Engineers|KaW

\section{Goblinoid Coalition}
Most domains consider enemy goblinoid realms a nuisance rather than a true military threat. Outside of urban environments, goblins, bugbears, and hobgoblins usually live in small clan groups, rising and falling in strength as the strongest leaders jockey for power. But sometimes an ambitious goblinoid leader is able to unite a number of clans as a singular community—and a dangerous fighting force. By pooling resources, goblinoid clans and bands can form a coalition capable of launching devastating attacks on all who threaten their lands and people, bringing goblins, hobgoblins, and bugbears into a unified and disciplined force of capable warriors.


\begin{itemize}
\begin{multicols}{2}
\item \textbf{Skills}
\item \textbf{Diplomacy:} +0
\item \textbf{Espionage:} +1
\item \textbf{Lore:} -1
\item \textbf{Operations:} +3
\item \textbf{Defenses}
\item \textbf{Communications:} 12
\item \textbf{Resolve:} 10
\item \textbf{Resources:} 12
\end{multicols}

\end{itemize}

\subsection{Domain Title}
The goblinoid coalition leader gains the following title and the noted additional features.

\paragraph{Optional Feature}
King Boss|KaW

\subsection{Domain Powers}
Gnomish kingdom officers gain access to the following domain powers.

\paragraph{Optional Feature}
Cheap Shot|KaW

\paragraph{Optional Feature}
Worg Pack|KaW

\subsection{Domain Features}
Gnomish kingdom officers gain access to the following domain features.

\paragraph{Optional Feature}
Mumbo Jumbo|KaW

\paragraph{Optional Feature}
Prisoners of War|KaW

\paragraph{Optional Feature}
Spike Pits|KaW

\paragraph{Optional Feature}
The Red Howl (Special Unit)|KaW

\section{Hag Coven}
A hag's deception and vile reckoning makes for a dangerous and unpredictable foe. Every encounter with one of these fiends is a gamble that might deal results worse than death. When several hags form a hag coven, their combined magical, seductive, and scheming power increases exponentially. Patient and cruel, the members of a hag coven sign pacts with the most desperate mortals, using the might of the horrid fey, fiends, and monstrosities that follow them as an army to enforce their dark bargains against even war leaders and emperors. Be wary when facing or allying with a coven, for every move the hags make furthers their singular goal of seeing the world suffer.


\begin{itemize}
\begin{multicols}{2}
\item \textbf{Skills}
\item \textbf{Diplomacy:} +2
\item \textbf{Espionage:} +0
\item \textbf{Lore:} +2
\item \textbf{Operations:} -1
\item \textbf{Defenses}
\item \textbf{Communications:} 12
\item \textbf{Resolve:} 12
\item \textbf{Resources:} 10
\end{multicols}

\end{itemize}

\subsection{Domain Title}
The goblinoid coalition leader gains the following title and the noted additional features.

\paragraph{Optional Feature}
Queen of Night|KaW

\subsection{Domain Powers}
Gnomish kingdom officers gain access to the following domain powers.

\paragraph{Optional Feature}
Curse of the Coven|KaW

\paragraph{Optional Feature}
Know Thy Place|KaW

\subsection{Domain Features}
Gnomish kingdom officers gain access to the following domain features.

\paragraph{Optional Feature}
A Comforting Voice|KaW

\paragraph{Optional Feature}
Clever Bargain|KaW

\paragraph{Optional Feature}
Drained Source|KaW

\paragraph{Optional Feature}
Swamp's Wrath (Special Unit)|KaW

\section{Infernal Echelon}
Fiends are evil incarnate. Whether they seek to cheat mortals out of their souls or simply consume those souls by force, all fiends share the same drive—the suffering of the living. It's easy to lose sight of this simple truth when dealing with these creatures, as both the Abyss and the Nine Hells are entrenched in complicated politics. Demons and devils both follow strict hierarchies, creating complex power struggles whose machinations have consequences that ripple out across multiple worlds. And at the heart of these grand plots stands the infernal echelon—an organized force of fiends dedicated to powerful leaders, and committed to unleashing death, madness, and despair upon the mortal realm.


\begin{itemize}
\begin{multicols}{2}
\item \textbf{Skills}
\item \textbf{Diplomacy:} +2
\item \textbf{Espionage:} +1
\item \textbf{Lore:} +1
\item \textbf{Operations:} +0
\item \textbf{Defenses}
\item \textbf{Communications:} 12
\item \textbf{Resolve:} 10
\item \textbf{Resources:} 11
\end{multicols}

\end{itemize}

\subsection{Domain Titles}
The infernal echelon leader gains one of the following titles (based on their particular flavor of evil), each of which grants the noted additional features.

\paragraph{Optional Feature}
Deceiver|KaW

\paragraph{Optional Feature}
Defiler|KaW

\subsection{Domain Powers}
Infernal echelon officers gain access to the following domain powers.

\paragraph{Optional Feature}
Doom Curse|KaW

\paragraph{Optional Feature}
Hellfire|KaW

\subsection{Domain Features}
Infernal echelon officers gain access to the following domain features.

\paragraph{Optional Feature}
Face Dancers|KaW

\paragraph{Optional Feature}
Inferno Artillery|KaW

\paragraph{Optional Feature}
Led by the Nose|KaW

\paragraph{Optional Feature}
Summon the Pale Riders (Special Unit)|KaW

\section{Medusean Tyranny}
Those afflicted with the medusa's curse often retreat into isolation, sequestering themselves in remote lairs to hide their monstrous transformation. Most medusas spend their days collecting not just victims but knowledge, as they attempt to undo the fell magic that has transformed them. A medusean tyranny arises when these powerful beings emerge from the shadows to seek out others of their kind. Alone, a medusa is a terrifying threat. Together, they pool centuries of knowledge and power, becoming all but unstoppable.

Medusas working together form a democratic unit known as a Veiled Council, with a medusean tyranny coming into full power when a council seizes a first stronghold. Once a base of operations is established, the council begins to expand across surrounding regions, ruling through fear and squashing resistance with powerful sorcery. The larger the hubris of a bordering domain, the more likely a medusean tyranny is to set their ensorcelled armies upon it.


\begin{itemize}
\begin{multicols}{2}
\item \textbf{Skills}
\item \textbf{Diplomacy:} +2
\item \textbf{Espionage:} -1
\item \textbf{Lore:} +3
\item \textbf{Operations:} +0
\item \textbf{Defenses}
\item \textbf{Communications:} 10
\item \textbf{Resolve:} 12
\item \textbf{Resources:} 11
\end{multicols}

\end{itemize}

\subsection{Domain Title}
The medusean tyranny leader gains the following title and the noted additional features.

\paragraph{Optional Feature}
Veiled Sovereign|KaW

\subsection{Domain Powers}
Medusean tyranny officers gain access to the following domain powers.

\paragraph{Optional Feature}
Doom Curse|KaW

\paragraph{Optional Feature}
Unrelenting Gaze|KaW

\subsection{Domain Features}
Medusean tyranny officers gain access to the following domain features.

\paragraph{Optional Feature}
Allies to Stone|KaW

\paragraph{Optional Feature}
Cult of Secrets|KaW

\paragraph{Optional Feature}
Living Statuary (Special Unit)|KaW

\paragraph{Optional Feature}
Promises of Power|KaW

\section{Orc Clan}
Fearless and unwavering in defense of their homelands or in pursuit of plunder, the legions of an orc clan are not to be trifled with. Many folk call themselves superior warriors, with their heavy plate mail and firearms. But countless tales speak of foolish champions who underestimated the power of an orc clan—and the dear price they paid for that ignorance.

Most orc clans are composed of numerous smaller tribes and extended families, creating a bond that sees clan members defend their fellows with ruthless tactics and little patience for insult. Clan leadership is not hereditary, but is only rarely determined by age or strength. Instead, the mantle of chieftain is usually passed down to whoever can inspire the most loyalty in their fellow warriors.

Many of those who lose a campaign against an orc clan tell tales of the orcs' brutality, but this paints a picture that is far from accurate. Orc clans connect to a proud warrior culture, but their social infrastructure is no more savage than that of any other civilization. Although some clans are driven by ambition that sees them challenge and conquer other realms, most are dedicated to defending the orcs' expansive homelands, attacking only when provoked.


\begin{itemize}
\begin{multicols}{2}
\item \textbf{Skills}
\item \textbf{Diplomacy:} +1
\item \textbf{Espionage:} -1
\item \textbf{Lore:} +1
\item \textbf{Operations:} +2
\item \textbf{Defenses}
\item \textbf{Communications:} 10
\item \textbf{Resolve:} 13
\item \textbf{Resources:} 11
\end{multicols}

\end{itemize}

\subsection{Domain Title}
The orc clan leader gains the following title and the noted additional features.

\paragraph{Optional Feature}
Chieftain|KaW

\subsection{Domain Powers}
Orc clan officers gain access to the following domain powers.

\paragraph{Optional Feature}
Outlast|KaW

\paragraph{Optional Feature}
Protect the Clan|KaW

\subsection{Domain Features}
Orc clan officers gain access to the following domain features.

\paragraph{Optional Feature}
Proud Warriors|KaW

\paragraph{Optional Feature}
Take No Insult|KaW

\paragraph{Optional Feature}
Tales of Clan Heroes Past|KaW

\paragraph{Optional Feature}
The Unbreakable (Special Unit)|KaW

\section{Planar Invaders}
Beyond the life known to the folk of the world, countless other worlds exist, filled with riches, mysteries, and danger. Planar invaders travel between these dimensions, seeking and consuming rare resources, obliterating those who get in their way, and too often leaving behind a trail of suffering. They pay no heed to the mundane conceptions of war understood by worldly folk, with their alien mindset and refusal to follow rules of engagement. Rather, they exploit every advantage they have, with no care to whether or not they leave survivors behind.


\begin{itemize}
\begin{multicols}{2}
\item \textbf{Skills}
\item \textbf{Diplomacy:} -1
\item \textbf{Espionage:} +2
\item \textbf{Lore:} +0
\item \textbf{Operations:} +3
\item \textbf{Defenses}
\item \textbf{Communications:} 11
\item \textbf{Resolve:} 12
\item \textbf{Resources:} 10
\end{multicols}

\end{itemize}

\subsection{Domain Title}
The planar invaders' leader gains the following title and the noted additional features.

\paragraph{Optional Feature}
Overlord|KaW

\subsection{Domain Powers}
Planar invaders' officers gain access to the following domain powers.

\paragraph{Optional Feature}
Resistance is Futile|KaW

\paragraph{Optional Feature}
You'll Be of Use Yet|KaW

\subsection{Domain Features}
Planar invaders' officers gain access to the following domain features.

\paragraph{Optional Feature}
Interplanar Blip|KaW

\paragraph{Optional Feature}
Recall the Fleet (Special Unit)|KaW

\paragraph{Optional Feature}
Spatial Anchor|KaW

\paragraph{Optional Feature}
We Will Take What We Need from Your Mind...Directly|KaW

\section{Reptilian Band}
Reptilian folk vary widely in their cultures and outlooks, with each preoccupied with unique concerns and internal politics. But every so often, a charismatic leader steps forth to bind disparate factions of kobolds, lizardfolk, dragonborn, and other peoples of reptilian ancestry into a mighty band. When united, these brave peoples make an unbreakable bulwark against those who would intrude into or steal resources from their lands—or a furious swarm dedicated to leveling whole kingdoms and reshaping the world as it was in primordial times.

\begin{DndSidebar}[float=!b]{Frogfolk}
The illustrious history of the world's oldest fantasy roleplaying game is filled with various frog peoples. While such creatures are technically amphibians, the reptilian band domain works great to create a villainous or NPC realm using such creatures.



\end{DndSidebar}


\begin{itemize}
\begin{multicols}{2}
\item \textbf{Skills}
\item \textbf{Diplomacy:} -1
\item \textbf{Espionage:} +2
\item \textbf{Lore:} +3
\item \textbf{Operations:} +0
\item \textbf{Defenses}
\item \textbf{Communications:} 11
\item \textbf{Resolve:} 12
\item \textbf{Resources:} 10
\end{multicols}

\end{itemize}

\subsection{Domain Title}
The reptilian band leader gains the following title and the noted additional features.

\paragraph{Optional Feature}
Matriarch Rex|KaW

\subsection{Domain Powers}
Reptilian band officers gain access to the following domain powers.

\paragraph{Optional Feature}
Regenerate|KaW

\paragraph{Optional Feature}
Vengeful Shriek|KaW

\subsection{Domain Features}
Reptilian band officers gain access to the following domain features.

\paragraph{Optional Feature}
Booby Traps|KaW

\paragraph{Optional Feature}
Ensorcelled Scales|KaW

\paragraph{Optional Feature}
Prehistoric Wrath (Special Unit)|KaW

\paragraph{Optional Feature}
Primal Magic|KaW

\section{Undead Dominion}
Undeath is perhaps the oldest form of fell magic, its foul stain marring history again and again despite the efforts of those who pledge to wipe it out once and for all. Not all undead are evil by nature, but most leaders who seek to live forever are, as they manipulate the magic of life and death to eternally exercise their will over others. And soon enough, some of those undead leaders begin to ask themselves: "Why shouldn't my people also share in the gift of immortality...?"

When an undead dominion rises, it creates a blight on the living landscape, turning soil to ash and draining the energy from all living things. Its people are a labor force that never needs rest. A military that needs no pay or medical care. They cannot be broken. They cannot be bought. They do not need to eat, or sleep, or breathe. They—and the threat they bring to bear on neighboring lands—are perpetual.


\begin{itemize}
\begin{multicols}{2}
\item \textbf{Skills}
\item \textbf{Diplomacy:} -1
\item \textbf{Espionage:} +0
\item \textbf{Lore:} +1
\item \textbf{Operations:} +3
\item \textbf{Defenses}
\item \textbf{Communications:} 11
\item \textbf{Resolve:} 13
\item \textbf{Resources:} 10
\end{multicols}

\end{itemize}

\subsection{Domain Title}
The undead dominion leader gains the following title and the noted additional features.

\paragraph{Optional Feature}
Deathlord|KaW

\subsection{Domain Powers}
Undead dominion officers gain access to the following domain powers.

\paragraph{Optional Feature}
Knee Deep in the Dead|KaW

\paragraph{Optional Feature}
Rise Up|KaW

\subsection{Domain Features}
Undead dominion officers gain access to the following domain features.

\paragraph{Optional Feature}
Mass Animate Dead (Special Unit)|KaW

\paragraph{Optional Feature}
The Curse Spreads|KaW

\paragraph{Optional Feature}
The Dead Do Not Falter|KaW

\paragraph{Optional Feature}
Vile Resistance|KaW

\section{Undersea Colony}
Some ancestries whose folk dwell below the waves make valuable allies, but not all pelagic creatures are friendly toward the people of the surface world—and few monsters stoke fear so visceral as those that hide in the unfathomable depths of the sea. The home of such creatures is a quiet, crushing, everlasting night that holds more in common with the deep void between the stars than the lands the waves crash against. Undersea folk live among ancient cities that sank long ago, filled with prehistoric secrets best left forgotten.

Land-dwellers who come into conflict with an undersea colony often see its residents as detached and remorseless. Their strange languages often hide even stranger morals, such that slaying the crew of a ship to retrieve a stolen artifact of the deep might seem no different to them than peaceful diplomacy.


\begin{itemize}
\begin{multicols}{2}
\item \textbf{Skills}
\item \textbf{Diplomacy:} +0
\item \textbf{Espionage:} +0
\item \textbf{Lore:} +2
\item \textbf{Operations:} +1
\item \textbf{Defenses}
\item \textbf{Communications:} 13
\item \textbf{Resolve:} 11
\item \textbf{Resources:} 10
\end{multicols}

\end{itemize}

\subsection{Domain Title}
The undersea colony leader gains the following title and the noted additional features.

\paragraph{Optional Feature}
Tideweaver|KaW

\subsection{Domain Powers}
Undersea colony officers gain access to the following domain powers.

\paragraph{Optional Feature}
The Bends|KaW

\paragraph{Optional Feature}
Uncanny Gift|KaW

\subsection{Domain Features}
Undersea colony officers gain access to the following domain features.

\paragraph{Optional Feature}
Dark Figures Move and Twist|KaW

\paragraph{Optional Feature}
Disrupt Shipping|KaW

\paragraph{Optional Feature}
Eerie Insight|KaW

\paragraph{Optional Feature}
Salt Golems (Special Unit)|KaW

\section{World Below City-State }
Most people think of the World Below as a vast network of caverns stretching beneath the surface of the everyday world. But this is another realm entirely—a separate manifold in reality that is home to strange creatures all vying for control. Only the craftiest folk survive the host of subterranean horrors and the constant conflict between factions in the deep. Drow, duergar, deep gnomes, and more all build city-states strong enough to weather any setback or assault. The people of these city-states rarely venture beyond their well-defended territories—but when they do, it is often in search of the resources that will help guarantee their continued security. While most World Below city-states have no interest in conquering the folk of the surface, their highly trained forces can strike quickly when necessary—or when provoked by aggression on any side.


\begin{itemize}
\begin{multicols}{2}
\item \textbf{Skills}
\item \textbf{Diplomacy:} -1
\item \textbf{Espionage:} +2
\item \textbf{Lore:} +0
\item \textbf{Operations:} +2
\item \textbf{Defenses}
\item \textbf{Communications:} 10
\item \textbf{Resolve:} 11
\item \textbf{Resources:} 13
\end{multicols}

\end{itemize}

\subsection{Domain Title}
The World Below city-state leader gains the following title and the noted additional features.

\paragraph{Optional Feature}
Strike Leader|KaW

\subsection{Domain Powers}
World Below city-state officers gain access to the following domain powers.

\paragraph{Optional Feature}
Close-Quarters Magic|KaW

\paragraph{Optional Feature}
Crawling Toxin|KaW

\subsection{Domain Features}
World Below city-state officers gain access to the following domain features.

\paragraph{Optional Feature}
Blades in the Dark|KaW

\paragraph{Optional Feature}
Subterranean Bounty|KaW

\paragraph{Optional Feature}
Tunnel Up|KaW

\paragraph{Optional Feature}
Worm Knight (Special Unit)|KaW

\chapter{Warfare}
\DndDropCapLine{W}{}hen great heroes like Ajax, Hector, or Achilles participate in a war, the story is about them fighting other heroes. Enemy heroes. The Trojan War is just a backdrop for these individuals' epic stories. So if that's the story you want to tell—of great armies clashing in the background while heroes and villains duel each other in single combat—then you don't need rules for warfare. You can just play out normal encounters while a huge battle is described happening around the heroes and their enemies, like some particularly dramatic background music.


This system, by contrast, assumes that even in a fantasy world, armies are important. Heroes can fight other heroes, but it takes an army to capture and hold territory—and to defend a domain against a villain's army. Each army's tactics, the units deployed, and who won is important—as is how they won. So while the characters still need to stop the villain in a thrilling combat, the characters' army also needs to defeat the villain's army to fully secure victory.

\section{Using These Rules}
These rules don't worry overmuch about where the player characters are during a battle, or exactly what they're doing. Most are probably off adventuring or fighting the leaders of the villain's army. They're not on the battlefield trying to micromanage their own army. But that army is still an extension of the characters. Each unit gains benefits in battle based on a player character's class (see \textit{Martial Advantages} on page 111), but a character's features, traits, spells, and feats mostly help make them better at fighting monsters. Characters aren't designed to fight armies. So these rules let the characters focus on fighting monsters and villains while their army tries to take and hold the villain's territory, or to stop the villain's army from taking territory of its own.

Players and GMs who like these rules and want to use this system in their games should make sure that everyone understands wars are about armies. The players and the GM control their own side's units, but the characters, the monsters, and the NPCs should always be focused on undertaking the personal challenges and one-on-one combat they're best at.

With all this in mind, we think everyone will find a lot of value in these rules. As an overview, this section presents just a few of the ways players and GMs can use this warfare system in your games.

\begin{DndSidebar}[float=!b]{Combat and Battle}
To keep clear the differences between warfare and characters engaging in combat, this book uses "combat" to refer to characters and monsters skirmishing using the game's normal rules. "Battle" is used to refer to armies clashing using the warfare rules.



\end{DndSidebar}

\subsection{We're Adventuring!}
Just because the heroes go off on adventures doesn't mean the rest of the world stops short, waiting in suspended animation for them to return. Armies are a fantastic tool that can be used to create tension in a game while the heroes are off adventuring.

Giving the villain of the adventure an army to attack the heroes' domain, the domains of their allies, or even just the place they live is a great way to get the players engaged with warfare. A GM can pause and say "Meanwhile..." as they describe a scenario in which the heroes' domain is threatened and the players must use the heroes' army to defend it.

\subsection{Downtime Warfare}
One of the most canonical uses of warfare is to give the players more—and more interesting—things to do during downtime, or at any time between group adventures when the characters are free to be more self-directed. In the time between one adventure and another, the players have a chance to come up with their own goals and motivations. And having characters be officers in a domain with a standing army, even a small one, opens up all sorts of new opportunities for story and game play.

The players might decide to use the characters' army proactively. Send it somewhere to liberate a town or temple, or to clear out a forest infested with cultists. Send it to blockade a road and force a confrontation with a local warlord. Or they might use it reactively based on challenges the GM places in front of them, like a rebellion or an uprising. Between adventures, unlimited possibilities for drama and action arise, and warfare can easily factor into these.

\subsection{Epic Finale}
The most spectacular use of this system is to simulate a climactic confrontation between a heroic organization and a villainous realm. The easiest and most fun way to do this is to run and resolve the battle first, then the combat, even though both are actually happening simultaneously. This works much the same way that players in combat resolve their characters' actions in distinct turns, even though in reality the characters are all moving and attacking at the same time.

When this is done, mark down who wins the battle and on which turn. Then, on that turn of the combat, the winning officers each gain a Morale Surge they can use at any time (see the sidebar). The warfare rules have been designed to, on a whole, resolve faster than character combat. This is not wholly realistic, but it is dramatic—and it allows the game to continue to focus on the heroes and their actions. Their army will win or lose before the end of the combat, and the victors gain a sudden and dramatic bonus reflecting the rush of adrenaline.

Of course, this requires some willing suspension of disbelief, since the players will have already resolved the battle before starting the combat. So they know who's going to win—and on what turn— going into their own final fight. But even if this causes characters to change their behavior in what is technically a metagaming sense, it can easily be explained as the characters being the ones who trained the forces fighting outside—and as such, being instinctively aware that those forces are going to win (or lose!), even if that wasn't strictly obvious during the actual running of the battle.

A battle can, of course, take longer than a combat! This isn't a problem if the winner of the battle is the same domain as the winner of the combat. But even if that's not the case, when a domain's commanders lose a combat (whether they died, surrendered, were banished to another dimension, and so on) their troops automatically undertake the Retreat maneuver. (The rest of this chapter explains maneuvers and all the other rules for warfare.)

\begin{DndSidebar}[float=!b]{Morale Surge}
A character experiences a sudden rush of adrenaline upon hearing the herald's news of victory on the field of battle. As a bonus action, the character can move half their speed and take one additional action, on top of their regular action and movement. The character can use this benefit once, and loses it if they don't use it before taking a long rest.



\end{DndSidebar}

\subsection{Bringing the Siege}
When the GM decides that it's time for a battle to begin, each side can attempt to bring the siege. This section of the rules vaguely emulates the behavior of European military conflicts from the eleventh century to the fifteenth century. During these times, the majority of battles were sieges—one army attacking a fortification, rather than two armies facing each other in the field. The major question of the time was: Which side would mobilize first and lay siege to the other side's castle or keep?

The warfare rules simulate this historical setup by requiring both domains to make an opposed Operations test at the start of each battle. The winner decides whether they will attack the villain's stronghold, defend their own stronghold, or meet on a different field of battle. In the case of a tie (or if the GM determines that it's not appropriate for a particular battle to involve a stronghold), both armies meet at a location away from either stronghold.

\subsection{Seizing the Initiative}
The army that defends its domain's stronghold gains the benefits of the fortifications of that stronghold. The army that attacks has more flexibility in its tactics and is said to have seized the initiative. All units in that army's vanguard rank get a free activation at the start of the battle. While the attacking army is taking this free activation, the units of the defending army can use reactions only.

\subsection{Using the Grid}
One of the keys to keeping battles fun is making sure everyone has room to maneuver. Each side in the battle has a 4 × 5 grid—twenty spaces in total— on which to place their infantry and artillery units. Cavalry and aerial units do not use the grid.

If one side of the battlefield is completely filled with troops, that side is left with no room to maneuver and the game turns into a slog. As such, these rules work best with each side controlling twelve or fewer units.

\subsection{Making Tests}
Throughout these rules, you'll see references to units making tests. Just like with the rules for running domains, tests in warfare are analogous to your character making ability checks—you roll a d20, add a modifier, and compare the total to a fixed DC or some other value. If the result is equal to or greater than the value it's compared to, the test succeeds. Tests can likewise be made with advantage and disadvantage, just like ability checks.

When a unit attacks another unit, it first makes an Attack test opposed by the opposed unit's Defense. To successfully execute a special maneuver, a unit must succeed on a Command test with a DC based on the complexity of the maneuver. (Attack, Defense, and Command are three of the modifiers that are part of every unit's battle statistics, discussed in more detail below.)

\subsection{Starting Small}
The best way to test out these rules is to start with a small battle, wherein each player controls just one unit. There are then a commensurate number of units on the opposing side, controlled by the GM. A battle of this size—perhaps two opposing forces trying to take over a small town—gives everyone, including the GM, the chance to learn the system and get their feet wet.

It might seem like a fun idea to start with a huge battle! An epic confrontation with dozens of units on each side! But this system will quickly bog down under that kind of setup, and there are better ways to create this epic feel. These rules don't model individual soldiers, which means that players and GMs are free to redefine the scale of battle to suit their needs. If you want to run an epic battle, simply say that each unit represents a legion of five thousand soldiers rather than a squad of fifty or some other smaller number.

Alternatively, many of what are considered huge historical battles were actually a series of smaller related skirmishes. So if you want to replicate this kind of massive battle, you can use this system to play out specific skirmishes within a larger war. Using different battlefield setups each time, you'll play out the battle for the port, then the battle for the tower gates, and finally the battle for the city, with the results of each battle affecting the next.

\begin{DndSidebar}[float=!b]{Allied and Opposed}
In the warfare rules, "allied units" refers to all units that are part of an army or are friendly to that army. "Opposed units" refers to all units that are part of the army fielded by the other side in a battle, or which are friendly to that other side's army. The same usage applies when talking about opposed commanders. This allows the rules to read properly from the perspective both of players running their characters' army and the GM running a villain's army, by avoiding terms like "enemy" that might be misread as referring only to the units of the villain or GM's side.



If any part of the rules allows a unit to affect allied units, the unit can also affect itself (assuming it meets all other criteria for the specific effect).



\end{DndSidebar}

\subsection{Larger Battlefields}
Each player character involved in a warfare battle can command a number of units equal to their proficiency bonus, which means players can eventually control up to six units each! If your group has five or six players, that can mean battles with thirty units or more on a side!

Any time you play a battle with more than fifteen units on a side, it's recommended that the battlefield be expanded by adding ranks (rows) and files (columns) to the field, so that there are always openings and empty spaces. As a general rule of thumb, add one new column for every three units beyond fifteen. Once you've added two columns, add another row: a second vanguard row with the same rules for deployment as the original vanguard row.

\begin{DndTable}[header=Expanding the Battlefield]{XXX}

Units per Side & Add... & \\
15 or fewer & Nothing; use the 4 × 5 grid for each side & \\
16-18 & One column & \\
19-21 & Another column & \\
22-24 & One vanguard row & \\
25-27 & Another column & \\
28 or more & Another column & \\
\end{DndTable}

\subsection{Domain Armies}
Not all domains need an army. Depending on the campaign or the nature of the specific adventure you're playing, domains for the characters and the villains are still useful for the additional features they grant, even without these warfare rules. A heroic organization might act behind the scenes, or lend aid to other domains in the form of special units that can be mustered and sent to help allies. But if an enemy has an army marching on the characters' stronghold—or vice versa—creating an army is the best way to play out that scenario!

\subsection{Units}
Each army is made up of units. Each unit is a collection of soldiers (including monsters) with the same ancestry, training, gear, and type.

An army defends a domain and fights battles for the domain's officers (either the player characters or the villains controlled by the GM). While the characters confront an opposed domain's leaders in an epic combat, the armies loyal to both sides clash in an epic battle. The characters and the villains issue orders, but the soldiers who fight and die in response to those orders are the ones who take and hold the field—the ones who win or lose the war.

\subsection{Anatomy of a Unit}
The following is a typical unit, so let's check out its stats.

There are a lot of numbers and symbols at use here, but you'll quickly get used to them. And you'll find much more information on all these stats and how they're used in the sections that follow.

\subsubsection{Name}
Every unit has a name. For normal units, that name is just a combination of the unit's ancestry and type—for example, Human Cavalry, Elf Artillery, or Hobgoblin Infantry. Some units have more descriptive names, though, such as Human Shield Maniple or Bugbear Heavy Claw. Special units might follow either format. And unique units have quotation marks around their names, to indicate that there is only one of these special units anywhere in the world.

\subsubsection{Commander}
Every unit also has a commander. This is the creature that issues orders to the unit in battle—usually one of the officers in charge of the domain fielding the unit.

The rules don't worry too much about how exactly a human paladin locked in combat with an evil enemy cleric is able to make their intentions known to the unit they command out on the battlefield. The system is deliberately abstract. Runners and heralds are tasked with conveying such orders, and minor magics can facilitate communication between a commander and their troops. But beyond that, a unit knows what its commander wants because the commander has trained that unit. As well, each unit is assumed to have sergeants and subcommanders who know what to do if the commander is incapacitated, killed, or turned into a slimy toad.

\subsubsection{Battle Statistics}
 

\textbf{Attack (ATK)} is a measure of a unit's tactical effectiveness—its ability to successfully execute an attack order and engage an opposed unit. A successful Attack test can inflict casualties on an opposed unit, but is only half of the process of making an attack. Attack is opposed by Defense and improved by experience.

\textbf{Defense (DEF)} is a measure of a unit's ability to maneuver in such a way that it can avoid an opposed unit's attack, or minimize it so much that it fails to make a noticeable dent in the unit's casualties. Defense is improved by experience.

\textbf{Power (POW)} is a measure of a unit's physical prowess in battle. When the halberds strike, the arrows fall, or the lances pierce, how much force is behind them? Making a Power test is the second part of making an attack, and inflicts even more casualties on an opposed unit. Power tests made for other reasons do not deal damage unless the rules say so. Power is opposed by Toughness and improved by equipment.

\textbf{Toughness (TOU)} is a measure of a unit's physical hardiness—its ability to withstand successful attacks and keep fighting. Toughness is improved by equipment.

\textbf{Morale (MOR)} measures a unit's ability to maintain discipline in the face of overwhelming odds, powerful magic, exotic enemies, and death. Failing a Morale test inflicts casualties just as Attack tests and Power tests can, since a soldier who is terrified and flees is just as ineffective as a dead soldier. Morale is improved by experience.

\textbf{Command (COM)} is a unit's ability to correctly interpret complex orders and execute them successfully. Command is improved by experience.

\textbf{Number of attacks (the sword icon)} represents the number of times a unit can organize itself to attack other units in one round of battle. Most units have only one attack. Number of attacks is improved by experience.

\textbf{Damage (DMG)} is the number of casualties inflicted on a successful Power test made after a successful Attack test. If some part of the rules grants a bonus to damage, that bonus applies only to the Power test made as part of an attack. Damage is improved by equipment.

\textbf{Size} is an abstract measure of a unit's ability to maintain morale and unit cohesion, and to carry out orders. Size determines a unit's \textbf{casualty die} (see below), and so determines how much damage a unit can take. The bigger a unit's size, the more effective it is at staying in the fight.

\textbf{Movement} is another of a unit's battle statistics, but does not appear on the card because all units have a movement of 1.

\subsubsection{Casualties}
During warfare, the number of casualties a unit can suffer but still keep on fighting is represented by the \textbf{casualty die}. A casualty die is based on a unit's size, and is placed on the unit card during battle with the number of casualties the unit has remaining facing up. For example, a size 6 unit would begin with a d6 casualty die, with the 6 facing up.

If a unit is ever reduced to half its maximum casualties or fewer, it is \textbf{diminished}. The first time a unit is diminished in a battle, it must succeed on a DC 13 Morale test or immediately suffer 1 casualty.

If a unit's casualty die ever reaches 0, the unit is \textbf{broken} and is removed from the battle. A broken unit might be able to be rallied to return to the battle, or it might be disbanded and permanently lost.

These rules don't worry about the exact number of soldiers or monsters in a unit. Likewise, casualties usually mean soldiers killed by enemy action, but not always. A soldier who is disoriented or too terrified to act, or who cannot see or hear or understand their commanders' orders, is just as ineffective as a dead soldier. For this reason, even though a unit might be removed from the battlefield when broken, there is a chance that it can be rallied later by gathering together scattered soldiers, bolstering their resolve, and organizing them for battle again.

\subsubsection{Type}
A unit's type is represented by a weapon symbol. Type determines what a unit can do in battle, who it can attack, and who it can be attacked by. There are four types of basic unit: infantry, artillery (including siege weapons), cavalry, and aerial. For the purpose of improving experience and equipment (see below), cavalry and aerial units use the same tables.


\begin{DndTable}{XcXcX}

Type & Symbol & Type & Symbol & \\
Aerial &  & Cavalry &  & \\
Artillery &  & Infantry &  & \\
Artillery (Siege) &  & Infantry (Levy) &  & \\
\end{DndTable}

\subsubsection{Experience}
Experience is a combination of how much training a unit has and how much fighting it's seen. More experienced units gain bonuses to Attack, Defense, Morale, and Command, as well as additional attacks, as indicated on the Unit Experience Bonuses table. But different types of units learn different lessons in battle.

Units gain experience by surviving battles, according to how experienced they already are (see \textit{Improving Units} on page 104). Each section of the table starts with the lowest level of experience at the top, with each subsequent level below indicating the next improvement.

The bonuses on these tables are \textbf{not cumulative}. If a veteran infantry unit survives a battle and improves to become elite, its Attack modifier goes up by 1, not by 2.

The number of stars in a unit card's upper left corner indicate the unit's level of experience. Levies and regular units have no stars, veteran units have one star, elite units have two stars, and super-elite units have three stars.

\begin{DndSidebar}[float=!b]{Levies}
Levies are infantry troops who have zero training and no experience, typically laborers and townsfolk. They know their likely fate in battle, but they believe enough in a domain's cause to fight and die for it—or they fear the price of not fighting even more.



Levies cannot have their gear or experience improved. They automatically disband after every battle and must be mustered again for the next battle.



\end{DndSidebar}

\subsubsection{Equipment}
Equipment describes a unit's arms and armor. Heavier units have better weapons and armor, granting them bonuses to Power, Toughness, and Damage (except for artillery units), but they might not be as flexible in battle as a lighter, more mobile unit. As with experience, different types of units gain different benefits from improving their equipment.

Units improve their equipment by having their domain pay for those improvements (see \textit{Improving Units} on page 104). Each section of the Unit Equipment Bonuses table starts with the baseline level of equipment at the top, with each subsequent level below indicating the next improvement.

\begin{DndTable}[header=Unit Experience Bonuses]{XcccccX}

Infantry Experience & Additional Attacks & Attack & Defense & Morale & Command & \\
Levies & +0 & -1 & -2 & -2 & -1 & \\
Regular & +0 & +0 & +0 & +0 & +1 & \\
Veteran & +0 & +1 & +2 & +2 & +2 & \\
Elite & +1 & +2 & +4 & +4 & +2 & \\
Super-elite & +1 & +3 & +6 & +6 & +3 & \\
\end{DndTable}


\begin{DndTable}{XcccccX}

Cavalry and Aerial Experience & Additional Attacks & Attack & Defense & Morale & Command & \\
Regular & +0 & +0 & +0 & +0 & +0 & \\
Veteran & +0 & +1 & +1 & +1 & +2 & \\
Elite & +1 & +2 & +2 & +2 & +4 & \\
Super-elite & +1 & +3 & +3 & +3 & +6 & \\
\end{DndTable}


\begin{DndTable}{XcccccX}

Artillery Experience & Additional Attacks & Attack & Defense & Morale & Command & \\
Regular & +0 & +0 & +0 & +0 & +0 & \\
Veteran & +1 & +2 & +1 & +1 & +1 & \\
Elite & +1 & +4 & +2 & +2 & +2 & \\
Super-elite & +1 & +6 & +3 & +3 & +3 & \\
\end{DndTable}

\begin{DndTable}[header=Unit Equipment Bonuses]{XcccX}

Infantry Equipment & Power & Toughness & Damage & \\
Light & +0 & +0 & +0 & \\
Medium & +2 & +2 & +0 & \\
Heavy & +4 & +4 & +0 & \\
Super-heavy & +6 & +6 & +1 & \\
\end{DndTable}


\begin{DndTable}{XcccX}

Cavalry & Aerial Equipment & Power & Toughness & Damage & \\
Light & +0 & +0 & +0 & \\
Medium & +1 & +1 & +0 & \\
Heavy & +2 & +2 & +0 & \\
Super-heavy & +3 & +3 & +1 & \\
\end{DndTable}


\begin{DndTable}{XccX}

Artillery Equipment & Power & Toughness & \\
Light & +0 & +0 & \\
Medium & +1 & +1 & \\
Heavy & +2 & +2 & \\
Super-heavy & +3 & +3 & \\
\end{DndTable}

\subsubsection{Ancestry}
Ancestry affects all of a unit's stats, and in many ways is the defining attribute of a unit. Ancestry also determines what traits a unit has.


\begin{DndTable}{XcXcX}

Ancestry & Symbol & Ancestry & Symbol & \\
Dragonborn &  & Human &  & \\
Dwarf &  & Kobold &  & \\
Elf &  & Lizardfolk &  & \\
Fiend &  & Monstrous &  & \\
Giant &  & Orc &  & \\
Gnoll &  & Special &  & \\
Gnome &  & Undead &  & \\
Goblinoid &  &  &  & \\
\end{DndTable}

The bonuses a unit gains because of its ancestry have very little to do with any of the stats that creatures of those ancestries have in the core rules. The warfare rules are concerned with how well creatures organize, take orders, and maintain morale and unit cohesion, none of which are necessarily related to how well creatures fight in single combat.

\subsubsection{Tier}
A unit's tier, represented by a Roman numeral on its card, is a measure of the unit's overall power or nastiness. Tier I units are the least powerful (and thus the easiest to put onto the battlefield), while Tier V units are the most powerful.

All units after Tier I have unit dependencies that determine how many units of a specific tier are needed to field a unit of a higher tier. For example, you can field an unlimited number of Tier I units, but you must have more Tier I units than Tier II units.

\subsubsection{Traits}
Each unit also has a collection of traits—maneuvers or special features that it can employ in battle. The names of each unit's traits are listed on its unit card, keying to the descriptions in the \textit{Unit Traits} section on page 126 in this chapter. This keeps the unit cards small and useful in battle, and means players don't have to constantly pick each card up to read what its traits can do in the middle of a battle.

\subsection{Building an Army}
As soon as a group of characters founds an organization, they muster four Tier I units of the players' choice from any ancestry the GM agrees the organization has access to. These are leaderless soldiers—probably regular light troops from other failed domains, looking for someone to serve.

As an option, any time a domain musters a Tier I unit, it can muster two levies instead. This can boost the size of a starting army by making it some combination of levies and regular troops. However, levies automatically disband at the end of the next battle they fight in.

Six ancestries are presented in this book—humans, elves, dwarves, orcs, goblinoids, and undead. Each ancestry features nine units covering different types and tiers. Whenever an organization musters new units, they must be from an ancestry the GM agrees the organization has access to, chosen from the nine units noted and following the rules for unit command and unit dependencies (detailed below).

(If you use \textit{Strongholds \& Followers} in your game, you can optionally use the rules that allow characters to attract a variety of followers, including new units, when they build or take over a stronghold.

Units mustered by a heroic organization belong to that organization. Each unit needs a commander (see \textbf{Unit Command and Dependencies} below) who must be an officer in the organization. These units gain experience through battle (unless they're levies), and the organization can spend gold to improve any unit's equipment.

Gaining additional units can be accomplished in three ways:


\begin{itemize}
\item An officer of the organization can use a domain action during intrigue to muster one unit from any ancestry the GM agrees the organization has access to (see the \textbf{Operations} section on page 23 in the \textbf{Domains \& Intrigue} chapter). Units mustered in this way belong to the organization and persist from battle to battle until they are disbanded. They must conform to the rules for unit dependencies.
\item An officer of the organization can muster the organization's special unit by succeeding on an Operations test during intrigue (see the \textbf{Operations} section on page 23, and the details of special units in each organization's write-up in the \textbf{Domains \& Intrigue} chapter). This unit disbands after the next battle it is used in, and must be mustered again before it can be used again. This unit ignores unit dependencies.
\item An officer of the organization can recruit new units from NPC realms by succeeding on a Diplomacy test during intrigue (see the \textbf{Diplomacy} section on page 20 in the \textbf{Domains \& Intrigue} chapter), gaining the use of that realm's special unit. Units gained by diplomacy disband after the next battle they are used in, and return to their home domain. Additional diplomacy is necessary if they are to be fielded again. These units ignore unit dependencies.
\end{itemize}

All units, regardless of how they are raised, need a commander. Units that conform to the rules for unit dependencies can be used to unlock new units based on their tier, but units that ignore unit dependencies cannot. For example, a Tier III unit mustered as an organization's special unit or recruited through diplomacy does not count toward being able to muster a Tier IV unit.

\subsection{Unit Ancestry}
The GM determines which ancestries an organization has access to when raising units. It's reasonable to assume a dwarf officer would know other dwarves, and that the officer's heroism has become known in their homeland. So when this heroic dwarf puts out the call, a unit of dwarves arrives! That said, a dwarf officer might be an exile, or might have been raised in a human city and have no real connection to any ancestral home. If so, the GM might decide that a different ancestry is available instead.

If the characters' organization is on good terms with a nearby NPC elf domain—perhaps because of interactions during a previous adventure—it's reasonable to assume that the organization can muster elven units. But the only way to muster the Court Jesters—a special aerial unit of veteran sprites with medium equipment belonging to a fey court domain—is through Diplomacy.

\subsection{Unit Types}
There are four types of units—infantry, artillery, cavalry, and aerial. Each has different base stats, and because of the rules governing which units can attack other units (see \textbf{Attacking}) and the rules for deployment, each has a specific role to play in warfare.

\subsubsection{Infantry Units}
Infantry are the meat-and-potatoes troops for any domain. They are not as flashy as artillery or cavalry, but they are harder to kill, having higher Toughness than other units of the same tier and ancestry. Infantry are often used to protect artillery units.

Infantry have very few legal targets in a battle. They can attack only adjacent units (in front, behind, to the left, or to the right of the infantry unit), and can't attack cavalry or aerial units at all. But infantry also have access to more maneuvers than other units. Only infantry units can use the Follow Up maneuver, which gives them free movement after an adjacent opposed unit breaks or moves away from them. This helps infantry quickly move into position to get closer to the opposed side's center rank—where all the squishy archers are usually hiding!

Infantry can also make use of the Set for Charge maneuver, which gives them a chance of inflicting casualties on any cavalry or aerial unit that attacks them. They cannot attack those more mobile units back, but they are not completely defenseless against them.

Infantry also include a domain's levies—the untrained troops who have either volunteered to support a domain or been pressed into service. Levies are cheaper than other Tier I units, with officers able to muster two units of levies instead of any other Tier I unit. But they have poor stats compared to other infantry units, and those stats cannot be improved. Levies disband after every battle.

The best use of levies is to defend a domain's rear rank against enemy cavalry. They often won't last long, but they'll force an opponent to waste actions wearing them down.

\subsubsection{Artillery Units}
Artillery units include both archers and siege engines. Mostly archers.

Artillery units can attack any other unit on the field, which means they're a domain's best defense against enemy aerial units. Those are thankfully rare, but archers' ability to pick targets from across the battlefield makes them incredibly good at two things—forcing Morale tests by diminishing opposed units, and breaking opposed units that have only 1 casualty remaining, regardless of where those units are.

Additionally, archers and cavalry can be coordinated to pick out a single enemy cavalry unit and wear it down, since archers and cavalry can attack any cavalry unit. Many battles begin with both sides using archers and cavalry in this way, trying to eliminate opposing cavalry before they can go on the offensive.

The downside, of course, is that archers are incredibly squishy. They have low Defense and Toughness compared to other troops of the same tier and ancestry. For this reason, they must be deployed in the center rank of the battlefield, and a large part of the strategy of a domain's officers typically involves keeping their archers alive.

Siege Engines are usually Tier II units, so they're harder to muster. They can attack any other units and inflict a lot of casualties—but most require an entire round between attacks doing nothing (loading the catapult or trebuchet, winding the ballista, and so forth). Combined with the fact that they're the only unit that can damage fortifications (which they automatically hit!), siege engines are typically only mustered when a domain is fighting an opposing force with fortifications, and are then used to batter down those walls and towers.

\subsubsection{Cavalry Units}
Cavalry units move fast and hit hard. They are so mobile that they belong to no specific rank, ranging across the battlefield to attack other cavalry, and pick off those infantry and artillery units foolish enough to leave themselves exposed.

Cavalry units have greater Power than other units of the same tier and ancestry, and they're one of the few Tier I units that deal 2 damage on a successful Power test. This means they can cause most Tier I units to become diminished with a single successful attack—a size 6 unit taking 1 casualty from a successful Attack test, and then 2 more casualties from a successful Power test!

Many battles begin with both sides trying to neutralize the other's cavalry. Units with the Archers trait are particularly good at this.

\subsubsection{Aerial Units}
Flying units are rare, with most of them in Tier III and having stats comparable to an average Tier II unit. This is because aerial units can attack any other unit (in the same manner as artillery units) but very few units can return fire against them! Only artillery and other aerial units can attack an aerial unit.

Mustering aerial units gives a domain a major advantage over an opposed domain that has no aerial units with which to defend itself. But while they have clear advantages in battle, aerial units are squishy. They have higher Attack modifiers than other units of the same tier and ancestry, but their need to be lightly armed and armored compared to other units gives them lower Defense and Toughness.

Properly used, aerial units can make a big difference in a battle, but they do not represent  instant victory. Rather, much of their usefulness is psychological. An army without aerial units facing an army with aerial units often feels as though it's on the defensive, even if it has archers enough to handle the situation.

\subsection{Unit Command and Dependencies}
A couple of restrictions affect how many units a domain can field, and which units it can field. This stops a battle from growing too large to manage, and simulates the fact that any army will have a very small number of elite units supported by a broad base of regular units. It also prevents domains from fielding an army made up of only one or two extremely nasty units that can't be defeated by a regular army.

\subsubsection{Unit Command}
An officer can command a number of units equal to their proficiency bonus. Every unit a domain fields—both regular units and special units, no matter how those units are gained—must have a commander.

Normally, only a domain officer (a player character member of a heroic organization, or the NPC leader and lieutenants for an enemy realm or some other opposed domain) can command units. But the GM might decide to let NPCs allied with a heroic organization command units as well. This is useful in scenarios that demand larger armies but have a small number of players.

\subsubsection{Unit Dependencies}
Each unit's tier represents not only its power level relative to other units. It represents the relationship wherein higher-tier units require many lower-tier units to support them, as follows:


\begin{itemize}
\item There is no limit to the number of Tier I units (including levies) a domain can field.
\item A domain can field a number of Tier II units one less than its current number of Tier I units.
\item A domain can field a number of Tier III units one less than its current number of Tier II units.
\item A domain can field a number of Tier IV units one less than its current number of Tier III units.
\item A domain can field a number of Tier V units one less than its current number of Tier IV units.
\end{itemize}

Remember that whenever a domain raises a regular Tier I unit, it can raise two units of levies instead. Because these levies are also Tier I, they are a quick way of increasing unit numbers to unlock more powerful units—at the cost of fielding much squishier units at the base of an army!

An example large army of fifteen units lead by a Tier V unit would look like this:

A domain's special unit, as well as units raised through diplomacy, ignore unit dependencies. These units can take to the battlefield regardless of what other units are in play for a domain.

\subsection{Improving Units}
For each of the units in play for a domain, experience and equipment can both be improved. Units gain experience through battle, and better equipment with money.

\subsubsection{Improving Experience}
The only way to improve a unit's experience is through battle, as follows:


\begin{itemize}
\item Regular units become veteran units if they survive their first battle. Easy!
\item Veteran units become elite units if they survive three more battles. Much harder.
\item Elite units become super-elite units after surviving four more battles (eight total).
\end{itemize}

Levies cannot be improved, as the untrained troops who make up those units are considered to go back to their lives as soon as their task in battle is done.

\subsubsection{Improving Equipment}
A domain can spend money to upgrade a unit's equipment from light to medium, medium to heavy, and so forth, requiring only a domain action. (Modifying Units on page 23 in the Domains \& Intrigue chapter has all this information.) A domain can upgrade as many units as desired with a single domain action, and those units can be upgraded to any level the domain can afford, from light to super-heavy. A domain can upgrade any units it controls.


\begin{DndTable}{XcX}

Upgrade from... & Cost\\
Light to medium & 500 gp × the tier of the unit & \\
Medium to heavy & 1,000 gp × the tier of the unit & \\
Heavy to super-heavy & 2,000 gp × the tier of the unit & \\
\end{DndTable}

\subsection{Voluntarily Disbanding Units}
Players and GMs might end up with an army fit to a purpose that decisively wins a climactic battle! But now a domain needs a different army for the next threat. To create the opportunity to muster new units, a domain's officers can voluntarily disband any units the domain controls during an extended rest. (The GM determines what an extended rest means in the campaign, but one week of downtime spent managing a domain is a good rule of thumb.) Units must be disbanded carefully, though, or the relationships of unit dependencies might be broken.

\subsection{Warfare Feats}
In campaigns that use the optional rule for feats, player characters can take any of the following feats to improve their units' stats in warfare.

Tactical Leader

Agile Leader

Lead from the Front

Unstoppable Leader

Persuasive Leader

Strategic Leader

\subsection{The Battlefield}
The battlefield is a grid five spaces wide, divided into two sections for the two opposing forces in a battle. Each section is four rows deep, with each row called a rank.

The topmost row in each section is the \textbf{vanguard}, which faces the opposed army. Behind the vanguard is an army's \textbf{reserve} rank, followed by the \textbf{center} rank, followed by the \textbf{rear}.

In addition to these four ranks, each side also has a \textbf{front}, which comprises the vanguard plus all enemy ranks. If a unit controlled by a domain moves from that domain's vanguard into the opposing side's vanguard, or even all the way into the opposing side's rear rank, they are still in the controlling domain's front.

With the exception of the reserve rank (see \textbf{Stacking} below), a space can only hold one unit.

\subsection{Setup}
At the start of a battle, all officers commanding units roll initiative. After initiative is determined for a battle, each character and NPC acting as a commander deploys their units in reverse initiative order. In other words, whoever acts last in the combat gets to place all their units first in the battle. If you are playing a battle alongside combat, use the same initiative order for both.

\subsection{Casualty Dice}
During deployment, each unit commander places a die on every unit they control. This is the casualty die representing the number of casualties the unit can suffer before it breaks. The number of casualties the unit has remaining should be facing up. So a size 6 unit would begin with a d6 casualty die, with the 6 facing up.

Some domain features can increase the size of a unit's casualty die. \textbf{Increasing a casualty die} means upgrading from a d4 to a d6, a d6 to a d8, and so forth. Decreasing a casualty die means downgrading in the same way.

As a unit suffers casualties, its casualty die is \textbf{decremented}. This means the number on the die is decreased by 1, so that a 6 becomes a 5, a 5 becomes a 4, and so forth. Under certain circumstances, the die is \textbf{incremented}—increased by 1—to indicate that a unit's battle readiness is being restored. A unit that has its casualty die decremented from 1, indicating that it has lost its last casualty, is \textbf{broken} and removed from the field.

Unless otherwise noted, no unit can gain casualties beyond the maximum value on its casualty die.

Once all units are placed, battle can begin. During battle, units are activated by the characters and NPCs controlling them in initiative order.

\subsection{Collapsing Ranks}
At the start of any round (including the first), if a rank contains no units or fortifications (see below), that rank collapses and is no longer part of the battlefield. The ranks in front of and behind the collapsed rank are now next to each other. Ranks that remain on the field keep their name and rules. For instance, if a rear rank collapses, the center rank next to it remains the center rank.

\subsection{Deployment}
When deploying units, artillery must go in the center rank. Infantry can be deployed in the vanguard, center, reserve, or rear ranks. Cavalry and aerial units are not placed in any rank, but can move freely across or above the battlefield.

\subsection{Stacking Units}
Units in the reserve rank can be stacked when deployed. When units are stacked, only the top unit can be targeted. Units under a stack are treated as though they were not on the battlefield for the purpose of being subject to attacks or effects.

Only the top unit of a group of stacked units can be activated. If that unit leaves the stack, the unit beneath it is now on top of the stack and can be targeted. Because a unit can't move into a space occupied by other units, units can only leave a stack during a battle, they can't enter one.

\subsection{Activation}
On each commander's turn (whether that commander is a player character or an NPC controlled by the GM), they activate all the units they control in any order. Each unit must finish its activation before another unit can be activated.

Each unit can take actions and move, in either order. The actions a unit can undertake during its activation include the following:


\begin{itemize}
\item Attack another unit
\item Use a magic item
\item Attempt a maneuver
\item Move 1 additional space (march)
\item Hold (do nothing)
\end{itemize}

Unless a unit holds and does nothing, it can move and take an action, take an action and move, or move and move (called a march).

\subsection{Attacking}
When a unit attacks, its choice of targets is determined by its unit type. (These parameters are sometimes referred to as the "order of battle.") The units of both sides are able to attack as follows:


\begin{itemize}
\item \textbf{Infantry units} can attack any adjacent unit— one in front, behind, to the left, or to the right. They cannot attack diagonally.
\item \textbf{Artillery and aerial units} can attack any unit.
\item \textbf{Cavalry units} can attack other cavalry units, and any infantry or artillery units that are exposed (see \textbf{Unit Conditions} on page 107).
\item \textbf{Siege engines} can attack any unit, and can attack fortifications.
\end{itemize}

Some units, typically infantry in the rear rank, might have no opposed units that are legal targets. (Units can attack allied units if compelled to do so by some effect, or just for fun.)

With a legal target selected, an attacking unit makes an Attack test against the target's Defense. If the Attack test succeeds, the attacking unit inflicts 1 casualty on its target and then makes a Power test against the target's Toughness. If the Power test succeeds, the attacking unit inflicts additional casualties equal to the unit's noted damage.

The attacks of artillery units are ranged attacks. The attacks of infantry, cavalry, and aerial units are melee attacks.

\subsection{Movement}
When a unit moves, it must move into an empty adjacent space—a space in front, behind, to the left, or to the right of its current space. Units cannot move diagonally. Units that have movement greater than 1 can move into successive empty spaces. A unit can move through other units only if it has a special feature that says so.

Though the front ranks of both sides in a battle are usually separated a little to make it easier to tell one side from another, units in one front can move into the opposing side's front if there is an empty space there. Units can then move normally through the opposing side's ranks as long as there are empty adjacent spaces for them to move into.

If a unit is in the rear rank of its army and moves backward, it is permanently removed from the battle.

\subsection{Bonus Movement}
A unit might gain "+1 to movement," which means the unit moves 1 extra space both when it moves and if it uses its action to move. (Which is to say, if it marches, it moves 2 extra spaces.)

A unit might suffer a penalty to movement (indicated by "−1 to movement"). If a unit with the normal speed of 1 and no bonuses to movement suffers a −1 movement penalty, it cannot move or march.

\begin{DndSidebar}[float=!b]{Unit Conditions}
Just like the conditions that affect characters, a number of unit conditions summarize the state of an army's units during the battle.



\textbf{Broken.} A unit that breaks becomes broken. It has lost its last casualty and is removed from the battle. Broken units can be reformed, usually by rallying.



\textbf{Disbanded.} A disbanded unit is removed from the game and cannot be reformed by normal means.



\textbf{Disorganized.} A disorganized unit does nothing on its next activation while it attempts to regain unit cohesion.



\textbf{Disoriented.} A disoriented unit can either take an action or move, but not both. Unless otherwise stated, this unit condition lasts until the end of the unit's next activation.



\textbf{Exposed.} A unit is exposed if there are no units between it and the leftmost or rightmost battlefield edge, or if there are no units in any rank to the rear of the unit. Units in the center and reserve of an army's ranks cannot be exposed as long as that army has its own units in both its rear rank and anywhere in its front. (See \textbf{Exposed} below for more information.)



\textbf{Hidden.} When a unit is hidden, other units have disadvantage on Attack tests against it.



\textbf{Misled.} A unit that is misled cannot attack, and spends its next activation moving randomly into an available space. Cavalry and aerial units cannot be misled.



\textbf{Weakened.} A unit that is weakened has disadvantage on Attack tests and Power tests.



\end{DndSidebar}

\subsection{Actions, Reactions, and Bonus Actions}
Each unit has the same types of actions in battle as a character does in combat—one action, one bonus action, and one reaction.

Each unit gets only one action per activation. If a unit gets multiple attacks, it makes those attacks as one action. If it has access to multiple actions, perhaps as a result of maneuvers or unit traits, it must choose which action to use on its activation.

A unit can use a bonus action only if some trait or other feature grants it a bonus action. A bonus action is used on the unit's activation, either before or after it attacks or moves. If a unit gets multiple bonus actions from traits or features, it must choose which one to use on its activation.

If a unit has a reaction, it can use that reaction during its own activation or on another unit's activation. Once a unit uses a reaction, it cannot use another one until its next activation. If the reaction interrupts another unit's activation, that unit can continue its activation after the reaction resolves.

\subsection{Resolving Multiple Reactions}
When resolving multiple reactions between units, the defending unit gets the first chance to resolve a reaction. Then the attacking unit can use its reaction, followed by the defending unit's allies, then the attacking unit's allies.

\subsection{Exposed}
Cavalry can attack exposed units. A unit is exposed if there are no units between it and the leftmost or rightmost battlefield edge, or if there are no units in any rank to the rear of the unit. By definition, all units in the rear ranks are always exposed. Units in the center and reserve of an army's ranks cannot be exposed as long as that army has its own units in both its rear rank and anywhere in its front. As soon as either of those conditions is no longer true, units in the center or reserve can become exposed. See the diagram below for more information.

\subsection{"Forward" and "Back"}
If the rules ever refer to a unit moving forward or back, those directions are relative to the unit's rear rank. Forward is moving away from the rear and back is moving toward the rear.

\subsection{"Opposite"}
Some rules refer to the unit opposite a target. This means the unit on the other side of the target from the attacker.

\subsection{Morale Tests}
When a unit is subject to magic or specific events, the rules sometimes call for that unit to succeed on a Morale test. This includes such events as a unit being diminished, when it must succeed on a DC 13 Morale test or suffer 1 casualty. Likewise, whenever a unit is affected by an opposed unit's battle magic, in addition to resolving the other effects of the magic, that unit must succeed on a DC 13 Morale test or suffer 1 casualty.

Unless otherwise stated, the DC for any Morale test is 13.

\subsection{Tokens}
Some unit traits and magic items can put tokens on units, with the trait or item description noting what effect the tokens create. When a unit with tokens on it activates, it immediately suffers the effect of the tokens. If a token inflicts damage, the unit takes damage for each token on it. At the end of the unit's activation, make a DC 12 Power test for each type (acid, fire, etc.) of token on it. On a success, remove one token of that type.

For example, a unit targeted by another unit with the Flaming Weapons trait and in a space containing an acid pool created by a \textit{wand of acid pool} might have two fire tokens and one acid token on it at the beginning of its activation. It suffers 3 damage (1 for each fire token and 1 for the acid token), then removes one fire token and one acid token.

\subsection{Diminished, Broken, and Disbanded}
A unit is \textbf{diminished} when its current casualties are half or less than its starting casualties. The first time a unit becomes diminished, it must succeed on a DC 13 Morale test or suffer another casualty. Each unit does this only once per battle.

A unit is \textbf{broken} and removed from the battle after it suffers its last casualty. This does not mean that every soldier in the unit is dead, but those who survive are left confused, panicked, squabbling—or might simply have fled. Broken units can be reformed with the Rally maneuver and certain martial advantages.

If a unit breaks and later fails a Morale test to rally, it is \textbf{disbanded}. (Certain martial advantages and other effects can also cause a unit to disband.) A disbanded unit is permanently destroyed. It is removed from the battle and cannot be rallied or reformed by normal means.

\subsection{Defeat}
At the end of each round of battle, after all officers on both sides have activated all their units, each side determines the current \textbf{point value} of their army. This total is based on the tier of each unit in the army, with Tier I units (including levies) worth 1 point, Tier II units worth 2, Tier III units worth 3, and so on.

At the end of any round, if one side has \textbf{more than twice the point value} of the other, that side has won the battle. Each commander on the losing side must immediately attempt the Retreat maneuver for each of their units (see \textbf{General Maneuvers} below).

\subsection{After the Battle}
Immediately after the battle, all units that were broken but not disbanded can execute the Rally maneuver. (This excludes units that have already been rallied during the battle, as a unit can rally only once per battle.)

Units that succeed at the Retreat or Rally maneuvers gain casualties at a rate of 1 per week, incrementing the casualty die until they have their full casualties once again.

Units that survive the battle, including those that rallied at the conclusion of the battle, can increase their experience, as discussed in \textbf{Improving Units} (page 104).

Levies immediately disband after the battle, as they are always temporary units. New levies must be mustered from scratch for the next battle. Likewise, special units go their own way, and units mustered through alliances or diplomacy return to their home domains.

\subsection{Ongoing Effects}
Any ongoing effects, tokens, or unit conditions automatically end for all units once the battle is over. It is assumed that once battle is no longer a priority, surviving troops and their officers can quickly resolve any problems a unit is having.

\section{General Maneuvers}
Maneuvers represent complex operations that a unit can attempt during its activation. Some maneuvers are made as reactions in response to specific events or conditions, as described in the maneuver.

Some maneuvers require a successful test, with the battle statistic and DC noted in the maneuver. On a failed test, the unit cannot perform the intended maneuver. However, it can still move or attack during its activation (though it can't do both).

Some maneuvers can be attempted only by certain types of units. Other maneuvers can be attempted by units of any type. Certain maneuvers (Follow Up, Rally, Retreat, Set for Charge, and Withdraw) can be undertaken by any appropriate units. Other maneuvers can be undertaken by appropriate units only if a martial advantage allows units to be trained to use the maneuver (indicated by the prerequisites).

This section collects maneuvers that are either generally available, or which can be assigned by martial advantages. Additional maneuvers are available as traits assigned to specific units, and are found in \textit{Unit Traits} (page 126).

\subsection{Feint}
\textit{Prerequisite: Rogue martial advantage}

As an action, choose an adjacent opposed unit and an empty space adjacent to that unit. This unit moves into the adjacent space and makes an opposed Command test against the target unit. On a success, the opposed unit moves into the space this unit just vacated.

\subsection{Find Cover}
\textit{Prerequisite: Ranger martial advantage}

As a reaction to moving, this unit can make a DC 13 Command test. On a success, any opposed unit attacking this unit has disadvantage on Attack tests until its next activation.

\subsection{Follow Up}
\textit{Prerequisite: Infantry}

As a reaction to an opposed unit adjacent to this unit breaking or moving out of its space, this unit makes a DC 8 Command test. On a success, this unit moves into the opposed unit's former space.

\subsection{Mobility Trap}
\textit{Prerequisite: Barbarian or monk martial advanage}

As a reaction when an opposed cavalry or aerial unit successfully attacks this unit, this unit makes a DC 13 Command test. On a success, the attacking unit's type becomes infantry and this unit places the attacking unit in any empty space closest to this unit. The attacking unit's type reverts back at the beginning of its next activation, and it is removed from its temporary space.

\subsection{Rage Charge}
\textit{Prerequisite: Barbarian martial advantage}

As an action, this unit makes a DC 13 Command test. On a success, the unit moves up to 5 spaces. If it ends this movement adjacent to an opposed unit, it can attack that unit.

\subsection{Rally}
\textit{Prerequisite: Any broken unit}

After the battle, this broken unit makes a DC 13 Morale test (no action required). On a success, the unit reforms with 1 casualty. If a 20 is rolled on the d20 for the test, the unit reforms with 2 casualties. On a failure, the broken unit disbands.

\begin{DndSidebar}[float=!b]{Rallying}
Each unit can be rallied only once per battle—either at the end of the battle using the Rally maneuver, or beforehand by the use of a martial advantage or some other special feature. At the end of a battle, after a victor has been declared, all broken units automatically attempt the Rally maneuver, unless they have already been rallied by some other means. (As a battle progresses, players should keep track of which units have already been rallied.)



If a commander has a reaction that allows them to rally a unit during another unit's activation (for example, as a reaction to being broken by an opposed unit's successful attack) the rally is executed immediately. If the unit successfully rallies, the unit that triggered the reaction then finishes its activation. (In the example above, it would proceed to the Power test, potentially inflicting more damage on the just-rallied unit.)



\end{DndSidebar}

\subsection{Retreat}
\textit{Prerequisite: None (any unit)}

This unit makes a DC 8 Command test. On a success, the unit is removed from the battle and its current casualties are recorded. On a failure, the unit suffers 1 casualty as opposed units take parting shots at it. If this does not break the unit, it is removed from the battle and its current casualties are recorded.

A commander must attempt the Retreat maneuver for all their units if the commander's army is defeated in battle. But a commander can also attempt the Retreat maneuver with any unit on its activation in order to save the unit from a potentially worse fate.

\subsection{Set for Charge}
\textit{Prerequisite: Infantry}

As a reaction to suffering a casualty from a successful Attack test from a cavalry or aerial unit that is not using the Archers trait, this unit makes a DC 13 Command test. On a success, the attacking unit suffers 1 casualty.

\subsection{Strafe}
\textit{Prerequisite: Fighter martial advantage}

As a reaction to succeeding on a Power test made as part of an attack against an opposed artillery or infantry unit, this unit makes a DC 13 Command test. On a success, two adjacent opposed units in the same rank as the target unit each suffer 1 casualty.

\subsection{Volley}
\textit{Prerequisite: Fighter martial advantage}

As an action, choose two adjacent units and make a DC 13 Command test. On a success, this unit attacks both units. On a failure, the unit attacks only one unit of its choice.

\subsection{Withdraw}
\textit{Prerequisite: Artillery or infantry}

As a reaction to failing a Morale test, this unit can move back into an empty space instead of suffering a casualty. If there is no empty space behind this unit, it cannot use this maneuver.

\section{Martial Advantages}
Between adventures and during downtime, player characters work with the units they command, whether training them or crafting items for them. Infantry units trained by a barbarian character fight quite differently than infantry units trained by a rogue. Likewise, wizards and clerics both craft wands and scrolls for the sergeants who command their troops, but these items grant very different boons.

These traits and items are called martial advantages, and are unlocked as the character's domain grows in size (from winning battles). A character's martial advantages typically apply to all the units that character controls as a commander, unless the text of the martial advantage says otherwise.

Multiclass characters can choose martial advantages from any of their classes. For example, a wizard/druid can choose from either the Wizard Martial Advantages table or the Druid Martial Advantages table each time they gain new a martial advantage based on their domain's size. However, those choices are permanent once made, and can't be swapped out.

Many of the martial advantages described in the tables below use the shorthand "DS" for the domain size of the commander's realm.

\subsection{Battle Magic}
Some commanders gain access to battle magic, allowing them to create magic items designed for war that can be distributed among the units they control. Only one item can be used on a commander's turn. These items are limited use, and expend their magic either immediately or at the end of the battle. Each martial advantage creates a single item per battle.

Whenever a unit is affected by an opposed unit's battle magic, in addition to resolving the other effects of the magic, that unit must succeed on a DC 13 Morale test or suffer 1 casualty.

\subsection{Barbarian Martial Advantages}
Barbarians favor lightly armored, highly mobile troops. These soldiers sometimes don't last long, but they can carve a path through any battlefield.


\begin{DndTable}{cXX}

Domain Size & Martial Advantages & \\
1 & Furious Assault, Mobility Training\\
2 & Berserkers\\
3 & Mobility Trap\\
4 & Rage Charge\\
5 & Blood Fever\\
\end{DndTable}

\subsection{Bard Martial Advantages}
As characters who excel at inspiring comrades in the middle of battle, bards make excellent leaders. They focus mostly on battle magic to inspire their soldiers or trick the forces of the enemy, but their troops also enjoy the benefits of being well taken care of.


\begin{DndTable}{cXX}

Domain Size & Martial Advantages & \\
1 & Song of Battles Won, Troops of Fame and Great Renown\\
2 & Scroll of Mass Hypnosis\\
3 & Scroll of Omund's Trumpet\\
4 & Tale of Heroes Bold\\
5 & Impassioned Speech\\
\end{DndTable}

\subsection{Cleric Martial Advantages}
As fighting priests, clerics master a dangerous combination of martial prowess and battle magic. They excel at healing troops and rallying them in battle—and the heavy infantry they train hit hard.


\begin{DndTable}{cXX}

Domain Size & Martial Advantages & \\
1 & Divine Rally, Wand of Healing\\
2 & Heavy Training\\
3 & Exorcizers\\
4 & Blessing of Strength\\
5 & Scroll of Mass Healing\\
\end{DndTable}

\subsection{Druid Martial Advantages}
Druids are dedicated to defending not just their groves but all of the natural world. Their battle magics muster the power of nature to protect their armies, and to devastate their enemies.


\begin{DndTable}{cXX}

Domain Size & Martial Advantages & \\
1 & Nature's Boon, Wand of Grasping Roots\\
2 & Scroll of Torrential Rain\\
3 & Lynx Soldiers\\
4 & Bark and Root\\
5 & Scroll of Storms\\
\end{DndTable}

\subsection{Fighter Martial Advantages}
Fighters excel in warfare. Their soldiers are the most highly trained on the battlefield, and they know an astonishing range of techniques to direct and engage troops in the middle of a battle.


\begin{DndTable}{cXX}

Domain Size & Martial Advantages & \\
1 & Heavy Training, Martial Rally\\
2 & Strafe\\
3 & Charge\\
4 & Volley\\
5 & Field Promotion\\
\end{DndTable}

\subsection{Illrigger Martial Advantages}
The illrigger is a fighting hellknight whose soldiers call upon fell powers to terrify their enemies. (A PDF containing the Illrigger class is available on the MCDM store at shop.mcdmproductions.com.)


\begin{DndTable}{cXX}

Domain Size & Martial Advantages & \\
1 & Death Commandos, Infernal Rally\\
2 & Terror Troops\\
3 & Scroll of Hellfire\\
4 & Execution\\
5 & The Ruin of the World\\
\end{DndTable}

\subsection{Monk Martial Advantages}
Monks work closely with the common folk, training them to defend themselves against cruel overlords. They rarely train heavily armed or armored troops, but their levies fight as well as any infantry, and cavalry fear all units trained by a monk.


\begin{DndTable}{cXX}

Domain Size & Martial Advantages & \\
1 & Hidden Reserves, Levy Training\\
2 & Mobility Trap\\
3 & Focused Resolve\\
4 & Like Water\\
5 & Mind over Body\\
\end{DndTable}

\subsection{Paladin Martial Advantages}
When the paladin calls, soldiers answer. Troops loyal to these dedicated warriors fight with fervor, and enemies fear facing them—especially cavalry trained by a paladin, which are among the most effective forces on the field.


\begin{DndTable}{cXX}

Domain Size & Martial Advantages & \\
1 & Righteous, Templar's Rally\\
2 & Scroll of Clarity\\
3 & Cavaliers, Hell's Hammer\\
4 & Blessing of Speed\\
5 & Scroll of Templar's Blessing\\
\end{DndTable}

\subsection{Ranger Martial Advantages}
Taking advantage of a ranger's experience in the wilds, ranger-trained troops are used to operating behind enemy lines and in hostile terrain. Moreover, the skill of their archer units are legendary.


\begin{DndTable}{cXX}

Domain Size & Martial Advantages & \\
1 & Archery Training, Scout Training\\
2 & Rough Terrain Training\\
3 & Find Cover\\
4 & Pin Them Down\\
5 & Coordinated Fire\\
\end{DndTable}

\subsection{Rogue Martial Advantages}
Soldiers trained by a rogue are highly mobile and well skilled in the use of dirty tricks. Rogue-trained scouts might not last long behind enemy lines, but they will wreak havoc while they stand.


\begin{DndTable}{cXX}

Domain Size & Martial Advantages & \\
1 & Mobility Training, Skirmishers\\
2 & Poison Arrows\\
3 & Scout Training\\
4 & Feint\\
5 & Infiltrators\\
\end{DndTable}

\subsection{Sorcerer Martial Advantages}
Specializing in battle magic, sorcerers craft scrolls and wands designed to inflict maximum damage to as many units as possible, and can take down even the most hardened soldiers.


\begin{DndTable}{cXX}

Domain Size & Martial Advantages & \\
1 & Sorcerous Training, Wand of Fire\\
2 & Invisibility\\
3 & Scroll of Translocation\\
4 & Fire Shield\\
5 & Scroll of Earthquake\\
\end{DndTable}

\subsection{Warlock Martial Advantages}
A warlock's dread patron empowers highly diverse battle magic, letting them reshape the battlefield, protect their soldiers, and terrorize enemies.


\begin{DndTable}{cXX}

Domain Size & Martial Advantages & \\
1 & Scroll of Illusory Soldiers, Wand of Acid Pool\\
2 & Patron's Curse\\
3 & Scroll of Blood Magic\\
4 & Flaming Hooves\\
5 & Scroll of Hell's Maw\\
\end{DndTable}

\subsection{Wizard Martial Advantages}
Few armies can stand resolved to fight while a wizard's lightning rips through them or opposed units appear from thin air to attack. A wizard specializes in battle magic that can both protect the troops they train and let their units unleash cataclysmic power upon the battlefield.


\begin{DndTable}{cXX}

Domain Size & Martial Advantages & \\
1 & Arrows of Dancing Lights, Wand of Lightning Storm\\
2 & Wall of Fire\\
3 & Scroll of Translocation\\
4 & Antimagic Shield\\
5 & Scroll of Cataclysm\\
\end{DndTable}

\section{Opposing Commanders}
When running the army controlled by the villain or some other NPC, the GM must track which units belong to which NPC commanders. Each opposing commander facing off against the characters-as-commanders acts in initiative order, just like the characters. When an opposing commander acts, they can activate all the units they control.

Each opposing commander can command a number of units equal to their proficiency bonus (just like player character commanders), with that proficiency bonus determined by the opposing commander's Challenge Rating.

For example, Lord Saxton's despotic regime has five officers—Magus Therin, Prelate Ellingwood, Sir Anglim, Lady Morgant, and Lord Saxton himself—who are the commanders of that domain's forces during warfare. Lord Saxton is CR 8, giving him a proficiency bonus of +3 and allowing him to control three units. The rest of his officers are CR 3, with their +2 proficiency bonus allowing them to each control two units.

In general, the GM decides which units each NPC officer controls before the battle starts, to suit their own ideas of how the battle should go. (The adventure in this book assigns units to specific officers, but even there, the GM is free to move things around.)

\subsection{Martial Advantages for Opposed Commanders}
The units of the villain's army gain the benefits of martial advantages just as the units of the characters' army do. However, the GM does not track which martial advantages each unit has based on their commander. (That would be too much to ask of the GM, who is already doing a lot!) Instead, all opposed units benefit from the same martial advantages, with the number of those advantages determined by the opposed domain's domain size. The GM chooses the martial advantages for the villain's army based on the general character types of any of the commanders of the opposed domain—spellcaster, knight, or assassin.

\subsubsection{Spellcasting Commander}
Enemy mages, shamans, necromancers, cult leaders, and grand viziers all qualify as enemy spellcasting commanders, as do any NPCs built around a class with spellcasting ability.


\begin{DndTable}{cXX}

Domain Size & Martial Advantages & \\
1 & Sorcerous Training, Wand of Fire\\
2 & Patron's Curse\\
3 & Scroll of Torrential Rain\\
4 & Fiery Defense\\
5 & Scroll of Firestorm\\
\end{DndTable}

\subsubsection{Knight Commander}
Any heavily armed and armored enemy leader qualifies as a knight commander. Spellcasters who wear heavy armor might also play the part of a knight commander instead of a spellcaster commander, at the GM's determination.


\begin{DndTable}{cXX}

Domain Size & Martial Advantages & \\
1 & Death Commandos, Heavy Training\\
2 & Cavaliers\\
3 & Charge\\
4 & Volley\\
5 & Well-Motivated\\
\end{DndTable}

\subsubsection{Assassin Commander}
Master thieves, expert rangers, and any other stealthy, highly mobile, high-damage-output NPCs qualify as enemy assassin commanders.


\begin{DndTable}{cXX}

Domain Size & Martial Advantages & \\
1 & Scout Training, Skirmishers\\
2 & Poison Arrows\\
3 & Mobility Training\\
4 & Execution\\
5 & Coordinated Fire\\
\end{DndTable}

\subsubsection{Choosing Commander Roles}
The GM has no obligation to tie martial advantages to the classes or NPC stat block names of an opposed domain's officers—only to let those roles inspire them as desired. Likewise, for monstrous commanders, the GM should focus more on the monster's role as a leader, rather than automatically assigning them the role of a specific commander based on their stat block.

This makes the characters' army more flexible than the villain's army, because the units of the characters' army will have a much wider array of martial advantages available to them. However, each of the player's units benefits only from a single set of advantages. Having an entire army run under one set of martial advantages creates many efficiencies that make the villain's army easier to run, and ultimately just as effective as the characters' army.

As an example, consider that Lord Saxton and  Lady Morgant are evil knights. Saxton has two spellcasters in his court (Magus Therin and Prelate Ellingwood) and an assassin (Sir Anglim). This means that between the leader and his lieutenants, Saxton has access to any of the three opposed commander types.

Saxton's domain is size 2, so the GM gets to pick three martial advantages (two for size 1 and one for size 2), which can be chosen from any of the three opposed commander tables—one advantage from each table, all three from one table, or any other combination.

\subsubsection{Expanding Martial Advantages}
In addition to the martial advantages gained from any of the three opposed commander tables, the GM is free to choose martial advantages from the character class tables as well. If one of a realm's officers is known to do magical deals with a mysterious supernatural figure, a martial advantage from the warlock table makes sense. However, a GM who does this should avoid any advantages that would allow a realm's commanders to rally troops in the middle of a battle. The ability to rally a unit during the fight is intentionally framed as a heroic action. And more importantly, if both sides can rally their troops in the middle of battle, that battle will drag on tediously and quickly stop being fun.

\section{Fortifications}
A fortification is a structure occupying a number of spaces on the battlefield, typically on the side of the force that is the defender in the battle. Units in those spaces gain benefits based on the type of fortification. The GM determines which fortifications are available to each side in the battle (both existing fortifications and fortifications raised during play) and in which spaces those fortifications appear.

Each fortification provides a bonus to Morale for all allied units of the side that controls the fortification, regardless of where those units are on the battlefield. (If there are multiple fortifications on the battlefield, allied units use the highest bonus to Morale for any of their side's fortifications.) As well, each fortification provides a bonus to Defense for any units in the same space as the fortification, and a tower fortification provides a bonus to Power for any artillery units in the same space. (Siege weapons cannot be placed on towers.)

The Basic Fortifications table provides a breakdown of fortifications, their sizes, and the bonuses they provide. All bonuses provided by a fortification are lost if the fortification is destroyed. A rank with a Fortification in it does not collapse when there are no units in that rank.

Units in fortifications can't be exposed.

\begin{DndTable}[header=Basic Fortifications]{XXcccc}

Fortification & Size & Morale* & Defense** & Power*** & Hit Points\\
Stone fence & 1–2 spaces & +1 & +2 & — & 4\\
Guard tower & 1 space & +1 & +2 & +2 & 6\\
Town walls & 3–4 spaces & +2 & +2 & — & 8\\
City gates & 5 spaces & +2 & +2 & +2 & 8\\
Keep & 2×2 spaces & +3 & +2 & +2 & 10\\
Castle & 4×2 spaces & +4 & +2 & +2 & 12\\
\end{DndTable}

\subsection{Moving}
Leaving a fortification costs no extra movement, but it costs 1 extra movement to move into a space with a fortification. The units defending a fortification can break like any other unit, or can use the \textit{Retreat} maneuver to depart the fortification (page 110). Doing so might allow an opposed unit to move into that space and gain the benefits of that fortification.

\subsection{Damaging a Fortification}
Only units with specialized traits can damage a fortification, with those traits specifically calling out fortifications in their descriptions. (See \textit{Unit Traits} on page 126.) Units with a specific siege trait—Siege Engine, Siege Weapon, Fast Siege Weapon, and Fast Siege Weapon (Heavy)—can attack and automatically hit a fortification, with no Attack test or Power test necessary.

\subsection{Contested Fortifications}
A fortification can be empty and still grant its bonus to Morale for the units of the army that control the fortification. But as soon as a unit moves into the area of a fortification controlled by the opposing side, that fortification becomes contested and no one gains a Morale bonus for it. If the only units in the area of a fortification are opposed units, then the fortification becomes occupied by that opposing army, and only the opposing army gains the fortification's bonus to Morale.

When two units attack each other and both are in the area of the same fortification, neither gains any of the fortification's bonuses to Defense or Power.

\subsection{Hit Points and Siege Weapons}
Fortifications have hit points based on their size. When a unit damages a fortification, that damage is assigned to the entire fortification, not any specific part of it. This means that fortifications occupying more than 1 space can't be targeted in any specific space, and they remain standing in all their spaces until the fortification is wholly destroyed.

When a fortification is reduced to 0 hit points, it is destroyed in all its spaces, and any unit in those spaces must make a DC 13 Morale test with disadvantage. On a failure, the unit suffers 1 casualty.

\subsection{Placement}
The GM decides each fortification's position on the battlefield, though fortifications are almost always on the defending side. Walls and gates do not have to be in the vanguard, but many scenarios will place them there to represent the idea that the defenders are behind those fortifications for the protection they offer. That being said, the GM might well decide to put the city gates in the center rank, so the setup of fortifications isn't meant to be taken too literally.

\subsection{Types of Fortifications}
The following are the most common fortifications, showing their typical placement on the battlefield. These are not the only possibilities, though, and the GM is free to use towers and walls in any combination to build unique fortifications, calculating those fortifications' hit points accordingly.

\subsubsection{Stone Fence}
This short wall provides protection for only one or two units. It might be used in a short battle in which a handful of units fight over a small town.

\subsubsection{Guard Tower}
This is a lone tower, atop which an artillery unit can leverage height and superior position to maximize the chance of inflicting casualties on enemy forces.

\subsubsection{Town Walls}
This larger version of the stone fence offers protection for even more units.

\subsubsection{City Gates}
A sturdy town wall is abutted by towers on either side of it! A large city gate is shown below, while a smaller one would have two towers with just a oneor two-space wall between them—and would have fewer hit points.

Each gate space is either a tower space or a wall space. Artillery units that move from a tower space to a wall space lose the bonus to Power that a tower provides.

Converter's Note: Apologies for the quality difference for this image. It had some weird transparency issues when I extracted it.

\subsubsection{Keep}
A keep is a collection of towers clustered closely together, and essentially acts as a large tower.

\subsubsection{Keep}
A castle represents the largest fortification typically found on a battlefield, built up of both walls and towers to provide maximum benefit for infantry and artillery.

\section{Terrain}
Terrain represents general effects that change the nature of the battlefield, from fog to rain to dense forest. Units of certain ancestries feature soldiers that are historically used to fighting in certain types of terrain, allowing those units to ignore the penalties terrain imposes. At the same time, certain terrain offers advantages for specific types of units.

\subsection{Fog}
Fog obscures vision and creates strange auditory effects. A sergeant's command shouted from the other side of the battlefield might echo through the fog to sound like it's coming from just a few feet away, making it difficult for troops to hear their orders.

Any artillery unit in fog has disadvantage on Attack tests. Any unit that rolls a 1 on a Command test while in fog becomes disorganized.

\subsection{Forest}
Trees act as natural cover, making archers almost useless in forest fighting. While in forest terrain, any artillery unit has disadvantage on Attack tests.

\subsection{Mud}
Mud, including terrain near or in a swamp, stymies movement. Any unit in mud must march in order to move 1 space, and it does not benefit from any bonuses to movement.

\subsection{Rain}
Rain obscures vision just like fog, and creates muddy ground to combine aspects of both those terrain types. While in rain, each unit must march in order to move 1 space, any artillery unit has disadvantage on Attack tests, and any unit that rolls a 1 on a Command test becomes disorganized.

\subsection{Plains}
The wide-open expanses of plains are excellent for cavalry units, providing them with maximum maneuverability. While in plains terrain, any cavalry unit has +2 to Command.

\subsection{Mountains}
Mountains feature treacherous slopes, steep cliffs, and falling rocks, and units unaccustomed to moving in such conditions routinely lose soldiers. While in mountain terrain, any unit that moves must succeed on a DC 8 Power test or suffer 1 casualty.

\section{Scenarios}
Scenarios change the shape of the battlefield, creating unique developments that might challenge forces on one or both sides, negate one side's advantage in numbers, or make some units useless!

\subsection{Bridge}
When two armies fight over a bridge, each side's vanguard rank is reduced to a single space, and cavalry units cannot be used. This scenario often comes down to a battle between artillery and aerial units.

\subsection{Defend the Pass}
In this scenario, one side occupies the pass while the other side must break the first side's defense, minimizing the field for one side. Cavalry units cannot be used in the battle, and artillery units can target only enemies in the same column. This sort of battle lends itself to a small army taking on a much larger one.

\subsection{Charge the Hill}
This scenario uses a normal battlefield, except the defender's rear rank is at the top of the hill and each rank in front of it is farther down the slope. Any unit moving uphill must succeed on a DC 8 Command test to march 2 spaces. On a failure, the unit moves 1 space instead. (If a unit's movement is limited, as when moving through mud, it might be able to march only 1 space and might not be able to move at all on a failure.)

Any infantry unit that moves 1 space downhill before attacking inflicts 1 additional casualty on a successful Power test made as part of an attack against a target unit that is downhill from them. Cavalry units attacking from the high side of the hill have advantage on Attack tests against any units.

This scenario allows a small number of units at the top of the hill to defend against a larger army at the bottom of it.

\subsection{Commanders on the Field}
Sometimes two armies meet in a field and neither side is defending anything. In this scenario, the commanders of the armies are leading their troops against each other, rather than fighting face to face.

Each commander chooses a unit to lead. If two units each led by opposing commanders end up adjacent to each other, a commander can forego issuing an order to their unit to instead challenge the opposing commander to single combat. This is an opposed Command test. The victor decides whether the single combat proceeds.

\subsubsection{Single Combat}
The battle is paused, and the two leaders enter combat together, rolling initiative again after each round. Either commander can yield as an action, in which case their unit suffers 1 casualty. If one commander drops to 0 hit points or dies outright, the unit they commanded immediately breaks and the single combat is over, though the battle can continue.

\subsection{Battle of Three Armies}
This scenario imagines two different allied armies attacking a third army. This scenario allows one large army to fight off two smaller armies.

The defending army has no rear rank, but it has two vanguard ranks and two reserve ranks. If one of those vanguard ranks collapses, units in the defending army's center rank and the reserve rank behind the collapsed vanguard are both exposed to the army the collapsed vanguard faced.

Converter's Note: Same issue here as with the fortifications, this one just isn't as intrusive so I left it alone.

\section{Unit Traits}
Over and above the statistics found on each unit card, a unit's traits define what it is able to do during battle. \textit{Unit Traits Descriptions} (below) breaks out every trait found across all units in Kingdoms \& Warfare, plus those in the \textit{special unit deck} (page 144). Each unit has the name of each of its traits printed on its card but not the full description, for two reasons. First, it would make the cards too large. And second, it would require players to pick a unit card up every activation to read what its traits can do. Easier to print the traits out for all players involved in the battle.

If the text of a trait conflicts with the baseline rules in this book, the text of the trait takes precedence.

\subsection{Recharge}
The notation "Recharge X–Y" means a unit can use a trait once, and that the trait then has a random chance of recharging during each subsequent activation. Each time the unit activates after using the trait, roll a d6. If the roll is one of the numbers in the recharge notation, the unit regains the use of the trait.

\subsection{Legal Target}
Where a trait refers to a legal target, this means any target that a unit could normally make an Attack test against.

\subsection{Unit Trait Descriptions}
Traits are presented in alphabetical order.

\section{Ancestral Units}
Many of the units that make up the armies in the warfare rules are organized by ancestry. Six of the most common ancestries are discussed below—humans, elves, dwarves, orcs, goblinoids (including hobgoblins and bugbears), and undead. These ancestries cover many of the peoples and forces that players and GMs might commonly recruit from, with other ancestries available through the special unit deck (see below).

Organizing units by ancestry is completely arbitrary, but still useful. If it makes sense for a specific campaign, any of these ancestries can be easily reskinned and referred to as specific nations, tribes, or domains. In an all-human campaign, dwarf units could represent a line of human mountain folk, while elf units stand in for the soldiers of a steadfast forest realm. In the same way, not all elf units need the same traits. A GM might develop three or four additional traits for different courts of elves in their world, so that only one court gets the Eternal trait.

Each of the standard ancestries features nine units, including a variety of infantry, cavalry, and artillery units covering Tiers I, II, and III. It's possible to build a solid army using just the units from those ancestries, making use of the characters' and NPCs' backgrounds and alliances.

\subsection{Units, Not Characters}
It's important to note that the way members of a given ancestry organize for war can be very different from how they perform during adventuring or when fighting a dragon single-handed. In other words, the stats a member of an ancestry has in combat are not related to the ancestry's stats in warfare. The things that make a good adventurer don't necessarily make a good soldier. Instead, these rules focus on how members of a particular ancestry might organize in large companies, how good they are at taking orders, how seamlessly they coordinate with dozens of other soldiers, and how they maintain their cool while taking casualties.

As such, the typical stats for an elf, dwarf, or zombie in the core rules don't translate mathematically into a unit's stats. Rather, a unit focuses on culture, training, and behavior as those things apply to the requirements of a large battle. By the same token, the soldiers represented by units are not adventurers. Though they have solid combat skill, their training is mostly focused on fighting together and not breaking under pressure. A battle in a dungeon or a dragon's lair is not a place for a hundred soldiers trained with pike and shield, accustomed to deploying in formation and to count on the soldiers in front of, behind, and beside them. Army battles take place in wide-open spaces to allow this sort of unit to make best use of its soldiers' training. Cram them all into a series of 10-foot-wide corridors without light, and it drastically reduces their effectiveness.

An army might have some kind of special operations section whose soldiers function more like adventurers, at least in terms of their skill sets. But for the most part, those sorts of soldiers are not the soldiers in a standard unit.

\subsection{Humans}
Humans are feeble and weak compared to most of the other ancestries in the book. They have no ancestral bonus to Attack, Power, Defense, or Toughness. What pathetic creatures! However, they are better than most other ancestries at two things: Morale and Command. And as it turns out, these are critical elements to success on the battlefield. Human units are thus less likely to break and run when the battle gets tough or things get weird. And they're better at executing complex commands, which means their infantry are more likely to succeed on maneuvers such as Follow Up and Set for Charge. Taken in combination with the Adaptable trait that all human units have—granting advantage on Morale and Command tests—and humans are an ancestry optimized for warfare.

Many are the commanders who've looked out at a pathetic band of spindly humans and imagined them to be easy pickings—only to experience real dread and terror when the humans don't break. They don't run. Rather, they stand fast. They hold their ground. And they think tactically. Each ancestry has its own specialty in the field of war. But for humans, their specialty is war.

Human armies are versatile. The best armies have a variety of units: infantry, cavalry, and artillery, and humans are equally good at all of these. This means you can field a balanced army, evenly divided between these disciplines. However, it's your army. You should muster the units you think are fun to field. Try an army with lots of archers! Or use your domain's Espionage skill to spy on the enemy, learn their army's composition, and muster units to thwart them!

If you don't know which units your enemy is fielding, try mustering the Hounds of Dalrath. They gain a bonus against most of the units the GM is likely to use.

\subsection{Elves}
Elf units are better on Attack than other ancestries in this book, representing their hereditary affinity with sword and bow. Elf units include some of the best archers in the game, though they are otherwise physically unimpressive on the battlefield. That said, elf units are highly mobile, making it easier for them to exploit any holes in an enemy's defenses. If a unit breaks or uses the Retreat maneuver, elf infantry will very likely move into position to take advantage of that gap.

As well, elf units have the Eternal trait. Undead and infernal creatures are no different to them than humans, dwarves, or goblins. As such, the soldiers of an elf unit don't experience existential dread at seeing a zombie horde shuffling toward them any more than they do in response to a human horde. This makes them better at holding their ground in the face of units that strike fear into the hearts of soldiers of other ancestries.

Elf armies favor archers! No surprise there. The Elven Thorns are one of the best artillery units in the game. Weakening opposed units with the Thorns is a powerful strategy likely to contribute to your victory. But the Forest Mist unit is also not to be underestimated.

Like many well-trained infantry units in the book, the Forest Mist gains a bonus to Defense when it moves into a gap left by an opposed unit after a successful Follow Up maneuver. In addition to the Forest Mist's high Defense, they can plunge into gaps of the opposed infantry without worrying too much about taking casualties as a result.

\subsection{Dwarves}
Dwarves are a hardy folk, granting their units greater natural Toughness than the units of other ancestries. This makes them excellent infantry, and dwarf infantry units are among the best in the game. Dwarf units also have the Stalwart trait, meaning that even after a dwarf unit is diminished, it is very hard to break it.

Infantry should be the meat and potatoes of any dwarf army, but dwarves field good artillery and cavalry too! Dwarves favor crossbows for their foot soldiers, and these are devastating at close range. But artillery must deploy into the center rank, so a skilled dwarf commander focuses on using their infantry to open some space in the front rank, and then move their artillery units adjacent to opposed units to maximize their close-range effectiveness.

\subsection{Orcs}
Orcs share the elves' affinity with weapons (granting a bonus to an orc unit's Attack) and have increased Toughness that makes them excellent soldiers. But their most potent asset is the Relentless trait, which allows an orc unit in danger of breaking to keep fighting and stay on the field. An orc unit reduced to the very last of its soldiers can keep fighting almost indefinitely!

Orcs are the traditional enemies of human, elf, and dwarf armies, known for their proud and fierce culture and traditions. But a domain that can avoid conflict to ally with orcs (typically in exchange for assurances of respect for territorial boundaries and mutual defense in times of need) can gain access to powerful units. An entire orc army with the Relentless trait is a thing few enemies are prepared to face. But orc units have a wide array of other traits.

The Bloodwalkers can outflank any enemy, making them the equivalent of super scouts. And the orc cavalry known as the Hellriders ride nightmares, which grant that unit the Harrowing trait. Most enemies find the presence of these hell steeds on the battlefield a serious blow to morale, and opposed artillery struggle to focus on targets whose Jaunt trait allows them to attack, then shift into the Ethereal Plane where they cannot be reached.

\subsubsection{Proxying Other Steeds}
Cavalry units gain no bonuses to their stats based on the mount they ride, but different mounts contribute unique traits to their unit. As you can infer from the Hellriders, nightmares add the Harrowing and Jaunt traits to this unit.

A GM who wants to give cavalry units different mounts should easily be able to look at existing mounted units in the game and reverse engineer what that mount does. Nightmares grant the Harrowing and Jaunt traits. Ram-mounted cavalry get the Ram Riders trait, units with elk mounts get the Spit Upon Their Horns maneuver, and so on.

\subsection{Goblinoids}
Often at the low end of the pecking order under bigger, stronger species, goblins tend to be pressed into service as levies more so than trained troops. But goblinoids also include the powerful bugbears and the hobgoblins, two peoples whose love of regimentation makes them exceptional soldiers. As such, the proper mix of goblins, hobgoblins, and bugbears on the field can prove a match for any other ancestry.

Bugbear units have huge bonuses to Attack and Defense (and are all Tier III units as a result). They also inspire other goblinoid units to greater deeds, so that an army of goblins with a bugbear unit leading them can be a formidable threat! For their part, hobgoblins are fearless combatants and sound strategists, giving hobgoblin units a better Attack than units of other ancestries, and making them the equal of human units for Morale. Hobgoblin units also make use of the Warbred trait, granting an additional attack after a successful Power test. This can be used to devastating effect as hobgoblin infantry chew through opposing forces—and makes those units easily the equal of humans on the battlefield.

Although goblins are typically no one's first choice for infantry units, their cleverness more than makes up for a lack of physical strength. Goblin Sappers make use of the Dig trait to undermine enemy fortifications, while Goblin Smokers use the Smoke Screen trait to leave opposed units unable to properly retaliate to goblinoid attacks.

An all-goblin army normally focuses on hobgoblins for infantry, artillery, and cavalry. Goblins are useful as levies and maybe one or two specialized units, such as Goblin Sappers if the enemy has a fortification.

Because the Bugbear Heavy Claw unit has traits that can grant free movement and attacks to goblinoid units, it's best to use that unit for an all-goblin army. If you want to spice up an existing army with some goblinoid units, hobgoblins are best.

\subsection{Undead}
Not an ancestry so much as a necromantic tool used by the commanders of evil domains, undead units feature in many NPC-led armies. Low-tier undead units are intended to be expendable, but higher-tier units are quite nasty—and many interesting undead units can be found in the special unit deck.

The baseline units of an undead army are skeletons and zombies. Skeleton units are destroyed relatively quickly, but zombie units can be as hard to take down as orc units. What makes all undead units especially formidable, though, is the Dead trait and the Harrowing trait. The Dead trait means an undead unit never worries about Morale tests and can't be diminished. Because its soldiers aren't alive, they don't worry about their fellow bonebags being ground up by enemy forces.

The Harrowing trait means that most units of mortal creatures are freaked out by the presence of undead on the battlefield. Elf units pay undead no mind, but the units of other ancestries don't like the prospect of fighting against their own forebears risen from the grave. A sense of dread and foreboding fills the field when undead march, and especially at the beginning of a battle, undead units can have a major advantage!

Undead armies shine once they get past Tier I units. The main purpose of fielding skeletons and zombies is to unlock ghoul, wight, and wraith units! Ghoul Infantry regain lost casualties thanks to the Feast trait, allowing them to feed on the corpses of fallen enemies or allies. And the ghoul artillery unit known as the Fireborn can inflict ongoing damage with fire tokens from their special maneuver.

Deathstormers are wight cavalry units that race across the battlefield on skeletal steeds. And the unique Tier III wraith unit known as the Black Wind is even more powerful. Along with the Dead and Harrowing traits, this unit gains bonus movement and can move through other units thanks to its Ethereal trait. This also means units that the Black Wind attacks gain no benefit from fortifications. If you've amassed a large enough army to field the Tier III units, you probably don't need that Zombie Catapult—but you should probably field it anyway, because it's dope.

\section{The Special Unit Deck}
Of course there was no way to fit every possible unit anyone might think of into a single book, so MCDM had to produce a supplement to Kingdoms \& Warfare: the special unit deck. This features sixty-four nonstandard units from a variety of additional ancestries and monstrous lineages, including dragons, fiends, more undead, kobolds, giants, black puddings, centaurs, gnomes, elementals, and more!

Ultimately though, the three tiers of units across six ancestries in this book, plus the special units that can be mustered by domains, are enough to run any number of interesting, unique battles without resorting to additional special units—especially when those battles include fortifications and different kinds of terrain.

Moreover, it's our hope that the warfare system and the setup of units is clear enough to make creating new units fun. There's no need to buy the special unit deck just to field a single dragon unit for a battle. Rather, just think about the system as presented, imagine how dragons might work in it, and create your own unit—which will probably end up very close to the official dragon units in the special unit deck!

\section{For the GM: Building the Villain's Army}
Just like the regular game focuses on the characters, these warfare rules are not a competitive exercise. It's still the GM's responsibility to curate content in order to create drama. So while the GM can use the same rules for building the villain's army as the players use to build the characters' army, it is perhaps more important that the army the heroes face is sufficiently challenging. Likewise, the villain's army should reflect the choices the players and characters have made. The heroes can always get in over their heads, bite off more than they can chew, and end up facing overwhelming odds on the battlefield—just as they can in combat with individual enemies.

The rules used for designing encounters with monsters apply to the warfare rules as well. Thus "more units" is typically a greater challenge than "better units," just like a standard encounter with many more monsters than characters can quickly become a TPK. The villain's army should obey the rules for unit command and dependencies (part of the \textit{Building an Army} section starting on page 100). But knowing that an army can field one Tier IV unit does nothing to suggest which unit is best for a particular scenario—and not all units of the same type are equal in battle.

Players can muster units only from ancestries they have connections or alliances with, or from ancestries predisposed to lend aid through good roleplaying. But the GM \textbf{should make liberal use of monstrous units}, as well as any units from any ancestry that seems reasonable and dramatic. The GM is not restricted by which ancestries they can field, as long as it seems reasonable that the leaders of an opposed domain would have access to those monsters or ancestries.

Deciding which units to field is not a science. It will take interpretation, guesswork, and experience to properly balance a battle—and to figure out that not all battles should be balanced. And this process doesn't stop once the battle begins! A GM might find that they need to discreetly fudge some die rolls or have units use less than optimal tactics because they realize only midway through the battle that they've overloaded the villain's forces—and that the characters' army is about to have a bad day.

In addition to the general advice for players, GMs should keep the following points in mind when building armies for a villainous realm or any other domain fighting against the characters:


\begin{itemize}
\item If one side has Tier V units, the other side probably should too. But a few heavily upgraded Tier IV units could balance things out.
\item As a general rule, avoid scenarios where one side has many more or much better units than the other side.
\item If one side has a super-heavy, super-elite unit, that's fine—as long as it's clear how the other side can make up for that. A good number of veteran medium units might be enough. Strategically, a good army is one with infantry, artillery, and cavalry units. About half an army should be infantry.
\item Tactically, an army should focus on breaking the opposing side's front. If one side's artillery and cavalry are used to try and get to the opposing side's center through the rear while infantry attack the opposed vanguard, that side will probably lose because it's dividing its forces.
\item Once an opposing side's front is broken, the attacking side's units can move around more. This opens up many different strategies, and lets cavalry and artillery units become more flexible.
\item The best use of archers is for picking off units that are close to breaking. (Assuming there are no aerial units to attack!)
\item It's reasonable for a GM to design a battle meant to be a rout of the characters' forces. But that should happen only if the players have gotten in over their heads in spite of the GM's best efforts to the contrary.
\end{itemize}

These rules can always be adjusted for drama and style. Perhaps a particular villain relies on a horde of infantry! But the rules are designed to make it so that an all-infantry army won't be able to bring enough troops to bear to be effective, and will likely inevitably lose to a better-balanced army. Still, GMs should always remember that as with the characters facing off against single foes, \textbf{the rules for warfare expect the players to face reasonable challenges and win}. The only time this should change (again, as with the regular game) is if the characters have gotten themselves in over their heads because they've done something foolish—particularly against the GM's advice.

\subsection{When Should a Monster Be a Unit?}
As a rule, opposed units should not be creatures so mighty that the characters would prefer to handle them personally. For example, a 5th-level party's army should not face off against a Red Dragon unit, both because that unit will probably smoke (perhaps literally) the characters' small, relatively inexperienced army—and because the players would most likely prefer that their characters fight that dragon face to face.

Eventually, though, as the characters level up, what was once a major threat becomes a minor issue. A high-level party with a large, experienced army might see a Black Dragon unit flying around the battlefield and shrug. "Our army can take care of it." These things are relative.

\subsection{Reinforcements: The Secret Weapon}
Because this is a brand-new system and no one is well familiar with it, players and GMs might find themselves in the middle of a battle that's going badly. "Going badly" means something undramatic is happening. One side is being routed, while the other side faces few challenges leading to an unearned victory. This scenario is bad no matter whether the player characters are the ones routing or being routed.

If this happens, it's a good time for the losing army to deploy their reinforcements. These are new units that arrive on the battlefield in the nick of time to save the day. (Cavalry are traditional reinforcements, but there's no need to be a martyr to tradition.)

Reinforcements should be used sparingly, and should always be a dramatic reveal. As just one example, an NPC domain whose leaders decided to stay neutral despite being courted by both the characters and the villain during intrigue makes an excellent source of unexpected reinforcements for either side. If the reinforcements come to aid the heroes, then the players and their characters gain control over them. The GM decides if reinforcements arrive and which units show up. This is a narrative tool used to keep the tension high, though, not a part of the battle setup.

\chapter{Monsters \& Magic Items}
\section{Monsters}
\DndDropCapLine{S}{}trongholds \& Followers introduced the idea of \textit{Concordance}, a numerical value every character possesses that measures how closely aligned they are with their god's or patron's teachings and desires. Anyone who worships a god or follows the teachings of a church or a greater power has Concordance. In a fantasy world where these beings are objectively real and grant their followers access to spells and abilities, any character should be able to petition their god and, if they are found sufficiently pious, receive aid. You can use the Concordance rules in Strongholds \& Followers to summon some of the creatures in this book.


\subsection{Action-Oriented}
Many of the creatures presented here are designed to be boss monsters: enemies that can take on an entire party by themselves or with a handful of minions. But boss monsters that are just "higher CR" often fail to challenge the characters. We've got a solution: \textbf{action-oriented monsters}.

The key to effective solo villains is giving them plenty of opportunities to act and move when it's not their turn. Thus, each action-oriented creature in this book has at least one special bonus action and reaction, as well as "villain actions" that let them dominate the battlefield.

These actions make the boss monsters action-oriented, dynamic, and formidable. Whether fought as an exciting solo challenge or alongside a few easy-to-run minions, action-oriented monsters challenge the characters with dramatic and powerful actions in combat.

\subsubsection{Villain Actions}
Every action-oriented monster has three villain actions that they can use after an enemy creature's turn. Villain actions are similar to legendary actions with the following exceptions:


\begin{itemize}
\item A creature can only use one villain action per round (as such, villain actions tend to be more powerful than legendary actions).
\item Each set of villain actions has a recommended round order. These abilities give the battle a logical flow and a cinematic arc.
\end{itemize}

The first villain action is an \textbf{opener}, which shows the characters that they are not battling a typical monster. Openers generally deal some damage, summon a minion, buff the creature, debuff the characters, or move the creature into an advantageous position. They're just a taste of what's to come.

The second villain action provides \textbf{crowd control}. It typically fires after the heroes have had a chance to respond once or twice, get into position, and surround the villain. This second action helps the villain regain the upper hand. Like an opener, this action comes in many flavors, but it is even more powerful than an opener.

The third and final villain action is an \textbf{ultimate move} or "\textbf{ult}"—a showstopper the villain can use to deal a devastating blow to the characters before the end of the battle.

While every creature has a recommended order of actions, you can take villain actions in the order most dramatic for your fight. You could push back a villain action if a monster is stunned or might stay alive for more than three rounds, or you could perform the recommended third action in round two after several surprise critical hits!

\subsection{Court of Decay}
The Last City, a hulking necropolitan ruin in the center of the Abyss ruled by the Demon Prince of Undeath, is home to insidious undead creatures— called the Deathless and the Soulless, they are creatures of hate, memory, and mind who refuse to yield to law and surrender to the grave.

The Court of Decay and the Court of the Deep maintain a taut, uneasy truce. The Deathless have no souls, so they yield no power to the demons of the depths. But they do not fall into \textit{lethe} (page 165) and therefore pose a constant threat.

Ancient kings of the Mundane World, these Deathless Lords mastered unholy lore in life so their unfailing will would overcome the sleep of death.

\subsubsection{Zaar Anathema, the Shackled King, the King in Gold}
Zaar, the Thane of Kalas Mithral and Lord of House Q'or, once ruled the empire of steel dwarves and personally arbitrated disputes between the Celestial domains and the Terran Empires. He was a wise king, just and true. Now Kalas Mithral is dead, forgotten, and he who once ruled now serves in the Court of Decay.

The end of Kalas Mithral came with the war against ██████. The astral celestials' first Army of Night sieged the walls of Kalas Mithral for months and weeks and days, six and six and six. The Legion of Adamant was mighty in warfare and impervious to sorcery and sword. But the greatest weapon the astral celestials wielded was treachery and corruption, and the Legion of Adamant, once called the Legion Unyielding, faltered.

The King in Gold foresaw the ruin of his house, his capital, and his people. He swore a vow that Kalas Mithral would not fall while he lived and in answer came a howl from below. A bane placed deep within the earth at its creation called out to Zaar. Long had the elementals known that Ord placed hideous blights within his masterwork, entombing them within the world. Fringe sages among the dwarves wrote that all Orden was a prison for something poisonous to the timescape.

Zaar summoned the Order of Fabrication, a team of elite mechanists. He pointed to the floor of their deepest chamber and uttered one command: "To the heart of the world." And thus the Order of Fabrication undertook their greatest making.

A mighty engine did the mechanists build, its corkscrew tip sharp enough to pierce the shell of the world, its hull clad in a metal never seen before or since, its core able to withstand the titanic pressures and temperatures. Thirty elementals could the vessel hold, engineers to man the controls and Zaar's elite wardens to battle whatever they encountered below, whatever guarded the bane at the heart of the world. War-ender did they name the great machine, for they believed the power below would end the wretched siege and bring their people salvation.

None now live who remember the War-ender's expedition, but this much is known: many days after it pierced the shell of the world, the Warender emerged, its hull shattered, its crew gone. Only Zaar remained, manning the controls alone.

And gripped tightly in his hand, exhumed from the depths of the earth, was the \textit{Tellac Oranic}, a stone that radiated death and turned the living into shadows.

At first, the lords of the other houses believed Zaar had unearthed their salvation. For three days Zaar led the Legion of Adamant with the verdant stone in hand, and the Army of Night melted away before him.

But the other houses quickly realized the price. While the Army of Night retreated, the dwarves of Kalas Mithral died, turned into dark umbras by the hideous stone. What seemed at first a victory was instead devastation—the Army of Night simply withdrew while the \textit{Tellac Oranic} feasted on legions of life.

Zaar did not heed his noble court's pleas as they begged him to surrender the stone. When he turned his back on them, they declared his house Anathema and stripped him of his rank and title. But Zaar did not notice this shame.

Eventually, the Legion of Adamant was gone, their souls fed to the stone. Only as he stood alone on the battlefield, the Army of Night an oncoming flood of death, did Zaar realize his folly. He called out to Ord, to Aan, to Myr and Kul, the makers of earth and air, of water and fire, and the four gods heard him. But their answer was not mercy—it was condemnation.

Sea, Sky, Mountain, and Magma cursed Zaar. And Ord, maker of the world, shackled the King in Gold, bound the \textit{Tellac Oranic} to him for eternity. As his flesh and mind decayed, he died howling, watching the mile-high walls of Kalas Mithral, which had never known tarnish nor blemish, fall to Night.

The star elves slaughtered the steel dwarves to the last, sparing none, and pulled their gleaming silver city down around their corpses. Zaar, the Shackled King, fell into the Unthinking Depths, where he drags the green stone through the Last City, remembering nothing of his days of glory—only the death of his people and the king who failed them.

\textbf{Zaar Anathema}

\subsubsection{Khorak, Jarl of Drengrheim}
When the Nine Clans of Vanigar rose in the north, Khorak the Stone Giant warned the other giant tribes that humans would spread, multiply, and drive the giants into extinction. Some agreed, and though they joined Khorak against the spread of the clans, their forces were not enough and the Vanirmen repelled them.

Khorak took his remaining followers into the arctic wastes of Vanigar and there made bargain with Olkhar, the Vanir God of Trickery. Olkhar gave Khorak and his followers the power they needed to conquer the Vanirmen, but his bargain cost them an eye. A ragged band of giants entered the frozen desert, but a cyclopean horde emerged—a force whom none could stand against.

Armed with Olkhar's gift, a prophetic power called \textit{deadsight}, Khorak and his followers won many battles against men, and more giants joined his growing legion. Khorak personally slew each of the chieftains of the Nine Clans in battle and claimed the lands of men. An empire of giants emerged, called Drengrheim, with Khorak the Far-Seeing at its helm. But these victories and this empire could not sate Khorak's hunger, fed by the visions of the \textit{deadsight}. It told him everything, showing endless lands and treasures to claim—all while hiding from him the means of his undoing.

He turned against the other giants, seeking to control their resources, and did what the humanoids could not: united the clans of giants against him. The Cloud, Fire, and Ice tribes united under the banner of Jarl Manghild, a fire giant princess Khorak orphaned and thought left for dead. In the end, Khorak died in battle against the jarls of the combined tribes. Never before had they allied, and never would they do so again.

\textbf{Khorak, Jarl of Drengrheim}

\subsubsection{Telcezalco XXVIII, the Blood River}
Telcezalco was the last king of the inantzicatl, the serpent folk of Orden. The serpent king's devotees gladly followed him into wars on other islands, worshipping him as a god. But Telcezalco was mortal. Every passing day reminded him of that fragility.

The inantzicatl believed the terror lizards of the island—the tetzahuitl—were sacred. They offered food, jewels, and other sacrifices to the tetzahuitl before hatchings, long journeys, and battles to ensure good fortune, and never harmed the lizards lest they risk angering the gods. Telcezalco suspected their blood contained great power—power over life and death. He was right...but the only death he gained power over was his own.

Telcezalco and his followers slaughtered thousands of tetzahuitl, pouring a river of blood that they bathed in. They emerged immortal and undying, but their victory over death was short-lived.

Defiling the sacred pact between the inantzicatl and the tetzahuitl woke Goxomoc, the god of the terror lizards. Telcezalco's newfound power was nothing in the face of Goxomoc. The titan lizard decimated the island and the inantzicatl empire. No more would the snake-men of Ix build or smelt or write, for Goxomoc destroyed all they had. The few survivors left Ix in search of more habitable lands.

Telcezalco, twenty-eighth and last, ruled over the death of his empire. Now he serves in the Court of Decay, and all fear the arrival of the final king.

\textbf{Telcezalco XXVIII}

\subsubsection{Máelodor Rhyllearnán, King of the Foxglove Court}
\DndQuote%
{}
{''"Men will be the end of this world.''}
{Máelodor Rhyllearnán}
Long before the first empire of man, Máelodor Rhyllearnán watched the Gol spread across the land. Where once only elves, dwarves, and dragons lived, now a fourth mortal people walked among them.

They did not farm, nor build roads, nor cut wood, all the offenses elves would later levy against men. But they worshipped many gods. Máelodor refused to share his kingdom with the humans, so the Gol went to war. They had warred against elf and dragon and dwarf—they had even warred against each other. Máelodor didn't scare them. Such defiance filled the king with rage.

"They will consume the world," Máelodor said.

"Leave them. They will consume themselves," the Queen of the Foxglove Court observed.

But the king became obsessed with the Gol. Their expansion unnerved him, and he spoke incessantly about his never-ending war with these humans. The Queen in Digitalis urged her husband to seek a truce and appeal for peace. Máelodor ignored her wisdom. After years of watching her people die for Máelodor's obsession, the Queen of the Foxglove Court left in the night, leaving her husband to his lonely war. She and her guard and favored courtiers slipped away in the moonlight and established the Orchid Court, which rules the Great Wode still.

Máelodor raged. He claimed the Gol ensorcelled his wife. It was their fault the battle continued, and now it was their fault he fought this bloody conquest alone. He howled a vow to his god that day, and Val heard his cry. The Oath of the Hunt bound the King in Digitalis to his own fury. On the longest night of the year, he gathered his devilhound cavalry and began the Solstice Hunt, the hunt against humans.

For weeks did the Solstice Dogs hound the Gol, their bows and hawks and jaws hewing men before them like weeds razed from a garden. Máelodor's forces caught the Gol surprised and pressed the advantage.

Finally, a Gol leader summoned the demon Taamgul, who carved into men sorcerous tattoos. These Gol, their skin still bleeding from the arcane ink, rode against the Solstice Dogs with Taamgul at the lead and cleaved the hunt's devilhounds out from under them.

Taamgul slew Máelodor Rhyllearnán, the King in Digitalis, but the King's vow sustained him. Dead, he fought on, and Taamgul was forced to embrace him. When she returned home to the Abyssal Waste, she dragged Máelodor with her.

Máelodor Rhyllearnán now serves in the Court of Decay, but legends tell of powerful necromancers summoning the ancient king, and his dread cavalry sometimes rides out against all humans, thirsting for vengeance.

\textbf{Máelodor Rhyllearnán}

\subsubsection{Tristan Vaslor, the King Below}
Three thousand years before the Council of Aberdanon, men from the Commonwealth flooded into the region then known as the Central Plains. They were organized and bore advanced technology. From Khoursir they brought the stirrup, and from Kalas Mithral they brought the secret of steel. Riding at the head of their armies was Tristan Vaslor, a warrior unmatched. In a search for great power, he led them to war with the Gol, the first humans.

But the Gol were lorecrafty and called upon the demons Khuurtok the All-Seeing, He of Five Hundred Eyes; Bytorr the Wolf; and Ullrok the Blighted, Consort of Flies. These demons fought alongside the Gol and pushed back Tristan Vaslor and his army. An uneasy truce settled, but Tristan Vaslor desired conquest.

The Central Plains were wild then, populated with hidden temples to unknown gods and mystic towers where hermit sages delved into knowledge man was not meant to know. The wizard Pellum lured Tristan to the White Tower, a tower of necromancy, and offered him aid in exchange for service. "Seek you the Temple of Primordial Chaos," the wizard called out from atop his tower. "There you will find the sword \textit{Maethelgas}, which shall aid you in your domination. And also a book I desire, one of little value to you."

Tristan and his knights rode out, found the temple, and made to sack it, but the temple was vast and many things hostile to light and life dwelled and warred within. For one year did Tristan's knights siege the temple, building a makeshift city around it, delving deep, disappearing for weeks, until Tristan emerged with the codex Pellum desired...and a great black iron runesword in his hand.

Pellum was true to his word—bound to \textit{Maethelgas}, Tristan was able to drink the souls of his enemies. Many Gol leaders did he slay, and each time \textit{Maethelgas} grew in power.

Finally, the vast plains were conquered. The region was named Vasloria, and Tristan ruled for 300 years, his life extended by the powerful sword at his side. But the blade required souls—Tristan's conquest could not cease. Soon the humanoid tribes were brought under his sway, and this was not enough. The elves and dwarves did he fight, and still the blade called for more.

Eventually, Tristan turned against his own dukes and fed their souls to the blade. After this massacre, his knights damned him, entombing him in the crypt he built for his many dead wives.

For three days, the knights of Vasloria thought their nightmare ended, but in his desperation Tristan fed his own soul to \textit{Maethelgas}. The crypt quaked with a necrotic surge, splitting asunder as Tristan emerged from his crypt as the Dead King.

For ten years, the King Below, the One Below All, ruled over Vasloria. Under his dread hand, Vasloria descended into a nightmarish realm of cruelty and suffering. The people cowered and lived in fear. None dared stand against their dread sovereign.

Not until two brothers, Adun and Cavall—one a farmer, the other a soldier—preached a new way of living. Not to remain slaves to fear, but to resist, to fight, even if it meant death. Even the grave was better than this wasteland they endured.

They led their growing congregation to the Dead King's castle and there confronted him. While none imagined they might live, much less prevail, Adun and Cavall prayed. Their prayer rose in righteous light and banished Tristan Vaslor to the Abyssal Wasteland, and as his undead howling faded from castle walls, the rule of the King Below ended and Vasloria knew peace.

The One Below All once ruled a continent, and in power he remains after death—he serves in the Court of Decay at the right hand of the Lord of Decay, the Ultralich, \textbf{Khorsekef the Infinite}.

\textbf{Tristan Vaslor}

\subsubsection{Khorsekef the Infinite}
First Pharaoh of the Khemharan Empire, Khorsekef was immortal and ruled for over 3,000 years. It was impossible for the Khem-hor to imagine a world without their god-king. Even now, five millennia later, the fashions and traditions Khorsekef ushered in are synonymous with Khemhara. The cometary impact of his rule still echoes loudly in the desert.

Under Khorsekef's rule, the people of Khemhara thrived in the desert. Khorsekef used his power to ensure this, providing an oasis for an entire nation. What his people could not know, because the desert isolated them from the world, was that Khorsekef's rule denied them the right to change or grow. They worshipped their god-king, who never died, and he watched while generation after generation turned to dust at his feet.

Children were born, grew old, married, had children, but nothing ever changed. Names, likenesses, even identities, inventions, and discoveries were endlessly recycled. The Khem-hor thought this was the cycle of the world. The truth was hidden from them.

At least, they did not until the coming of the great explorer Lady Agatamori from Higara, who crossed the desert in the ingenious sand ships her engineers built. Lady Agatamori, a cartographer eager to chart new lands, did not bring war, but she caused one. A civil war. A rebellion. A war against time.

Khorsekef captured Lady Agatamori and her retinue almost immediately and hid them in one of his pyramids, called hedrons. But it was too late—word spread. There were other humans in the world. Humans who strove, created, invented. Humans who were not bound to the Khem-hor's eternal cycles of rediscovery and reinvention, recycling the same knowledge, even the same identities, through eternity.

Still, this might have remained a rumor, soon forgotten, but for the captain of the Pharaoh's guard who was so overwhelmed by the brightsouled humanity in Lady Agatamori that he smuggled her and her crew out of the Great Tet and into the Atenopolis.

Like a virus, understanding began to spread. The people of Khemhara had been lied to. They were not alone in the world, and they did not have to die repeating the same generation over and over.

Armed with new knowledge from conversing with the Khem-hor, Lady Agatamori's vizier unlocked the final clue. Khorsekef's hedrons were not, as he had told his people, mausoleums for noble citizens. None were buried there. Instead, the hedrons had been set in certain astronomically significant locations for a secret purpose. The vizier thought he knew why.

While the Pharaoh's loyal guard closed in on Lady Agatamori's hiding place in the city, brave Khem-hor brought word of the vizier's discovery to the Heironauts, the Pharaoh's cadre of wizards. They quickly assembled the clues and realized they, too, had been lied to. Their position brought so much privilege and luxury that it was easy to trick them into thinking their limited knowledge was the sum total of world lore.

The Heironauts discovered the great lie: the hedrons were engines that stored time, absorbing it from the topology of the manifold, a time-looped reality created by the Pharaoh, and transmitting it into the Infinite Pharaoh to preserve him while stagnating everything else. Nothing in Khemhara changed because nothing could change. Time itself had stopped. Each generation wasn't merely like the previous—it was a single population giving birth to itself, burying itself, over and over through eternity. They thought their pyramids were tombs. But the desert was a tomb, and they were the embalmed.

Powerful magics erupted in the desert as the Heironauts broke the oppressive cycle and made massive strides in lore. Within days they discovered new spells, and soon reached a limit imposed by the hedrons. So they focused their lore on the hedrons and began to drain their power for themselves, using the stored time to fuel their magic.

The Pharaoh, their god-king, could not allow this. He used the power given him by the hedrons to bring battle against his own wizards.

Cause grappled with effect as the War Against Time engulfed the Atenopolis. Though the Pharaoh initially proved stronger, the Heironauts were growing more powerful by the hour. Soon, the very fabric of the manifold would rip, spilling the Atenopolis out into the Real.

The Pharaoh then did the unimaginable: He opened the Great Tet. The Great Tet was not like the other hedrons. The other pyramids stored time like batteries, releasing it to perpetuate the manifold Khorsekef created. But the Great Tet had never released its stores of time. The Pharaoh broke the seal on the massive hedron, freeing the millennia contained within. Even the Heironauts had not known of the deep stores of time kept inside—it was an emergency measure Khorsekef had built exactly for this contingency.

Opening the Great Tet, the Infinite Pharaoh drank deep of the eons stored within, but its power was so great no physical vessel could contain it. His atoms burst with stored time and shredded apart, leaving only his will.

The white-eyed, dessicated form that remained was the animated memory of the Will of Khorsekef. He was no longer the Pharaoh. He was now the Ultralich, a mind so powerful no force has yet been discovered that can end him.

Facing annihilation by the newly dead god before them, minds filled with new spells and ideas from Lady Agatamori and those of their own creation, the Heironauts did something no citizen of the Atenopolis had done in a millennium. They improvised. They dared stand against the pharaoh's stagnant rule.

"The gate swings both ways!" the First Heironaut declared, and her council understood. Using the residual power stored in the remaining hedrons, they invented a ritual to negate the Infinite Pharaoh's millennia-long program.

They took the geometries of the manifold he had created, the cul-de-sac of reality they had lived in for uncounted centuries, and inverted it. They spilled its contents across the timescape, and the Atenopolis reentered the Mundane World.

At the same time, the fabric of the pocket universe they lived in wrapped back against itself, against its center. The point of origin. Khorsekef, the being that had created it.

Caught within his own imploding manifold, Khorsekef was hurled downward into the Abyssal Wasteland. There, he ascended as the Lord of Decay and the King of Death. The Ultralich, commanding the loyalty of the Court of Decay. They skirmish with the other courts while Khorsekef seeks a way to return to the Atenopolis and restore his infinite rule.

Khorsekef vowed he would return, and the oracles of Khemhara predicted exactly that for the last several millennia. But in the meantime, an explosion of learning and innovation propels the Khemhor forward. The Pharaohs of the desert weave the Eternal Design in preparation for the Infinite Pharoah's return. If they can complete it before he returns, Khorsekef will find nothing but sand and dust and scarabs where once was a bright civilization—the Khem-hor will have saved themselves and found a paradise.

\textbf{Khorsekef the Infinite}

\subsection{Court of the Deep}
The Abyssal Wasteland is a dead plane filled with ruins of once-great cities. Or are they the degenerating reflections of existing cities in the timescape? Sages differ. Demons call this place home, this rotting reality where forgotten gods go to die. It is a land of entropy and the creatures that call the Waste home consider order a poison.

This is the deepest world in the timescape. Creatures who sink this far down find themselves unable to organize their thoughts, unable to plan or scheme beyond base emotion and survival. Only by feasting on the souls of others can a mind rise above this entropy and muster order enough to plan and, possibly, escape.

Cruel and unyielding, the demons who rank in the Court of Depths gorged on souls until their minds were sharp enough to plot and scheme and climb over one another in a desperate attempt to rise out of the depths and perhaps, one day, loose themselves on the Mundane World, where souls are rich and potent.

\subsubsection{Devourers of Souls}
Unique among monsters, demons of the Court of the Deep feast not upon food and water, but upon \textbf{souls}. Souls fuel their bloodthirsty powers; while starved for souls, a demon can scarcely even think.

Similar to how monster hit points are formatted, the number of souls a given demon has consumed is presented both as a die expression and as an average number. The amount of dice used to calculate a demon's soul count is dependent on the demon's challenge rating, as shown in the table below.

\begin{DndTable}[header=Soul Dice]{cccX}

CR & Soul Dice & Average (Rounded Down) & \\
1–10 & 1d4 & 2\\
11–15 & 2d4 & 5\\
16–18 & 3d4 & 7\\
19+ & 4d4 & 10\\
\end{DndTable}

During an encounter, a demon can burn souls to use or enhance certain abilities, decreasing their soul count. Demons gain more souls by slaying other creatures. When a demon reduces a creature to 0 hit points, the creature must immediately make a DC 11 Wisdom saving throw before it falls unconscious and begins dying. On a failed save, the creature's soul is consumed by the demon. The demon adds 1 to their soul count, and the creature immediately dies and subsequently can't be restored to life by any means short of a wish spell.

Soul devouring makes the demons of the Court of the Deep deadly opponents. You can remove this ability and only have a demon devour a soul when a creature dies through normal means within 60 feet of it or adjust the DC of the saving throw to give characters a better chance of living through an encounter with one of these opponents. Likewise, you can make the deadliness of soul devouring scale with a demon's power, requiring a saving throw with a DC equal to 8 + the demon's proficiency bonus or equal to the demon's challenge rating.

There is no upper limit to a demon's soul count. When a demon's soul count drops to 0, it falls into lethe (see \textit{Lethe}).

\begin{DndSidebar}[float=!b]{Souls and ActionOriented Monsters}
The Court of the Deep monsters use villain actions and soul devouring in combat. While actionoriented monsters are made to be fought alone, having a few low-challenge-rating demon or cultist minions in the encounter gives the big bad some souls to suck up! Consider adding some of these creatures during an encounter with a member of the Court of the Deep.



\end{DndSidebar}

\subparagraph{Soul-Fueled Power}
Each demon has a swathe of features that require souls to use. These can be broken down into two broad categories: passive features and active features.

\textbf{Passive features} are triggered whenever a demon's soul count crosses a certain threshold. The threshold is unique to each feature and detailed in the feature's description. The feature remains active until the demon's soul count drops below the feature's respective threshold.

\textbf{Active features} require a demon to burn souls and lower their soul count to use. Some active features buff common statistics, such as attack and damage rolls, while others grant the demon access to special—and oftentimes incredibly powerful—new actions.

Each active feature has a base cost, or the number of souls that must be burned in order to use the ability. This cost is noted in parentheses beside the feature name. Additionally, some active features allow demons to heighten the effect by burning additional souls on top of the ability's base cost.

\subparagraph{Soulsight}
Demons have a special sense called soulsight, which allows them to detect any sentient creature within a certain radius. Unless otherwise stated, neither physical matter nor magical effects can impede soulsight; however, creatures who lack sentience or have no soul are invisible to a demon's soulsight.

The radius of a demon's soulsight is denoted in their respective stat block.

\subparagraph{Lethe}
When a demon runs out of souls, it falls into lethe: a state of craven hunger, wherein a demon can only lash out in a desperate search for sustenance. Demons who have fallen into lethe are mindless and violent creatures whose only goal is to consume.

While a demon is in lethe, it is subject to the following effects:


\begin{itemize}
\item Its Intelligence score drops to 3 (−4), unless it is already lower.
\item It has advantage on all attack rolls and disadvantage on all Intelligence, Wisdom, and Charisma saving throws.
\item It can't cast spells.
\item On its turn, the demon automatically moves toward the nearest creature it can see with its soulsight, regardless of the terrain in its way. If the demon is still able to take an action that turn, it uses its action to attack that creature (or any obstacle in its way).
\end{itemize}

A demon is no longer in lethe when its soul count increases above 0, at which point the effects of lethe end.

\begin{DndSidebar}[float=!b]{Adding Soul Rules to Existing Demons}
This section presents the stat blocks for six of the most prominent demons in the Court of the Deep. However, a court is nothing without its subjects, and there are bound to be other demons roaming the planes, all thirsting for souls.



To add soul rules to an existing demon, such as a balor or vrock, first use the Soul Dice table to calculate the demon's initial soul count and increase its challenge rating by 1. Then give the demon two or three additional abilities it can use by burning those souls, such as the following:




\begin{description}
\item[Abyssal Shield (Costs 1+ Souls).]
As a bonus action, the demon can burn any number of souls and increase its AC by an amount equal to the number of souls burnt for 1 minute.

\item[Soul Rend (Costs 1–5 Souls).]
When the demon makes a successful attack against a creature, it can burn up to 5 souls and deal an extra 14 (4d6) damage for each soul burnt.

\item[Hungering Strike (Costs 1+ Souls).]
When the demon makes an attack roll and misses, it can burn any number of souls and add a +5 bonus to the attack roll for each soul burnt, potentially causing the attack to hit.

\end{description}



\end{DndSidebar}

\subsubsection{Groyle Fleshender}
Barely able to form coherent thoughts, Groyle has souls enough to hate and hungers for more. A dullard and a brute, Groyle has no capacity for cruelty—he has not the imagination or patience. He rends and demolishes his way across the battlefield, seeking out whichever creature's soul shines the brightest, endeavoring to consume them whole and alive. Their screaming desperation gives their immortal soul a sweetness the Fleshender can't resist.

Groyle begins combat by using Goring Charge to attack a nearby spellcaster and follows it up with a bite attack, preferring to target and swallow a creature with healing spells. The demon then focuses his attacks on the toughest enemy, using Mark of Abhorrence and Soul Rend to bring down foes quickly.

\textbf{Groyle Fleshender}

\subsubsection{Hara'antar the Soulthief}
Hara'antar the Blasphemite is the last acolyte of the Blood Sun cult. Effulgent in her unholy worship, she intends to complete the bloodshadow rite, assume the mantle of the High Priest, and start the cult anew.

The soulthief's prayers are unique in the Courtof the Deep; with tolling reverence empowering her hissing orisons, she can fuel her dark spells with the souls of nearby enemies without needing to consume them first. This makes her a pariah among the demons; if Hara'antar can tap into the dozens or hundreds of souls each great demon carries, what power could she gain?

\textbf{Hara'antar the Soulthief}

\subsubsection{Zor'yal Lifeswallower}
A cloud of blood surges across the battlefield. Within its crimson deluge is the Lifeswallower. The sanguine torrent cloaks Zor'yal from other demons' soulsight—without this invisibility, Zor'yal is weak. But do not underestimate this deadly assassin. Her claws are so potent, anyone caught unawares will die to their sting.

Zor'yal led the assault on Kham, her demons warring with the products of the Lords of Kham's beast-magic until the \textit{codex incabulum} was opened and the demons banished. Zor'yal is less now than she was, but she remembers. Her insatiable hunger for revenge against the Mundane World drives the gory squall that carries her into battle after battle until she can wreak her vengeance on the world above.

\textbf{Zor'yal Lifeswallower}

\subsubsection{Vorg'aut Harrowfist, the Lady Blight}
The Lady Blight is the greatest warrior of the Wasteland. She leads the Corrosion, a unit of elite demon heavy infantry responsible for singlehandedly repelling the Inexorable invasion of the Waste after the rest of the Depths broke morale and began to feed on each other. The Corrosion stood alone against the advance of Integral, Axiom's premiere cavalry—unable to sunder Vorg'aut's adamant stand, Integral was forced to retreat.

With hulking fists radiating crimson light, Vorg'aut commands the Wasteland, equal parts leader, warrior, and siege engine. Even the battlefield yields to her will, for the Lady Blight is unbending.

\textbf{Vorg'aut Harrowfist}

\subsubsection{Trall the Mindbreaker}
When sages across the timescape dare to turn the oracular eye downward, peering into the Abyssal Wasteland and seeking forgotten lore that long ago sunk beneath the surface of the Real, they find Trall's ivory void eyes, purging and all-consuming, staring back at them.

The Mindbreaker is attuned to any divination targeting the Waste and, rather than trawling the depths looking for souls with some spark left in them, she stabs leeching psychic tendrils through the arcane conduit, draining the soul from the poor wizard, sorcerer, or warlock foolish enough to peer into the Deep.

Thus does the Mindbreaker stay sated with souls. An augur unmatched in foresight and cruelty, she is both mage and anti-mage.

\textbf{Trall the Mindbreaker}

\subsubsection{Sylt Bloodheat}
The greatest demon of the Court of the Deep, Sylt commands the grudging obeisance of fiends within the Abyssal Wasteland. Unlike the others, Sylt has never fallen into lethe—she has transcended it, subjugated it, until the very flow of souls is hers to dominate.

Sustained by the Bloodheat, Sylt drinks souls in a never-ending flow, feasting on that which empowers the demons and rewarding the legions at her side. Demons under her command find their powers and intellect grow as their supply of howling souls is refueled, effortlessly sustained in a banquet given to her most loyal.

The Court hates Sylt—she is a dark and fell intelligence who prefers to wait and scheme rather than succumb to temporary hungers. They crave bloodshed and feasts of souls and sinew, impatient with her cunning and insidious plots. But none gainsay her right to rule. Many in the Waste wait for her word to lead them across the timescape in a bloody war of anti-life.

In combat, Sylt prefers to use hit-and-run tactics in environments with plenty of cover and places to hide.

\textbf{Sylt Bloodheat}

\subsection{Court of Seven Cities}
The Seven Cities of Hell scheme and plot endlessly, each ruled by an archdevil seeking domination over the others to become the King of Hell. Should this happen, the legions of Hell would erupt across the timescape, bathing the upper worlds in black soulblood.

Each archdevil sends an emissary to serve in the Court of Seven Cities, as required by the Onyx Edict, the founding document of Hell. But the court is infested with servants of these representatives, each one a point of connection in the ever-shifting web of deceit, subterfuge, and bloodshed.

\subsubsection{True Names}
All devils have true names that they keep secret, as speaking the true name of a devil within earshot of it can strip it of its supernatural defenses. True names are hard to discover—fiends go to great lengths to bury them. But with enough effort, your Domain could deliver one to you to use as a weapon at a crucial moment in the eternal war against evil. Proclaiming the true name of a devil takes one action, and if the devil can hear you, the proclamation removes its Magic Resistance trait, its damage immunities, and its damage resistances for 24 hours.

\begin{DndSidebar}[float=!b]{Why Seven Cities?}
Each of the Seven Cities of Hell is named after a real "hell" (or hell-like analog) from a real culture in our world.



This continues a tradition going all the way back to the original game in 1974, a tradition I have a lot of affection for. The idea that \textbf{our world}, the real world, Earth, is just another plane in the multiverse and it is "orbited" by other nearby Earth-adjacent planes, such as Arcadia, Hades, Asgard, and Niflheim.



Earlier editions imagined your cleric might worship gods from our real world! Thor or Ra, for example. Because it presumed those gods lived in their own planes, and those planes were accessible from any Prime Material plane, like Earth or your homebrew fantasy world. Earlier editions even mentioned characters who traveled from our real world to the fantasy RPG world.



So the MCDM multiverse, the timescape, imagines that an intrepid planar traveler could arrive in these hells and reimagines them as cities warring through eternity in the larger plane of Hell.



These hells are from many cultures, including Greek (Styx, Acheron) and Jewish tradition (Sheol), as well as Mongolian tradition (Kasyrgan). Naraka is from Hindu, Sikh, and Buddhist traditions. Finally, Dis is from Christian literature (Dante's Inferno).



These cities are named and included thus because I thought it was a cool opportunity to learn more about different cultures' notions of the afterlife. It was fun, doing this research. It should not be taken as an opportunity to imagine any of these cultures are in any way lesser. They should not be used as an opportunity to trot out hateful, backward, racist tropes. These were behaviorally modern humans like you and me. They led complex internal lives. They weren't ignorant, they just knew different things than we do.



And, hopefully, by representing different cultures among these hells that some readers may not be familiar with, some folks might read more about them! Learn how different real-world peoples throughout history have thought about the afterlife. That stuff is super cool and also, I think, makes us better GMs, because we are able to more realistically present our own worlds and peoples.



In the end, these aren't meant to be the real, authentic places described in our earthly cultures, merely a fantasy version appropriate to a pulp fantasy RPG.



—MC



\end{DndSidebar}

\subsubsection{Hate Devil}
Hate devils are brutal, ferocious infantry. They radiate antipathy so strongly that anyone in their aura begins to hate their friends, their allies—even themselves.

During the first round of combat, a hate devil attacks from above, using Spit Flame or casting \textit{phantasmal killer} to harm foes that resist its Aura of Enmity. On subsequent turns, the devil lands to focus melee attacks on a flying or ranged combatant before taking back to the skies for defense. The devil is willing to provoke an opportunity attack or two if it means returning to the safety of the skies.

\textbf{Hate Devil}

\subparagraph{Rakat, Baron of Styx}
Styx, the City of Blood, was sacked in the War of Forgetfulness and is now least among the Seven Cities. Their representative in the court is a lowly baron—but a baron who, since the war, is determined to not forget what Styx suffered.

A hate devil of sizable power, Rakat intends to pit the court against itself and destroy it from within. Their blinding rage hides a scheming mind.

Like other hate devils, Rakat begins combat by flying above their enemies, but remains low enough to ensure they are in the fiend's Aura of Enmity. On their first turn, they cast \textit{power word stun} on a foe with lower hit points than others to ensure the spell's success, and fire their Ray of Antipathy at a creature outside their Aura of Enmity. They use the Divide and Conquer villain action to cut any bard, cleric, or other supportive spellcaster off from a creature affected by \textit{power word stun} or charmed by their Aura of Enmity.

On subsequent turns, Rakat engages characters in melee with \textit{Zealcleaver}, though the devil continues to remain at least 10 feet in the air, benefitting from their long reach.

\textbf{Rakat, Baron of Styx}

\subsubsection{Stone Devil}
Nearly indestructible, stone devils are the bulwarks of the Seven Cities. No unit of stone devils has ever lost morale in the history of the timescape. Their hearts grind spirits into black soulblood that corrupts everything it touches.

A stone devil relies on its high AC to protect it from opportunity attacks as it moves around the battlefield to make the most of its Soulblood Trail trait.

The flaw imposed by the devil's Soulblood Trail trait and Foul Ichor action is a temporary affliction that is fun to roleplay for many players, but others truly dislike having a personality change imposed by a failed saving throw. Only use the flaw if you play with the type of players who enjoy incorporating it into the narrative.

\textbf{Stone Devil}

\subparagraph{Moranon, Marquis of Acheron}
Loyal and unyielding, Marquis Moranon steadily serves Asmodeus, archdevil of Acheron, the City of Whispers. Moranon's fidelity is admirable, Asmodeus just wishes his servant was a little more ambitious...

Unlike other stone devils, Moranon's villain actions and Disciplined Advance reaction allow him to create hazardous areas using his Soulblood Trail without provoking opportunity attacks. He is less reckless than his soldiers and makes far greater use of Profane Flame, since he can use it as a bonus action or as a reaction using Disciplined Advanced.

Moranon focuses his attacks on arcane spellcasters, believing they are the ones most likely to discover his true name and use magic to bypass his high AC. His advance is calculated and considered. If armored warriors get in the devil's way, he cuts them down, using Disciplined Advance and villain actions to get around such foes and engage spellcasters first, and using Profane Flame to attack those casters if warriors lock him down.

\textbf{Moranon, Marquis of Acheron}

\subsubsection{Steel Devil}
The razors of hell, steel devils are highly mobile terror weapons deployed to wreak maximum carnage in battle. It is impossible to fight one without suffering wounds; even psionic attacks cause the mind to bleed.

A steel devil uses its incredible speed to engage its enemies in melee combat, since its mere presence can dish out damage to any nearby creature. It prefers to fight warriors foolish enough to remain close to make the most of its Agonizing Visage and Thorns of Pain traits. Even at a distance, the devil's Agonizing Visage and Soul Scalpel can harm foes as it advances.

\textbf{Steel Devil}

\subparagraph{Sryz, Count of Naraka}
It's rare for a steel devil to enter the ranks of Hell's nobility, but Count Sryz sees the weakness brought about by the Wars of Forgetfulness and intends to quickly ascend to reigning as the prince—or to destroy the court in the process.

Sryz is the ultimate steel devil, and acts much like their counterparts do in combat. The devil engages other creatures in melee, surrounding themselves with as many combatants as possible to get the most out of their Agonizing Visage and Thorns of Pain traits and setting themselves up to devastate as many creatures as possible with their Sonic Attack villain action.

On their turn, Sryz spreads the damage around with their six Eviscerating Strikes, dashing with Lethal Agility if need be, then targets a creature they can't reach with Soul Scalpel. Sryz prefers to use Lacerating Presence on a foe that concentrates on a spell to make it more difficult for the creature to maintain a hold on the magic.

Sryz damages as many creatures as possible to maximize the effect of their Howl of Suffering villain action.

\textbf{Sryz, Count of Naraka}

\subsubsection{Lore Devil}
Lore devils are living compendiums of spent scrolls and dead codices. Keepers of lost arcane secrets, they wield cunning magics and feed on the energy of weaker mages.

A lore devil prefers to battle foes at range, beginning combat with its Multiattack action, first casting \textit{dominate person} to bring a strong bodyguard to the fiend's side, then using Arcane Bolt twice. The devil continues to use Multiattack, casting \textit{disintegrate} to harm an enemy spellcaster or \textit{reverse gravity} followed by \textit{fireball} and \textit{lightning bolt} on subsequent turns to damage as many foes caught in the upward pull as possible. If a creature gets close enough to make melee attacks against the lore devil, the fiend uses Spelldrain to regain hit points before casting \textit{misty step} or \textit{plane shift} to get away.

Note that the lore devil has a damage vulnerability uncommon among the other denizens of hell-fire. Characters don't normally encounter lore devils in flaming wastelands or volcanic lava pits the way they might some other fiends. These devils are smart and stay out of the places in hell that are hazardous to their health.

Some lore devils study ancient infernal rituals that allow them to bind the flesh of other devils to their papery skin. This fiendish essence infuses their pages with durability against flames, removing the vulnerability. Some lore devils gather enough fiendish flesh to perform the ritual several times, making them resistant or even immune to fire damage.

\textbf{Lore Devil}

\subparagraph{Indix, Earl of Sheol}
Lord Indix is composed of the burning pages of the once-powerful \textit{codex mutabilis}, the Book of Change. When the arcane tome was thrown into the caldera of Mount Vor, the book was incinerated, and deep within Sheol, the City of Death, Lord Indix was created.

Indix is the Book of Change given life—as such, they have access to now-abolished rituals, ones of immense and unmatched power. The Lords of Sheol, the City of Death, turn a blind eye to Earl Indix's machinations, fearing the lord's power even as they rely on it.

The Earl of Sheol is a complex boss monster, and their villain actions are maximally disruptive! To keep the game rolling, consider doing the math on these actions ahead of time. In the case of Transmute Aspects, having new stats and attack modifiers ready for your players to use will dramatically reduce the chance that they riot and set your table on fire!

\textbf{Indix, Earl of Sheol}

\subsubsection{Fang Devil}
Fang devils are the high priests of Hell, channeling unholy power to twist and subvert their enemies, mutating them into unrecognizable monstrosities. Horns erupt from their skulls and twist to pierce their eyes, their flesh splits open to reveal bloody gashes filled with teeth—blood that turns to poison as it drips from the maws in profane salivation.

\textbf{Fang Devil}

\subparagraph{Horat, Duke of Kasyrgan}
Horat is hierarch of the Envenomed, a cult of pain. His wicked power twists and corrupts everything it touches.

The Envenomed have spies throughout the Seven Cities, but thus far, Horat has been loyal to Pharyon. Many in Hell's cognoscenti believe there is a real alliance between the Duke of Kasyrgan and the Prince of Dis, though what this portends, no one knows.

Horat uses Multiattack to attack up to two spellcasters or lightly armored creatures with his Mouth Hand then attempts to poison those targets with Venomous Bayonet. He finishes his Multiattack by using Twist Flesh on a foe in his Mutation Field that isn't grappled by him.

If several foes get near Horat, he uses Hail of Teeth to deal as much damage as possible, even attacking any minions that might be in the line of fire.

Horat uses his Allow Me to Improve You reaction as often as possible to keep the toughest threats at bay while he tears apart foes with fewer hit points.

\textbf{Horat, Duke of Kasyrgan}

\subsubsection{Eye Devil}
Oracles of the Court of Seven Cities, the eye devils are inward-looking prophets. They speak in riddles, and only the Cult of the Mirrored Eye can decipher the proclamations of the Eyes of Hell.

An eye devil attacks from above, casting \textit{dominate monster} to bring the best melee combatant among its foes to its side, followed by \textit{incendiary cloud} on the fiend's second turn. Since the devil's Multiattack action allows it to cast spells and make attacks, it uses Flame Jet and Incinerate to attack enemies that can reach it. The devil uses Malvision on the first turn of combat, hoping to create maximum chaos among its enemies during the battle.

After casting its most powerful spells, the eye devil targets flying or ranged characters that can threaten it with the \textit{eyebite} spell. The eye devil concentrates attacks on one at a time, methodically bringing down threats.

\textbf{Eye Devil}

\subparagraph{Pharyon, Prince of Dis}
Once again, the power of Dis rises and the City of Lies rules over Hell. Dispater dispatches his newest lieutenant, Pharyon the Edictor, Revelator, Lord of Answers, Oracle to the Court.

For the first time in the (admittedly now ruined) lists of Hell, an eye devil leads the Court of Seven Cities, and the devil's dark portents guide the unholy court. The eye devils' power to see past, future, and alternate realities makes them rather unpredictable on the battlefield, but Pharyon is even more cryptic and inward-looking than most of his kin.

Pharyon carries the \textit{Omega Toll}, the Bell of Light's End. When it sounds its thunderous knell, the old epoch ends. What comes after...even Prince Pharyon can't see.

In combat, Pharyon flies to avoid melee and uses Malvision followed by Multiattack to attack with Flame Jet and a spell, typically \textit{dominate monster}, \textit{incendiary cloud}, \textit{delayed blast fireball}, or \textit{eyebite}. When a creature hits Pharyon with an attack, he uses his Look Again reaction, provided the attacker is in a relatively safe spot and not surrounded by the devil's foes.

\textbf{Pharyon, Prince of Dis}

\subsection{Gemstone Dragons}
A long-forgotten branch of the draco genus, \textit{draco crystallus} are primarily hexapods like their more common cousins \textit{draco metallus} and \textit{draco chromaticus}. They are older than the other draconic genera, which accounts for their wider diversity in body type.

You can read more about these psionic dragons in \textit{Strongholds \& Followers}. This book presents two unique gemstone dragons: Trudy, a ruby dragon wyrmling, and Cthrion Uroniziir, a mighty onyx dragon.

\subsubsection{Psionics}
The gemstone dragons are silico-organic organisms. When hatched, they are almost entirely organic. As they grow, their crystalline structure grows, as does their psionic power.

The crystals that grow on and thread through their bodies form a network of psionic filaments that capture, conduct, amplify, and store charges of psionic energies, which the dragons then convert into manifestations.

Though manifestations may seem like magic, they are an entirely different discipline and power. Spells such as \textit{antimagic field} and \textit{dispel magic} have no effect on psionic manifestations.

Each dragon begins with a number of charges based on their age. They can expend these charges to cast any psionic manifestation. Many of these manifestations have a base cost and can be enhanced by expending more charges.

Over time, the dragon's crystals conduct more psionic energies into its body, recharging its power. At the end of each of its turns, roll the recharge die and restore that many charges.

\subparagraph{Manifestations}
Psionic charges are used to manifest psionic abilities. These manifestations are the primary source of a gemstone dragon's power. All manifestations have a range of 30 feet, unless noted otherwise.

The following manifestations are by no means the only ones the dragons have access to. It is known that the most powerful gemstone dragons can alter their form at will in a manner that defies all inspection. But these are the most welldocumented and commonly used manifestations available to all psionic creatures.

The save DC against a dragon's manifestation is 8 plus the dragon's proficiency bonus plus the dragon's Intelligence modifier.


\begin{itemize}
\begin{multicols}{4}
\item Amplify
\item Another World
\item Believe
\item Distance
\item Elsewhere
\item Flay
\item Forget
\item Mindscape
\item The Real
\item Reflection
\item Sympathy
\item Weight
\end{multicols}

\end{itemize}

\subparagraph{Auras}
A dragon projects an aura 30 feet around it. The aura persists while the dragon has psionic charges remaining. The aura's effect depends on the dragon's gemstone variety.

\subparagraph{Crystal Hide}
Normal weapons bounce and spark off the crystals protruding from a gemstone dragon's skin. These dragons have resistance to nonmagical bludgeoning, piercing, and slashing damage.

\subparagraph{Sympathetic Vibrations}
All gemstone dragons are vulnerable to psychic damage. The same crystalline structure that makes them hard to damage with normal weapons conducts any psionic radiation, causing its latticework to crack and shatter.

\subparagraph{Mindspeech}
All gemstone dragons develop telepathy as they age. Most acquire it when they are young, but the topaz dragons are born with it.

\subparagraph{Matrix Mind}
The gemstone dragons can maintain several persistent manifestations simultaneously, but doing so requires spending 1 charge per manifestation at the start of each of its turns. Whenever a dragon takes damage while concentrating on a psionic manifestation, it must make an Intelligence saving throw to maintain concentration as if concentrating on a spell.

\subsubsection{Trudy}
When the great Cthrion Uroniziir rose from her slumber, many \textit{draco crystalli} were woken, hatching and crawling forth. One of them, Trudy the ruby wyrmling, is new to the struggle to preserve reality against the onslaught of the Würm of the World's End. But though she is, Trudy knows her duty, eager to join the battle and stop the Time Ender from consuming whole realities.

Trudy is still young as the crystal dragons reckon these things, but she is already an expert on the fabric of reality, the many worlds of the timescape, and the people who inhabit them. She talks to herself, constantly analyzes everything and everyone, and is a little oblivious to people's reactions to her. When she gets nervous, she breaks into rhyming couplets, often in Draconic. This can't be helped.

\textit{Quod est,}

\textit{Quod licet,}

\textit{Radix Malorum!}

\textit{Ad infinitum, quod erat demonstrandum!}

\textbf{Trudy}

\subsubsection{Cthrion Uroniziir}
Known to the elves as the Würm of the World's End, Cthrion Uroniziir, Lord of the Neutral Evil Dragons rises from her slumber to collapse the myriad manifolds of the timescape into one partition. If she succeeds, worlds will be annihilated. Collapsed together into one singular \textit{universe}.

Opposing her, a faction led by the slumbering Ballisantirax whose power is so great, should she wake, it would unmake the Mundane World. So she dreams, and from her dreams projects allies to aid the heroes who oppose Cthrion Uroniziir, the Time Ender.

Uroniziir is a world-ending campaign final boss requiring the greatest heroes of the timescape to join forces!

\textbf{Cthrion Uroniziir}

\subsection{Relg, the Descender, the Lord in Corpulect}
Before the Wars of Forgetfulness, Relg ruled the Court of the Deep as the King in Lethe. Almost invincible, Relg radiated lethe. No matter the souls they contained, no matter how powerful their will, any demon who opposed Relg fell into the unthinking entropy. He ruled for an eternity.

Until the coming of the Bloodheat, the Demon That Does Not Forget. Sylt's immunity to lethe meant she alone could stand against the Descender. Sylt overthrew Relg in a bloody, merciless clash and formed a new court in the foundations of blood left behind.

Relg was adrift in the timescape until the coming of Ajax, Invincible, the Iron Saint. Having heard the dark whispers about the Lord in Corpulect, Ajax summoned Relg, bound the once-king to him, and placed Relg at the head of the army of demons he now commands, the Infernium.

In combat, Relg positions himself among as many enemies as possible to make the most of his Aura of Lethe. The demon welcomes melee combatants and always tries to maintain a soul count of 3 or higher so his Pus-Filled Soul Boils harm those who threaten him.

On his turns, Relg first uses his tentacles to attempt to grab at least one creature so he can use his Reel trait, then uses his axe and horns to threaten other targets he can reach, using Soul Rend early and often to destroy his foes quickly. If a flying creature threatens the demon, he uses Weighted Soul to bring it to the ground.

\textbf{Relg}

\section{Magic Items}
Whenever we present an entirely new system (and this book contains two of them: intrigue and warfare), it seems like a good idea to present new magic items that support these systems. If organizations can act the way heroes do, and if characters can command units, then surely wizards and priests have developed magic to enhance these things!

This section presents new items that help officers better manage their domain and commanders better lead their troops. In addition, we include nine new codices, artifact-level magical tomes that bend the fabric of the timescape.


\begin{itemize}
\begin{multicols}{3}
\item \textit{Coat of the Crimson Guard}
\item \textit{Counterpoint}
\item \textit{Fate's Needle}
\item \textit{Grandmaster's Quiver}
\item \textit{Helmet of Dual Fates}
\item \textit{Heroic Soul's Chasuble}
\item \textit{Maethelgas}
\item \textit{Omega Toll}
\item \textit{Staff of Providence}
\item \textit{Tellac Oranic}
\item \textit{Zealcleaver}
\end{multicols}

\end{itemize}

\section{New Codices}
Because they were popular in \textit{Strongholds \& Followers}, we present nine new reality-bending codices. Each codex is an artifact-level item, written when the world was young and magic not yet tamed. Many codices grant their users incredible power, including the power to raise new units. They should not be looked at as rewards, but rather immense burdens. Threats to the natural order. These tomes are best found locked away in some obscure vault, guarded by a powerful entity who rightly fears what might happen should the codex fall into the wrong hands. Or any hands! Mortals, certainly, can't be trusted with this power. A high-level NPC might ask the heroes to retrieve one after a powerful enemy has wrested it from the bowels of obscurity. For a time, the heroes can use the power within—but this should always be temporary.

\subsection{Encountering a Codex}
Each codex grants whoever attunes to it enormous power. Before the players ever encounter a codex in the hands of a formidable enemy or sealed away for protection, you should make sure they have some idea what these books can do. Many of these abilities are dramatic and even over-the-top, but not wholly useful in most adventure scenarios. Even still, it becomes unfair, even unsportsmanlike, to surprise your players by making them the target of one of these abilities without warning them what the codices can do first.

These warnings could be presented as in-game knowledge conferred by an NPC or discovered through the heroes' own research.

\textit{Example: If you attempt to read the mind of a character attuned to the \textit{astral codex}, your head explodes. No saving throw. In one sense, this is an enormously powerful ability and conveys the severe stakes at hand. But the players should know before they encounter an NPC attuned to the Book of Stars that anyone who wields it has learned the mind mastery technique of the star elves, and what this means. It doesn't need to be specific: "Their minds can't be read. Anyone who attempts it winds up with their brain scrambled ...or worse!" When the heroes are forewarned, this alarming ability is much less ridiculous and startling and much more dramatic and foreboding.}

\subsection{Attuning to a Codex}
To attune to a codex, the reader must spend a month in uninterrupted study of no less than 8 hours per day, six days per week.

Once attuned, the owner of the codex need not take the book with them. They can leave it in a secure location with no limits to distance, but must refresh their knowledge through 8 hours of uninterrupted study once per month or lose its benefits.

\subsection{Sacrifice}
Some of these books require the attuned owner to surrender some of their life essence to satiate the fell powers constrained or described in the texts. Sacrificing requires rolling some number of Hit Dice, losing that many hit points and the Hit Dice immediately, and reducing your hit point maximum by the same amount. While the effect powered by the sacrifice persists, the lost hit points and Hit Dice can't be regained. When the effect ends, you can heal the damage, restore your hit point maximum to normal, and recover the Hit Dice normally.

If sacrificing Hit Dice causes you to drop to 0 hit points, you fall unconscious and all effects you created end. You can't sacrifice more Hit Dice than you have. If you must sacrifice a Hit Die to sustain an effect but you have none, the effect ends.

\chapter{Magic Items}

\chapter{Creatures}
\setlist[description]{nolistsep,listparindent=3pt,topsep=6pt,font={\bfseries\sffamily\small}}
\pagestyle{empty}

\begin{DndMonster}[float*=b,width=\textwidth + 8pt]{Vampire}
\begin{multicols}{2}
\DndMonsterType{Medium undead (shapechanger), lawful evil}
\DndMonsterBasics[
armor-class = {16 (natural armor)},
hit-points = {\DndDice{17d8 + 68}},
speed = {30 ft.}
]
\DndMonsterAbilityScores[
 str = 18,
 dex = 18,
 con = 18,
 int = 17,
 wis = 15,
 cha = 18
]
\DndMonsterDetails[
saving-throws = {Dex +9, Wis +7, Cha +9},
skills = {Perception +7, Stealth +9},
damage-resistances = {necrotic; bludgeoning, piercing, slashing from nonmagical attacks},
senses = {darkvision 120 ft., passive perception 17},
languages = {the languages it knew in life},
challenge = 13,
]
\DndMonsterAction{Shapechanger}
If the vampire isn't in sunlight or running water, it can use its action to polymorph into a Tiny bat or a Medium cloud of mist, or back into its true form.

While in bat form, the vampire can't speak, its walking speed is 5 feet, and it has a flying speed of 30 feet. Its statistics, other than its size and speed, are unchanged. Anything it is wearing transforms with it, but nothing it is carrying does. It reverts to its true form if it dies.

While in mist form, the vampire can't take any actions, speak, or manipulate objects. It is weightless, has a flying speed of 20 feet, can hover, and can enter a hostile creature's space and stop there. In addition, if air can pass through a space, the mist can do so without squeezing, and it can't pass through water. It has advantage on Strength, Dexterity, and Constitution saving throws, and it is immune to all nonmagical damage, except the damage it takes from sunlight.

\DndMonsterAction{Legendary Resistance (3/Day)}
If the vampire fails a saving throw, it can choose to succeed instead.

\DndMonsterAction{Misty Escape}
When it drops to 0 hit points outside its resting place, the vampire transforms into a cloud of mist (as in the Shapechanger trait) instead of falling unconscious, provided that it isn't in sunlight or running water. If it can't transform, it is destroyed.

While it has 0 hit points in mist form, it can't revert to its vampire form, and it must reach its resting place within 2 hours or be destroyed. Once in its resting place, it reverts to its vampire form. It is then paralyzed until it regains at least 1 hit point. After spending 1 hour in its resting place with 0 hit points, it regains 1 hit point.

\DndMonsterAction{Regeneration}
The vampire regains 20 hit points at the start of its turn if it has at least 1 hit point and isn't in sunlight or running water. If the vampire takes radiant damage or damage from holy water, this trait doesn't function at the start of the vampire's next turn.

\DndMonsterAction{Spider Climb}
The vampire can climb difficult surfaces, including upside down on ceilings, without needing to make an ability check.

\DndMonsterAction{Vampire Weaknesses}
The vampire has the following flaws:

\textit{Forbiddance.} The vampire can't enter a residence without an invitation from one of the occupants.

\textit{Harmed by Running Water.} The vampire takes 20 acid damage if it ends its turn in running water.

\textit{Stake to the Heart.} If a piercing weapon made of wood is driven into the vampire's heart while the vampire is incapacitated in its resting place, the vampire is paralyzed until the stake is removed.

\textit{Sunlight Hypersensitivity.} The vampire takes 20 radiant damage when it starts its turn in sunlight. While in sunlight, it has disadvantage on attack rolls and ability checks.

\DndMonsterSection{Actions}
\DndMonsterAction{Multiattack (Vampire Form Only)}
The vampire makes two attacks, only one of which can be a bite attack.

\DndMonsterAction{Unarmed Strike (Vampire Form Only)}
\textsl{Melee Weapon Attack:} 9 to hit, reach 5 ft., one creature. \textsl{Hit:} 8 (1d8 + 4) bludgeoning damage. Instead of dealing damage, the vampire can grapple the target (escape DC 18).

\DndMonsterAction{Bite (Bat or Vampire Form Only)}
\textsl{Melee Weapon Attack:} 9 to hit, reach 5 ft., one willing creature, or a creature that is grappled by the vampire, incapacitated, or restrained. \textsl{Hit:} 7 (1d6 + 4) piercing damage plus 10 (3d6) necrotic damage. The target's hit point maximum is reduced by an amount equal to the necrotic damage taken, and the vampire regains hit points equal to that amount. The reduction lasts until the target finishes a long rest. The target dies if this effect reduces its hit point maximum to 0. A humanoid slain in this way and then buried in the ground rises the following night as a \textbf{vampire spawn} under the vampire's control.

\DndMonsterAction{Charm}
The vampire targets one humanoid it can see within 30 feet of it. If the target can see the vampire, the target must succeed on a DC 17 Wisdom saving throw against this magic or be charmed by the vampire. The charmed target regards the vampire as a trusted friend to be heeded and protected. Although the target isn't under the vampire's control, it takes the vampire's requests or actions in the most favorable way it can, and it is a willing target for the vampire's bite attack.

Each time the vampire or the vampire's companions do anything harmful to the target, it can repeat the saving throw, ending the effect on itself on a success. Otherwise, the effect lasts 24 hours or until the vampire is destroyed, is on a different plane of existence than the target, or takes a bonus action to end the effect.

\DndMonsterAction{Children of the Night (1/Day)}
The vampire magically calls 2d4 swarms of bats or rats, provided that the sun isn't up. While outdoors, the vampire can call 3d6 wolves instead. The called creatures arrive in 1d4 rounds, acting as allies of the vampire and obeying its spoken commands. The beasts remain for 1 hour, until the vampire dies, or until the vampire dismisses them as a bonus action.

\DndMonsterSection{Legendary Actions}
Vampire can take 3 legendary actions, choosing from the options below. Only one legendary action can be used at a time and only at the end of another creature's turn. Vampire regains spent legendary actions at the start of its turn.

\begin{DndMonsterLegendaryActions}
\DndMonsterLegendaryAction{Move}{The vampire moves up to its speed without provoking opportunity attacks.
}
\DndMonsterLegendaryAction{Unarmed Strike}{The vampire makes one unarmed strike.
}
\DndMonsterLegendaryAction{Bite (Costs 2 Actions)}{The vampire makes one bite attack.
}
\end{DndMonsterLegendaryActions}
\DndMonsterSection{Regional Effects}
The region surrounding a vampire's lair is warped by the creature's unnatural presence, creating any of the following effects:


\begin{itemize}
\item There's a noticeable increase in the populations of bats, rats, and wolves in the region.
\item Plants within 500 feet of the lair wither, and their stems and branches become twisted and thorny.
\item Shadows cast within 500 feet of the lair seem abnormally gaunt and sometimes move as though alive.
\item A creeping fog clings to the ground within 500 feet of the vampire's lair. The fog occasionally takes eerie forms, such as grasping claws and writhing serpents.
\end{itemize}

If the vampire is destroyed, these effects end after 2d6 days.

\end{multicols}
\end{DndMonster}



\begin{DndMonster}[float=t]{Vampire Spawn}
\DndMonsterType{Medium undead, neutral evil}
\DndMonsterBasics[
armor-class = {15 (natural armor)},
hit-points = {\DndDice{11d8 + 33}},
speed = {30 ft.}
]
\DndMonsterAbilityScores[
 str = 16,
 dex = 16,
 con = 16,
 int = 11,
 wis = 10,
 cha = 12
]
\DndMonsterDetails[
saving-throws = {Dex +6, Wis +3},
skills = {Perception +3, Stealth +6},
damage-resistances = {necrotic; bludgeoning, piercing, slashing from nonmagical attacks},
senses = {darkvision 60 ft., passive perception 13},
languages = {the languages it knew in life},
challenge = 5,
]
\DndMonsterAction{Regeneration}
The vampire regains 10 hit points at the start of its turn if it has at least 1 hit point and isn't in sunlight or running water. If the vampire takes radiant damage or damage from holy water, this trait doesn't function at the start of the vampire's next turn.

\DndMonsterAction{Spider Climb}
The vampire can climb difficult surfaces, including upside down on ceilings, without needing to make an ability check.

\DndMonsterAction{Vampire Weaknesses}
The vampire has the following flaws:

\textit{Forbiddance.} The vampire can't enter a residence without an invitation from one of the occupants.

\textit{Harmed by Running Water.} The vampire takes 20 acid damage when it ends its turn in running water.

\textit{Stake to the Heart.} The vampire is destroyed if a piercing weapon made of wood is driven into its heart while it is incapacitated in its resting place.

\textit{Sunlight Hypersensitivity.} The vampire takes 20 radiant damage when it starts its turn in sunlight. While in sunlight, it has disadvantage on attack rolls and ability checks.

\DndMonsterSection{Actions}
\DndMonsterAction{Multiattack}
The vampire makes two attacks, only one of which can be a bite attack.

\DndMonsterAction{Bite}
\textsl{Melee Weapon Attack:} 6 to hit, reach 5 ft., one willing creature, or a creature that is grappled by the vampire, incapacitated, or restrained. \textsl{Hit:} 6 (1d6 + 3) piercing damage plus 7 (2d6) necrotic damage. The target's hit point maximum is reduced by an amount equal to the necrotic damage taken, and the vampire regains hit points equal to that amount. The reduction lasts until the target finishes a long rest. The target dies if this effect reduces its hit point maximum to 0.

\DndMonsterAction{Claws}
\textsl{Melee Weapon Attack:} 6 to hit, reach 5 ft., one creature. \textsl{Hit:} 8 (2d4 + 3) slashing damage. Instead of dealing damage, the vampire can grapple the target (escape DC 13).

\end{DndMonster}


\chapter{Spells}
\addcontentsline{toc}{section}{Antimagic Field}


\DndSpellHeader%
{Antimagic Field}
{8th level Abjuration}
{1 action}
{Self (10 feet sphere)}
{V, S, a pinch of powdered iron or iron filings}
{Concentration, up to 1 hour}
A 10-foot-radius invisible sphere of antimagic surrounds you. This area is divorced from the magical energy that suffuses the multiverse. Within the sphere, spells can't be cast, summoned creatures disappear, and even magic items become mundane. Until the spell ends, the sphere moves with you, centered on you.

Spells and other magical effects, except those created by an artifact or a deity, are suppressed in the sphere and can't protrude into it. A slot expended to cast a suppressed spell is consumed. While an effect is suppressed, it doesn't function, but the time it spends suppressed counts against its duration.

\subparagraph{Targeted Effects}
Spells and other magical effects, such as \textit{magic missile} and \textit{charm person}, that target a creature or an object in the sphere have no effect on that target.

\subparagraph{Areas of Magic}
The area of another spell or magical effect, such as \textit{fireball}, can't extend into the sphere. If the sphere overlaps an area of magic, the part of the area that is covered by the sphere is suppressed. For example, the flames created by a wall of fire are suppressed within the sphere, creating a gap in the wall if the overlap is large enough.

\subparagraph{Spells}
Any active spell or other magical effect on a creature or an object in the sphere is suppressed while the creature or object is in it.

\subparagraph{Magic Items}
The properties and powers of magic items are suppressed in the sphere. For example, a +1 longsword in the sphere functions as a nonmagical longsword.

A magic weapon's properties and powers are suppressed if it is used against a target in the sphere or wielded by an attacker in the sphere. If a magic weapon or a piece of magic ammunition fully leaves the sphere (for example, if you fire a magic arrow or throw a magic spear at a target outside the sphere), the magic of the item ceases to be suppressed as soon as it exits.

\subparagraph{Magical Travel}
Teleportation and planar travel fail to work in the sphere, whether the sphere is the destination or the departure point for such magical travel. A portal to another location, world, or plane of existence, as well as an opening to an extradimensional space such as that created by the rope trick spell, temporarily closes while in the sphere.

\subparagraph{Creatures and Objects}
A creature or object summoned or created by magic temporarily winks out of existence in the sphere. Such a creature instantly reappears once the space the creature occupied is no longer within the sphere.

\subparagraph{Dispel Magic}
Spells and magical effects such as \textit{dispel magic} have no effect on the sphere. Likewise, the spheres created by different \textit{antimagic field} spells don't nullify each other.

\addcontentsline{toc}{section}{Delayed Blast Fireball}


\DndSpellHeader%
{Delayed Blast Fireball}
{7th level Evocation}
{1 action}
{150 feet}
{V, S, a tiny ball of bat guano and sulfur}
{Concentration, up to 1 minute}
A beam of yellow light flashes from your pointing finger, then condenses to linger at a chosen point within range as a glowing bead for the duration. When the spell ends, either because your concentration is broken or because you decide to end it, the bead blossoms with a low roar into an explosion of flame that spreads around corners. Each creature in a 20-foot-radius sphere centered on that point must make a Dexterity saving throw. A creature takes fire damage equal to the total accumulated damage on a failed save, or half as much damage on a successful one.

The spell's base damage is 12d6. If at the end of your turn the bead has not yet detonated, the damage increases by 1d6.

If the glowing bead is touched before the interval has expired, the creature touching it must make a Dexterity saving throw. On a failed save, the spell ends immediately, causing the bead to erupt in flame. On a successful save, the creature can throw the bead up to 40 feet. When it strikes a creature or a solid object, the spell ends, and the bead explodes.

The fire damages objects in the area and ignites flammable objects that aren't being worn or carried.

\addcontentsline{toc}{section}{Disintegrate}


\DndSpellHeader%
{Disintegrate}
{6th level Transmutation}
{1 action}
{60 feet}
{V, S, a lodestone and a pinch of dust}
{Instantaneous}
A thin green ray springs from your pointing finger to a target that you can see within range. The target can be a creature, an object, or a creation of magical force, such as the wall created by \textit{wall of force}.

A creature targeted by this spell must make a Dexterity saving throw. On a failed save, the target takes 10d6 + 40 force damage. The target is disintegrated if this damage leaves it with 0 hit points.

A disintegrated creature and everything it is wearing and carrying, except magic items, are reduced to a pile of fine gray dust. The creature can be restored to life only by means of a \textit{true resurrection} or a \textit{wish} spell.

This spell automatically disintegrates a Large or smaller nonmagical object or a creation of magical force. If the target is a Huge or larger object or creation of force, this spell disintegrates a 10-foot-cube portion of it. A magic item is unaffected by this spell.

\addcontentsline{toc}{section}{Dispel Magic}


\DndSpellHeader%
{Dispel Magic}
{3rd level Abjuration}
{1 action}
{120 feet}
{V, S}
{Instantaneous}
Choose one creature, object, or magical effect within range. Any spell of 3rd level or lower on the target ends. For each spell of 4th level or higher on the target, make an ability check using your spellcasting ability. The DC equals 10 + the spell's level. On a successful check, the spell ends.

\addcontentsline{toc}{section}{Dominate Monster}


\DndSpellHeader%
{Dominate Monster}
{8th level Enchantment}
{1 action}
{60 feet}
{V, S}
{Concentration, up to 1 hour}
You attempt to beguile a creature that you can see within range. It must succeed on a Wisdom saving throw or be charmed by you for the duration. If you or creatures that are friendly to you are fighting it, it has advantage on the saving throw.

While the creature is charmed, you have a telepathic link with it as long as the two of you are on the same plane of existence. You can use this telepathic link to issue commands to the creature while you are conscious (no action required), which it does its best to obey. You can specify a simple and general course of action, such as "Attack that creature," "Run over there," or "Fetch that object." If the creature completes the order and doesn't receive further direction from you, it defends and preserves itself to the best of its ability.

You can use your action to take total and precise control of the target. Until the end of your next turn, the creature takes only the actions you choose, and doesn't do anything that you don't allow it to do. During this time, you can also cause the creature to use a reaction, but this requires you to use your own reaction as well.

Each time the target takes damage, it makes a new Wisdom saving throw against the spell. If the saving throw succeeds, the spell ends.

\addcontentsline{toc}{section}{Dominate Person}


\DndSpellHeader%
{Dominate Person}
{5th level Enchantment}
{1 action}
{60 feet}
{V, S}
{Concentration, up to 1 minute}
You attempt to beguile a humanoid that you can see within range. It must succeed on a Wisdom saving throw or be charmed by you for the duration. If you or creatures that are friendly to you are fighting it, it has advantage on the saving throw.

While the target is charmed, you have a telepathic link with it as long as the two of you are on the same plane of existence. You can use this telepathic link to issue commands to the creature while you are conscious (no action required), which it does its best to obey. You can specify a simple and general course of action, such as "Attack that creature," "Run over there," or "Fetch that object." If the creature completes the order and doesn't receive further direction from you, it defends and preserves itself to the best of its ability.

You can use your action to take total and precise control of the target. Until the end of your next turn, the creature takes only the actions you choose, and doesn't do anything that you don't allow it to do. During this time you can also cause the creature to use a reaction, but this requires you to use your own reaction as well.

Each time the target takes damage, it makes a new Wisdom saving throw against the spell. If the saving throw succeeds, the spell ends.

\addcontentsline{toc}{section}{Eyebite}


\DndSpellHeader%
{Eyebite}
{6th level Necromancy}
{1 action}
{Self}
{V, S}
{Concentration, up to 1 minute}
For the spell's duration, your eyes become an inky void imbued with dread power. One creature of your choice within 60 feet of you that you can see must succeed on a Wisdom saving throw or be affected by one of the following effects of your choice for the duration. On each of your turns until the spell ends, you can use your action to target another creature but can't target a creature again if it has succeeded on a saving throw against this casting of eyebite.

\subparagraph{Asleep}
The target falls unconscious. It wakes up if it takes any damage or if another creature uses its action to shake the sleeper awake.

\subparagraph{Panicked}
The target is frightened of you. On each of its turns, the frightened creature must take the Dash action and move away from you by the safest and shortest available route, unless there is nowhere to move. If the target moves to a place at least 60 feet away from you where it can no longer see you, this effect ends.

\subparagraph{Sickened}
The target has disadvantage on attack rolls and ability checks. At the end of each of its turns, it can make another Wisdom saving throw. If it succeeds, the effect ends.

\addcontentsline{toc}{section}{Fireball}


\DndSpellHeader%
{Fireball}
{3rd level Evocation}
{1 action}
{150 feet}
{V, S, a tiny ball of bat guano and sulfur}
{Instantaneous}
A bright streak flashes from your pointing finger to a point you choose within range and then blossoms with a low roar into an explosion of flame. Each creature in a 20-foot-radius sphere centered on that point must make a Dexterity saving throw. A target takes 8d6 fire damage on a failed save, or half as much damage on a successful one.

The fire spreads around corners. It ignites flammable objects in the area that aren't being worn or carried.

\addcontentsline{toc}{section}{Incendiary Cloud}


\DndSpellHeader%
{Incendiary Cloud}
{8th level Conjuration}
{1 action}
{150 feet}
{V, S}
{Concentration, up to 1 minute}
A swirling cloud of smoke shot through with white-hot embers appears in a 20-foot-radius sphere centered on a point within range. The cloud spreads around corners and is Vision and Light. It lasts for the duration or until a wind of moderate or greater speed (at least 10 miles per hour) disperses it.

When the cloud appears, each creature in it must make a Dexterity saving throw. A creature takes 10d8 fire damage on a failed save, or half as much damage on a successful one. A creature must also make this saving throw when it enters the spell's area for the first time on a turn or ends its turn there.

The cloud moves 10 feet directly away from you in a direction that you choose at the start of each of your turns.

\addcontentsline{toc}{section}{Lightning Bolt}


\DndSpellHeader%
{Lightning Bolt}
{3rd level Evocation}
{1 action}
{Self (100 feet line)}
{V, S, a bit of fur and a rod of amber, crystal, or glass}
{Instantaneous}
A stroke of lightning forming a line 100 feet long and 5 feet wide blasts out from you in a direction you choose. Each creature in the line must make a Dexterity saving throw. A creature takes 8d6 lightning damage on a failed save, or half as much damage on a successful one.

The lightning ignites flammable objects in the area that aren't being worn or carried.

\addcontentsline{toc}{section}{Misty Step}


\DndSpellHeader%
{Misty Step}
{2nd level Conjuration}
{1 bonus}
{Self}
{V}
{Instantaneous}
Briefly surrounded by silvery mist, you teleport up to 30 feet to an unoccupied space that you can see.

\addcontentsline{toc}{section}{Phantasmal Killer}


\DndSpellHeader%
{Phantasmal Killer}
{4th level Illusion}
{1 action}
{120 feet}
{V, S}
{Concentration, up to 1 minute}
You tap into the nightmares of a creature you can see within range and create an illusory manifestation of its deepest fears, visible only to that creature. The target must make a Wisdom saving throw. On a failed save, the target becomes frightened for the duration. At the end of each of the target's turns before the spell ends, the target must succeed on a Wisdom saving throw or take 4d10 psychic damage. On a successful save, the spell ends.

\addcontentsline{toc}{section}{Plane Shift}


\DndSpellHeader%
{Plane Shift}
{7th level Conjuration}
{1 action}
{Touch}
{V, S, a forked, metal rod worth at least 250 gp, attuned to a particular plane of existence}
{Instantaneous}
You and up to eight willing creatures who link hands in a circle are transported to a different plane of existence. You can specify a target destination in general terms, such as the City of Brass on the Elemental Plane of Fire or the palace of Dispater on the second level of the Nine Hells, and you appear in or near that destination. If you are trying to reach the City of Brass, for example, you might arrive in its Street of Steel, before its Gate of Ashes, or looking at the city from across the Sea of Fire, at the DM's discretion.

Alternatively, if you know the sigil sequence of a teleportation circle on another plane of existence, this spell can take you to that circle. If the teleportation circle is too small to hold all the creatures you transported, they appear in the closest unoccupied spaces next to the circle.

You can use this spell to banish an unwilling creature to another plane. Choose a creature within your reach and make a melee spell attack against it. On a hit, the creature must make a Charisma saving throw. If the creature fails this save, it is transported to a random location on the plane of existence you specify. A creature so transported must find its own way back to your current plane of existence.

\addcontentsline{toc}{section}{Power Word Stun}


\DndSpellHeader%
{Power Word Stun}
{8th level Enchantment}
{1 action}
{60 feet}
{V}
{Instantaneous}
You speak a word of power that can overwhelm the mind of one creature you can see within range, leaving it dumbfounded. If the target has 150 hit points or fewer, it is stunned. Otherwise, the spell has no effect.

The stunned target must make a Constitution saving throw at the end of each of its turns. On a successful save, this stunning effect ends.

\addcontentsline{toc}{section}{Reverse Gravity}


\DndSpellHeader%
{Reverse Gravity}
{7th level Transmutation}
{1 action}
{100 feet}
{V, S, a lodestone and iron filings}
{Concentration, up to 1 minute}
This spell reverses gravity in a 50-foot-radius, 100-foot-high cylinder centered on a point within range. All creatures and objects that aren't somehow anchored to the ground in the area fall upward and reach the top of the area when you cast this spell. A creature can make a Dexterity saving throw to grab onto a fixed object it can reach, thus avoiding the fall.

If some solid object (such as a ceiling) is encountered in this fall, falling objects and creatures strike it just as they would during a normal downward fall. If an object or creature reaches the top of the area without striking anything, it remains there, oscillating slightly, for the duration.

At the end of the duration, affected objects and creatures fall back down.

\end{document}
